\documentclass[10pt, oneside]{book}

\usepackage{amsmath}
\usepackage{amsfonts}
\usepackage{amssymb}
\usepackage{amsthm}
\usepackage{array}
\usepackage{bookmark}
\usepackage{bm}
\usepackage[justification = centering]{caption}
\usepackage{changepage}
\usepackage{color, soul}
\usepackage{enumitem}
\usepackage{fancyhdr}
\usepackage{float}
\usepackage[margin=1in]{geometry}
\usepackage{graphicx}
\usepackage{hyperref}
\usepackage[all]{hypcap}
\usepackage[utf8]{inputenc}
\usepackage{longtable}
\usepackage[version=4]{mhchem}
\usepackage{multicol}
\usepackage{pgfplots}
\usepackage[most]{tcolorbox}
\usepackage{titlesec}
\usepackage{tikz-cd}
\usepackage{wrapfig}
\usepackage[usenames,dvipsnames]{xcolor}

% ---------- Header / Footer ----------
\fancypagestyle{plain}{
    \fancyhf{}
    \fancyhead[L]{MATH 202A: Introduction to Topology and Analysis}
    \fancyhead[R]{William Nguyen}
    \fancyfoot[C]{\thepage}
}

\definecolor{pen}{HTML}{212120}
\definecolor{theoremFont}{HTML}{94330E}
\definecolor{theoremBg}{HTML}{FAF1EA}
\definecolor{exampleFont}{HTML}{B58900}
\definecolor{exampleBg}{HTML}{FFF8E6}
\definecolor{definitionFont}{HTML}{538ACA}
\definecolor{definitionBg}{HTML}{F3F5FC}
\definecolor{proofBg}{HTML}{FEFDFB}
\definecolor{proofFont}{HTML}{94330E}
\definecolor{solutionFont}{HTML}{B58900}
\definecolor{solutionBg}{HTML}{FFFBEF}
\definecolor{remarkFont}{HTML}{149B52}
\definecolor{remarkBg}{HTML}{ECFAEB}


\color{pen}


\newtcolorbox{theoremBox}[1]{
  enhanced,
  skin=enhancedmiddle,
  colback=theoremBg,
  sharp corners = all,
  boxrule=0pt,
  arc=5pt,
  left=8pt,
  right=8pt,
  top=8pt,
  bottom=8pt,
  grow to left by=11pt,    
  grow to right by=11pt,
  frame hidden,
  borderline west= {1mm}{0mm}{theoremFont},
  fonttitle=\large\sffamily,
  coltitle=theoremFont,
  title={#1},
  attach title to upper,
  after title = {\\ \vspace{-8pt} \\}  
}

\newtcolorbox{definitionBox}[1]{
  enhanced,
  skin=enhancedmiddle,
  colback=definitionBg,
  sharp corners = all,
  boxrule=0pt,
  arc=5pt,
  left=8pt,
  right=8pt,
  top=8pt,
  bottom=8pt,
  grow to left by=11pt,    
  grow to right by=11pt,
  frame hidden,
  borderline west= {1mm}{0mm}{definitionFont},
  fonttitle=\large\sffamily,
  coltitle=definitionFont,
  title={#1},
  attach title to upper,
  after title = {\\ \vspace{-8pt} \\}  
}

\newtcolorbox{remarkBox}[1]{
  enhanced,
  skin=enhancedmiddle,
  colback=remarkBg,
  sharp corners = all,
  boxrule=0pt,
  arc=5pt,
  left=8pt,
  right=8pt,
  top=8pt,
  bottom=8pt,
  grow to left by=11pt,    
  grow to right by=11pt,
  frame hidden,
  borderline west= {1mm}{0mm}{remarkFont},
  fonttitle=\bfseries\large\sffamily,
  coltitle=remarkFont,
  title={#1},
  attach title to upper,
  after title = {\\ \vspace{-8pt} \\}  
}

\newtcolorbox{proofBox}[1]{
  enhanced,
  skin=enhancedmiddle,
  colback=proofBg,
  sharp corners = all,
  boxrule=0pt,
  arc=5pt,
  left=8pt,
  right=8pt,
  top=8pt,
  bottom=8pt,
  grow to left by=11pt,    
  grow to right by=11pt,
  frame hidden,
  borderline west= {1mm}{0mm}{proofFont},
  fonttitle=\bfseries\fontsize{11pt}{13pt}\selectfont\sffamily,
  coltitle=proofFont,
  title={#1},
  attach title to upper,
  after title = { },
  breakable
}

\newtcolorbox{solutionBox}[1]{
  enhanced,
  skin=enhancedmiddle,
  colback=solutionBg,
  sharp corners = all,
  boxrule=0pt,
  arc=5pt,
  left=8pt,
  right=8pt,
  top=8pt,
  bottom=8pt,
  grow to left by=11pt,    
  grow to right by=11pt,
  frame hidden,
  borderline west= {1mm}{0mm}{solutionFont},
  fonttitle=\bfseries\fontsize{11pt}{13pt}\selectfont\sffamily,
  coltitle=solutionFont,
  title={#1},
  attach title to upper,
  after title = { },
  breakable
}

\newtcolorbox{exampleBox}[1]{
  enhanced,
  skin=enhancedmiddle,
  colback=exampleBg,
  sharp corners = all,
  boxrule=0pt,
  arc=5pt,
  left=8pt,
  right=8pt,
  top=8pt,
  bottom=8pt,
  grow to left by=11pt,    
  grow to right by=11pt,
  frame hidden,
  borderline west= {1mm}{0mm}{exampleFont},
  fonttitle=\large\sffamily,
  coltitle=exampleFont,
  title={#1},
  attach title to upper,
  after title = {\\ \vspace{-8pt} \\},
  breakable  
}

\renewenvironment{proof}[1][]{\vspace{-16pt}\begin{proofBox}}{\end{proofBox}\vspace{4pt}}
\newenvironment{solution}[1][]{\vspace{-16pt}\begin{solutionBox}}{\end{solutionBox}\vspace{4pt}}
\newenvironment{theorem}[1][]{\boxsubsection{#1}\begin{theoremBox}{#1}}{\end{theoremBox}\vspace{4pt}}
\newenvironment{definition}[1][]{\boxsubsection{#1}\begin{definitionBox}{#1}}{\end{definitionBox}\vspace{4pt}}
\newenvironment{remark}[1][]{\begin{remarkBox}}{\end{remarkBox}\vspace{4pt}}
\newenvironment{example}[1][]{\begin{exampleBox}}{\end{exampleBox}\vspace{4pt}}
\newenvironment{lemma}[1][]{\begin{theoremBox}{#1}}{\end{theoremBox}\vspace{4pt}}

\titleformat{\chapter}
  [display]
  {}
  {\huge\bfseries Chapter \thechapter}
  {1ex}
  {\raggedright\Huge\bfseries}






\hypersetup{
    hypertexnames=false,
    hidelinks
}

\pgfplotsset{compat=1.18}

\pgfplotsset{
    every axis/.style={
        axis lines=middle,   % axes through the origin
        ticks=none,          % no tick marks or numbers
        xlabel={}, ylabel={}, % no axis labels
        clip=true,          % allow plots beyond axes
    },
}


\graphicspath{{./Images/}}
\newcommand{\defcolor}{Goldenrod!35!white}
\newcommand{\theoremcolor}{Cerulean!15!white}

\newcommand{\R}{\mathbb{R}}
\newcommand{\Z}{\mathbb{Z}}
\newcommand{\Q}{\mathbb{Q}}
\newcommand{\N}{\mathbb{N}}
\newcommand{\C}{\mathbb{C}}
\renewcommand{\P}{\mathcal{P}}
\newcommand{\lda}{\lambda}
\newcommand{\norm}[1]{\lVert {#1} \rVert}
\newcommand{\abs}[1]{\left\lvert {#1} \right\rvert}
\newcommand{\dist}{\text{dist}}
\newcommand{\proj}{\text{proj}}
\newcommand{\Span}{\text{Span}}
\newcommand{\brac}[1]{\left\langle {#1} \right\rangle}
\newcommand{\Nul}{\text{Nul}\,}
\newcommand{\Col}{\text{Col}\,}
\newcommand{\Row}{\text{Row}\,}
\newcommand{\rank}{\text{rank}\,}
\newcommand{\quand}{\quad \text{and} \quad}
\newcommand{\qqquad}{\qquad \qquad}
\renewcommand{\div}{\mid}
\newcommand{\ndiv}{\nmid}
\renewcommand{\mod}[1]{\ (\text{mod } {#1})}
\renewcommand{\hat}[1]{\overline{#1}}
\newcommand{\Rel}{\ R\ }
\newcommand{\id}{\text{id}}
\newcommand{\del}{\partial}
\newcommand{\res}[1]{\left|_{#1}\right.}
\newcommand{\col}{\, :\, }
\newcommand{\codim}{\text{codim}\mkern1mu}
\renewcommand{\H}{\mathbb{H}}
\newcommand{\Int}{\text{Int}\mkern1mu}
\newcommand{\U}{\mathcal{U}}
\newcommand{\V}{\mathcal{V}}
\renewcommand{\cal}[1]{\mathcal{#1}}
\newcommand{\floor}[1]{\lfloor {#1} \rfloor}
\renewcommand{\tt}[1]{\text{#1}}

\renewcommand\arraystretch{}

\newcommand{\defspace}{\phantom{0} \vspace{0.2cm}\\}
\newcommand{\sectionspace}{\vspace{0.2cm}}
\newcommand{\chapterspace}{\vspace{2.4cm}}
\renewcommand{\indent}{\hspace{0.4cm}}
\newcommand{\exspace}{\vspace{0.6cm}}

\newcommand{\tocchapter}[1]{\chapter{#1}}
\newcommand{\tocsection}[1]{\section*{#1} \phantomsection \addcontentsline{toc}{section}{#1} \sectionspace}
\newcommand{\tocsubsection}[1]{\subsection*{#1} \phantomsection \addcontentsline{toc}{subsection}{#1} \sectionspace}
\newcommand{\boxsubsection}[1]{\phantomsection \addcontentsline{toc}{subsection}{#1}}
\newcommand{\tocsubsubsection}[1]{\subsubsection*{#1} \phantomsection \addcontentsline{toc}{subsubsection}{#1}}
\newcommand{\fig}[3]{\begin{figure}[H] \centering \includegraphics[scale = {#1}]{{#2}} \caption*{#3} \end{figure}}
\renewcommand{\qed}{\hfill\blacksquare\\}
\newcommand{\miniqed}{\hfill\square\\}
\newcommand{\longline}{\noindent\makebox[\linewidth]{\rule{\paperwidth - 2in}{0.4pt}}\vspace{0.4cm}}
\newcommand{\grid}[2]{\draw[step=1, gray!30] (-{#2},-{#1}) grid ({#2},{#1}); \filldraw (0, 0) circle (2pt);}
\newcommand{\commutes}{\arrow[phantom]{ul}{\circlearrowleft}}
\newcommand{\torus}[4]{\begin{scope}[shift = {(#1, #2)}, scale = {0.5*#3}, rotate = {#4}] \draw[thick] (-2, 0) to[out = 90, in = 180] (0, 1) to[out = 0, in = 90] (2, 0) to[out = -90, in = 0] (0, -1) to[out = 180, in = -90] cycle; \draw[thick] (-1, 0.1) to[out = -20, in = 200] (1, 0.1); \draw[thick] (-0.7, -0) to[out = 20, in = 160] (0.7, -0); \end{scope}}
\newcommand{\sphere}[4]{\begin{scope}[shift = {(#1, #2)}, scale = {#3}, rotate = {#4}] \draw[thick] (0, 0) circle (1); \draw[thick] (-1, -0.03) to[out = -25, in = 205] (1, -0.03); \draw[thick, dotted] (-1, 0.03) to[out = 15, in = 165] (1, 0.03); \end{scope}}
\newcommand{\transv}{\mathrel{\text{\tpitchfork}}}
\makeatletter
\newcommand{\tpitchfork}{%
  \vbox{
    \baselineskip\z@skip
    \lineskip-.52ex
    \lineskiplimit\maxdimen
    \m@th
    \ialign{##\crcr\hidewidth\smash{$-$}\hidewidth\crcr$\pitchfork$\crcr}
  }%
}
\makeatother


\let\vect\mathbf
\let\cal\mathcal

\titleformat*{\subsection}{%
    \fontsize{14}{15}\bfseries%
}
\titleformat*{\section}{%
    \fontsize{16}{19}\bfseries\filcenter%
}


\setlength{\parindent}{0px}
\setcounter{tocdepth}{2}


\begin{document}

\frontmatter
\begin{titlepage}
    \centering
    \vspace*{8cm}

    {\Huge \textbf{MATH 202A: Introduction to Topology and Analysis}}\\[1cm]
    {\Large By William Nguyen}\\[0.5cm]
    {\large Fall 2026}\\
    \vfill
    
\end{titlepage}

\tableofcontents
\newpage
\pagenumbering{arabic}

\clearpage
\mainmatter
\pagestyle{plain}

\tocsection{\textbf{8/25/26 Lecture 1: }Introduction}

\begin{definition}[\textbf{Definition 1.1:} Metric Spaces]
    Given a set $X$, $d : X \times X \to \R$ is a \emph{metric} if
    \begin{enumerate}
        \item $d(x, y) = 0 \iff x = y$
        \item $d(x, y) = d(y, x)$
        \item $d(x, z) \leq d(x, y) + d(y, z)$
    \end{enumerate}
    We then say $(X, d)$ is a \emph{metric space} 
\end{definition}

Lets start with $\Q$ with it's metric being $d\left(\frac{m_1}{n_1}, \frac{m_2}{n_2}\right) = \abs{\frac{m_1}{n_1} - \frac{m_2}{n_2}}$

\subsection*{Construction of the Reals}

First, why do we want $\R$? The answer is that in $\Q$, we don't have solutions to many polynomial equations, and that sets don't have well defined suprema and infima. 

\begin{definition}[\textbf{Definition 1.2: }Cauchy Sequences]
	We say a sequence $\{x_n\} \in X$ is \emph{Cauchy} if $\forall \epsilon > 0$, $\exists N > 0$ such that $\forall n, m > N$,
	\[ d(x_n, x_m) < \epsilon \]
\end{definition}
Now, the goal is to define real numbers as Cauchy sequences of $\Q$, but right now, we have a problem with unique representation. There are many Cauchy sequences that approach the same limit, so we can fix that.\\

Let $X' = \{\{x_n\} \in \Q \col \{x_n\} \text{ is Cauchy}\}$, and define the following equivalence relation. Say
\[ \{x_n\} \sim \{y_n\} \text{ if, and only if } d(x_n, y_n) \to 0 \text{ as } n \to \infty \]
With this, we can define $\R$ to be equivalence classes of $X'$. Note that with this definition, we can define a natural metric on $\R$ by
\[ d_{\R}([x_n], [y_n]) = \lim_{n \to \infty} d_{\Q}(x_n, y_n) \]
and this is well defined (choice of representative is irrelevant). 

\begin{definition}[\textbf{Definition 1.3: }Completeness]
	We say a metric space $X$ is \emph{complete} if all Cauchy sequences converge.
\end{definition}

\begin{theorem}[\textbf{Theorem 1.4:} Completeness of $\R$]
	With $d(x, y) = \abs{x - y}$, $\R$ is complete. 
\end{theorem}

Here are some interesting examples of metric spaces. 

\begin{itemize}
	\item Given a connected graph, let $d(x, y)$ be the length of the shortest path between the vertices $x$ and $y$. 
	\item Given a set of strings of length $n$ from an alphabet, let $d(w_1, w_2)$ be the number of entries in which $w_1$ and $w_2$ differ. This is called the \emph{Hamming distance}.
	\item Consider $\Q$, and for $x, y \in \Q$, write $\abs{x - y} = 2^{k}\frac{a}{b}$ where $a$ and $b$ are odd. This representation is unique, so define $d(x, y) = 2^{-k}$. This is called the 2-adic metric. It turns out by Ostrowski's theorem, every nontrivial absolute value on $\Q$ is either the normal one, or a $p$-adic one, which is cool!
	\item Consider $X = \{\text{functions on } [0, 1]\}$ with $d(f, g) = \sup\abs{f(x) - g(x)}$. This is well defined since $[0, 1]$ is compact. 
\end{itemize}

\begin{definition}[\textbf{Definition 1.5: }Completions]
	Given a metric space $(X, d)$, we say that $(X^{*}, d^{*})$ is a \emph{completion} of $(X, d)$ if it is complete and there exists an injective function $f : X \to X^{*}$ such that
	\[ d^{*}(f(x), f(y)) = d(x, y) \]
	and also that $f(X)$ is dense in $X^{*}$. By dense, we mean that for any $x^{*} \in X^{*}$, there is a sequence of points in $X$ that approach it. 
\end{definition}

With this definition, we see that $\R$ is a completion of $\Q$, and we take $f(x) = \{x_n = x\}$ to be the injective function. A quick question one may ask is why we need density? The reason is that we want uniqueness. For example, we don't want to say that $\R$ with an additional far point off the real line is a completion of $\Q$.

\begin{theorem}[\textbf{Theorem 1.6: }Existence of Completion]
	Every metric space has a completion
\end{theorem}
\begin{proof}{Proof Sketch:}
	Take the desired set to be equivalence classes of Cauchy sequences in $X$, and show that this is a completion.
\end{proof}

Now, we also want to show that completions of metric spaces are unique, but we need a notion of equivalence first. 

\begin{definition}[\textbf{Definition 1.7: }Isometries]
	A map $f : X \to Y$ of metric spaces is an \emph{isometry} if
	\[ d_{X}(x, y) = d_{Y}(f(x), d(f(y))) \]
	An \emph{isomorphism} is a bijective isometry.
\end{definition}

\begin{theorem}[\textbf{Theorem 1.8: }Uniqueness of Metric Spaces]
	A completion of a metric space is unique up to isomorphism.
\end{theorem}

An example of such an isomorphism is letting $X = \R$, $Y = \{(x, 0) \in \R^2\}$, and $f(x) = (x, 0)$, where $X$ and $Y$ have the Euclidean metric. It can be easily checked that $f$ is an isomorphism. 
\clearpage

\tocsection{\textbf{8/28/26 Lecture 2: }Introduction to Measure}

We want to assign a size to subsets $A \subseteq \R$. If $A = [a, b]$, it's quite natural to say its size is $b - a$. Also, we would want the size of $\cup I_{n}$ with $I_{i}$ are disjoint to be the sum of the sizes of $I_{n}$. 

\subsection*{Digression on Probability Theory}
\sectionspace

We can model a sequence of coin flips as an element if $\{T, H\}^{\N}$. However, here is another way to do it. For $x \in [0, 1]$, write $x$ in binary like $x = (0.1001\ldots)_{2}$ for instance. Then, associate coin tosses to points in $[0, 1]$ by mapping $T \to 0$ and $H \to 1$, so that $HTTH\ldots \to (0.1001\ldots)_{2}$. This allows us to interpret statements about probability as statements about sizes of subsets of $[0, 1]$.\\

For instance, 
\begin{align*}
    \{x \in \{T, H\}^{\N} \col x \text{ begins with an }H\} &= \{x \in [0, 1] \col x = 0.1\ldots\}\\
    &= [1/2, 1]
\end{align*}
Thus, one can conclude that the probability that a sequence of flips starts with an $H$ is $1/2$.\\

Here's another question, one that is harder to solve: How do we show that the number of head and tails get closer and closer to being equal with probability $1$? Let's try to represent this set of sequences as a subset of $[0, 1]$. It is clear that the set of such sequences is in correspondence with
\[ \mathcal{N} = \left\{x \in [0, 1] \col \lim_{n \to \infty} \frac{1}{N}\sum_{i = 1}^{N} b_{n}(x) = 1/2\right\} \]
where $b_{n}$ is the $n$th binary digit of $x$. However, this has limits and, so is there a way to make this more digestible? The answer is yes. Consider
\[ \bigcap_{k = 1}^{\infty} \bigcup_{N = 1}^{\infty} \bigcap_{n = N}^{\infty} \left\{x \in [0, 1] \col \abs{\frac{1}{n}\sum_{i = 1}^{n} b_{i}(x) - 1/2} < 1/k\right\} \]
Parsing what this set is, we see it says that belonging in this set means that ``for all $k$, there exists an $N$, such that for all $n \geq N$, $\abs{\frac{1}{n}\sum_{i = 1}^{n} b_{i}(x) - 1/2} = 1/k$''. That is, we just turned the notion of a limit into the intersection and union of simpler sets! Soon, we will see this technique makes it much easier to measure sets with properties like limits approaching a certain value. Before that, some topology.

\begin{definition}[\textbf{Definition 2.1:} Negligible Sets]
    A set $A$ is \emph{negligible} if, $\forall \epsilon > 0$, there is a countable collection of (open?) intervals $\{I_{n}\}_{n = 1}^{\infty}$ such that
    \[ A \subseteq \bigcup_{n = 1}^{\infty} I_{n} \]
    and $\sum_{n = 1}^{\infty} \abs{I_{n}} < \epsilon$. 
\end{definition}
\clearpage

The following is a theorem we won't prove, but we will state it here. 

\begin{theorem}[\textbf{Theorem 2.2: }Borel's Normal Number Theorem]
    The set $\mathcal{N}^{c} = [0, 1] \backslash \cal{N}$ is negligible. 
\end{theorem}
Thus, we can conclude that the probability that the number of heads and tails aren't equal in the long run occurs with probability 0. Moving past probability, let's explore some other useful results about this topological/measure theoretic term we just defined.

\begin{lemma}[\textbf{Proposition 2.3}]
    If $A$ is countable, then it is negligible. 
\end{lemma}
\begin{proof}{Proof:}
    Let $\epsilon > 0$ and $A = \{a_1, a_2, \ldots\}$. Then, letting $I_{n} = (a_{n} - \epsilon / 2^{n + 2}, a_{n} + \epsilon / 2^{n + 2})$, we see that $A \subseteq \cup_{n = 1}^{\infty} I_{n}$. Then, as $\abs{I_{n}} = \epsilon/2^{n + 1}$, we see
    \[ \sum_{b = 1}^{\infty} \abs{I_{n}} = \sum_{n = 1}^{\infty} \epsilon / 2^{n + 1} = \epsilon/2 < \epsilon \]
    Thus, $A$ is countable. 

    $\qed$
\end{proof}

\begin{lemma}[\textbf{Proposition 2.4}]
    Let $\{A_{n}\}_{n = 1}^{\infty}$ be a collection of negligible sets. Then $\cup_{n = 1}^{\infty} A_{n}$ is also negligible. 
\end{lemma}
\begin{proof}{Proof:}
    Let $\epsilon > 0$. For each $A_{i}$, we can find a collection of intervals $\{I_{n, i}\}$ such that $A_i \subseteq \cup_{n = 1}^{\infty} I_{n, i}$ and $\sum_{n = 1}^{\infty} \abs{I_{n, i}} < \epsilon / 2^{i}$.\\
    
    Now, consider the set $\{I_{n, i} \col i, n \in \N\}$. We see that this is a countable collection of intervals, $\cup_{i = 1}^{\infty} \subseteq \cup_{n, i \geq 1} I_{n, i}$, and
    \begin{align*}
        \sum_{n, i \geq 1} \abs{I_{n, i}} &= \sum_{i = 1}^{\infty} \sum_{n = 1}^{\infty} \abs{I_{n, i}}\\
        &< \sum_{i = 1}^{\infty} \epsilon / 2^{i}\\
        &= \epsilon
    \end{align*} 
    Thus, $\cup_{n = 1}^{\infty} A_{n}$ is negligible.

    $\qed$
\end{proof}
\clearpage

Now, small contingency. Is it possible that we can have a $\cal{N} \subseteq [0, 1]$ with both $\cal{N}^{c} = [0, 1] \backslash \cal{N}$ and $\cal{N}$ being negligible? Do non-negligible sets even exist? Our intuition says this shouldn't occur, but it is a suprisingly nontrivial result to show! 

\begin{lemma}[\textbf{Proposition 2.5}]
    $[0, 1]$ is not negligible. 
\end{lemma}
Before we prove this result however, we need some more tools!

\begin{definition}[\textbf{Definition 2.6: }Open Covers]
    Let $C \subseteq \R$. An \emph{open cover} $\cal{O} = \{\cal{U}_{\alpha}\}_{\alpha \in I}$ is a collection of open sets such that $C \subseteq \cup_{\alpha \in I} \cal{U}_{\alpha}$. 
\end{definition}

\begin{definition}[\textbf{Definition 2.7: }Compactness]
    Let $C \subseteq \R$. We say that $C$ is \emph{compact} if every open cover of $C$ has a finite subcover. That is, given an open cover $C \subseteq \cup_{\alpha \in I} \cal{U}_{\alpha}$, then we can write $C \subseteq \cup_{i = 1}^{n} \cal{U}_{\alpha_1}$ for some $\cal{U}_{\alpha_i} \in \{\cal{U}_{\alpha}\}$. 
\end{definition}

\begin{theorem}[\textbf{Theorem 2.8: }Heine-Borel Theorem]
    A set $C \subseteq \R$ is compact if, and only if, $C$ is closed and bounded.
\end{theorem}

Now, let's try to prove Proposition 2.5 To do this, we will prove the following lemma. The full proof will be shown next lecture, but here is the proof sketch. 

\begin{lemma}[\textbf{Lemma 2.9}]
    If $\cal{O} = \{I_\alpha\}_{\alpha \in J}$ is a collection of open intervals such that $[a, b] \subseteq \cup_{\alpha} I_{\alpha}$, then $\sum_{\alpha} \abs{I_{\alpha}} > b - a$. 
\end{lemma}
\begin{proof}{Proof Sketch:}
    Given an open cover (of intervals), pass to a finite cover. Thus, we only need to show this result for finite collections. To do this, we will use induction. For $k = 1$, we see that if $[a, b] \subseteq [c, d]$, indeed $d - c \geq b - a$. For the inductive step, wlog, write $[a, b] = [0, 1]$. Then, considering a cover of $k + 1$ intervals, note that one interval must contain the endpoint $1$. Let this interval be $I^{*}$. Then, apply the inductive step to the interval $[0, 1] \backslash I^{*}$ as this is covered with $k$ intervals. 
\end{proof}
\clearpage
\tocsection{\textbf{8/31/26 Lecture 3: }Introduction to Measure}

\begin{lemma}[\textbf{Proposition 3.1}]
    A set $\cal{U} \subseteq \R$ is open if $\cal{U}$ can be written as a union of countable many disjoint open intervals.
\end{lemma}

\begin{theorem}[\textbf{Theorem 3.2: }Existence of Non-Negligible Sets]
    The interval $[0, 1]$ is not negligible
\end{theorem}
Before we prove this theorem, we need a short lemma.

\begin{lemma}[\textbf{Lemma 3.3}]
    If $\{I_n\}_{n = 1}^{N}$ is a finite collection so that $[a, b] \subseteq \bigcup_{i = 1}^{N} I_{i}$, then $[a, b] \leq \sum_{i = 1}^{N} \abs{I_{i}}$. 
\end{lemma}
\begin{proof}{Proof:}
    We will use induction on $N$, the size of the collection. First, for the $N = 1$, we see clearly that if $[a, b] \subseteq (c, d)$, then $d - c > b - a$.\\

    Now, suppose it holds for $N - 1$, and suppose $\{I_{n}\}_{n = 1}^{N}$ cover $[a, b]$. Then, there must be an $n^{*}$ so that $b \in I_{n^*}$. Now, if $I_{n^{*}} = (c, d)$, then we see that $\{I_{n}\}_{n = 1}^{N - 1}$ must cover $[a, c] = [a, b] \backslash I_{n^{*}}$.\\

    By the inductive step, we know that $\sum_{i = 1}^{N - 1} \abs{I_{i}} \geq c - a$. Thus, we see that
    \[ \sum_{i = 1}^{N} \abs{I_{i}} \geq c - a + (d - c) = d - a \geq b - a \]
    which was to be proved.

    $\qed$
\end{proof}

Now, we proceed with the proof of Theorem 3.2.\\

\begin{proof}{\textbf{Proof of Theorem 3.2:}}
    Suppose that $[0, 1] \subseteq \cup_{n = 1}^{\infty} I_{n}$. By compactness, there exists a finite subcollection of the $I_{i}$ such that $[0, 1] \subseteq \cup_{n = 1}^{N} I_{n_k}$. Then, by our proposition, we know that $\sum_{k = 1}^{N} \abs{I_{n_k}} \geq 1$, so clearly $\sum_{n = 1}^{\infty} \abs{I_{n}} \geq 1$. Thus, all interval coverings have combined length greater than or equal to $1$, so it cannot be epsilon close to $0$. 

    $\qed$
\end{proof}

Now, moving past these topological notions, we want to try to construct a function $\lambda$ on some subsets $A \subseteq \R$ such that $\lambda(A)$ assigns a size to $A$. So far, we want that if $\cal{U} \subseteq \R$ and $\cal{U} = \cup_{n = 1}^{\infty} I_{n}$ with $I_{n}$ being disjoint intervals, then we should have $\lambda(\cal{U}) = \sum_{n} \abs{I_{n}}$.\\

Similarly, if we have an $A \subseteq [0, 1]$ that is closed, then as $A \backslash [0, 1] = \cal{U}$ is open, then we should have that $\lambda(A) = 1 - \lambda(\cal{U})$. \\

Here is a first attempt: For $A \in \P(\R)$ define
\begin{align*}
    \lambda^{*}(A) &= \inf\left\{\sum_{n = 1}^{\infty} \abs{I_{n}} \col A \subseteq \bigcup_{n = 1}^{\infty} I_{n}\right\}\\
    &= \inf\{\lambda(\cal{U}) \col \cal{U} \text{ is open and } A \subseteq \cal{U}\}
\end{align*}
This is actually not a bad attempt, and motivates the following definition!

\begin{definition}[\textbf{Definition 3.4: }Outer Measures]
    For a set $X$, a function $\mu^{*} : \P(X) \to [0, 1]$ is an \emph{outer measure} if
    \begin{enumerate}
        \item $\mu^{*}(\emptyset) = 0$
        \item $\mu^{*}\left(\bigcup_{n = 1}^{\infty} A_n\right) \leq \sum_{n = 1}^{\infty} \mu^{*}(A_n)$, $A_n \subseteq X$
        \item If $A \subseteq B$, then $\mu^{*}(A) \leq \mu^{*}(B)$
    \end{enumerate}
\end{definition}

\begin{lemma}[\textbf{Proposition 3.5}]
    $\lambda^{*}$, as defined above, is an outer measure.
\end{lemma}
\begin{proof}{Proof:}
    $1$ and $3$ follow immediately from the definition of $\lambda^{*}$, so we just need to prove $2$. Start by letting $\{A_n\} \subseteq \R$ and let $\epsilon > 0$. Then, for each $A_{n}$, pick a collection of intervals $\{I_{n, j}\}$ so that $A_n \subseteq \cup_{j} I_{n, j}$ and that $\sum_{j} \abs{I_{n, j}} \leq \lambda^{*}(A_n) + \epsilon/2^{n}$ (we can do this by the definition of an infimum). Then, considering the countable collection $\{I_{n, j}\}$ we see that $\cup_{n} A_n \subseteq \cup_{n, j} I_{n, j}$ and that
    \begin{align*}
        \lambda^{*}\left(\bigcup_{n} A_n\right) &\leq \sum_{n, j} \abs{I_{n, j}}\\
        &= \sum_{n}\left(\sum_{j} \abs{I_{n, j}}\right)\\
        &= \sum_{n} \lambda^{*}(A_n) + \epsilon
    \end{align*}
    Thus, $\lambda^{*}\left(\bigcup_{n} A_n\right) \leq \sum_{n} \lambda^{*}(A_n) + \epsilon$ for all $\epsilon$. Thus, $\lambda^{*}\left(\bigcup_{n} A_n\right) \leq \sum_{n} \lambda^{*}(A_n)$.

    $\qed$
\end{proof}

Now, we will come back to outer measures soon, but first, we should define what a measure is. To do that, we need a suitable class of subsets of $\R$ on which a measure can exist.

\begin{definition}[\textbf{Definition 3.6: }$\sigma$-Algebras]
    For a set $X$, a subset $\cal{A} \subseteq \P(X)$ is a \emph{$\sigma$-algebra} if
    \begin{enumerate}
        \item $\emptyset, X \in \cal{A}$
        \item If $A_n \in \cal{A}$, then $\cup_{n} A_n \in \cal{A}$
        \item If $A \in \cal{A}$ then $A^{c} = X \backslash A \in \cal{A}$
    \end{enumerate}
\end{definition}

\begin{lemma}[\textbf{Corollary 3.7}]
    If $A_n \in \cal{A}$, then $\cap_{n} A_n \in \cal{A}$
\end{lemma}
\begin{proof}{Proof:}
    This follows immediately from applying De Morgan's laws to the definition of a $\sigma$-algebra.

    $\qed$
\end{proof}

Some examples of $\sigma$-algebras on any $X$ are $\{\emptyset, X\}$ and $\P(X)$. With this definition, we can now define a measure.

\begin{definition}[\textbf{Definition 3.7: }Measures and Measure Spaces]
    A \emph{measure space} is a triple $(X, \cal{A}, \mu)$ where $\cal{A}$ is a $\sigma$-algebra and $\mu : \cal{A} \to [0, \infty]$ is a map that satisfies the following properties:
    \begin{enumerate}
        \item $\mu(\emptyset) = 0$
        \item If $A_n \in \cal{A}$ is a countable collection of pairwise disjoint sets, then $\mu(\cup_{n} A_n) = \sum_{n} \mu(A_n)$
    \end{enumerate}
    We call $\mu$ a \emph{measure}.
\end{definition}
\tocsection{\textbf{9/2/26 Lecture 4: }Measures and $\sigma$-Algebras}

Today, we will construct a measure on a $\sigma$-algebra using an outer measure.

\begin{example}{Examples}
    \begin{enumerate}
        \item For any set $X$, let $\mathcal{A} = \P(X)$. Then, we can let $\mu(A) = \#A$, where $\#A$ is the cardinality of $A$. We call $\mu$ the counting measure
        \item Let $X$ and $\mathcal{A} = \P(X)$. Then, fix $x \in X$, and define
        \[ \delta_{x}(A) = \left\{\begin{matrix} 1 & x \in A \\ 0 & \text{otherwise}\end{matrix}\right. \]
        Then, $\delta_{x}(A)$ is a measure, called a point mass.
    \end{enumerate}
\end{example}

\begin{lemma}[\textbf{Proposition 4.1}]
    Let $(X, \cal{A}, \mu)$ be a measure space. Then, the following hold.
    \begin{enumerate}
        \item Let $\{A_n\}_{n = 1}^{N}$ be a finite subset of $\mathcal{A}$ that are pairwise disjoint. Then, we have
        \[ \mu\left(\bigcup_{n = 1}^{N} A_{n}\right) = \sum_{n = 1}^{N} \mu(A_n) \]
        \item If $A, B \in \cal{A}$ and $A \subseteq B$, then $\mu(A) \leq \mu(B)$
        \item If $A_n \in \mathcal{A}$ for each $n \in \N$ not necessarily disjoint, then
        \[ \mu\left(\bigcup_{n = 1}^{\infty} A_n\right) \leq \sum_{n = 1}^{\infty} \mu(A_n) \]
        \item If $A_{1} \subseteq A_2 \subseteq \cdots$ are an increasing sequence of sets in $\cal{A}$, then
        \[ \mu\left(\bigcup_{n = 1}^{\infty} A_n\right) = \lim_{n \to \infty} \mu(A_n) \]
        \item If $A_1 \supseteq A_2 \supseteq \cdots$ is a decreasing sequence in $\cal{A}$ and $\mu(A_1) < \infty$, then
        \[ \mu\left(\bigcap_{n = 1}^{\infty} A_n\right) = \lim_{n \to \infty} \mu(A_n) \]
    \end{enumerate}
\end{lemma}
\begin{proof}{Proof:}
    
    \textbf{1)} Given $A_1, \ldots, A_N$, set $A_m = \emptyset$ for $m > N$. Then, we see that
    \[ \mu\left(\bigcup_{n = 1}^{\infty} A_n\right) = \sum_{n = 1}^{\infty} \mu(A_n) = \sum_{n = 1}^{N} \mu(A_n) \]
    and as $\bigcup_{n = 1}^{\infty} A_n = \bigcup_{n = 1}^{N} A_n$, we have our desired result.\\

    \textbf{2)} Since $B = A \cup (B \backslash A)$, then we see that
    \[ \mu(B) = \mu(A) + \mu(B \backslash A) \]
    as $A$ and $B \backslash A$ are disjoint. Then, as $\mu(B \backslash A) \geq 0$, we have that $\mu(B) \geq \mu(A)$.\\

    \textbf{3)} Define $B_1 = A_1$, $B_2 = A_2 \backslash A_1$ and $B_{3} = A_3 \backslash \left(A_1 \cup A_2\right)$. In general, let $B_{n} = A_n \backslash \left(\bigcup_{i = 1}^{n - 1} A_i\right)$. Now, we see that the $B_{i}$ are disjoint, so
    \[ \mu\left(\bigcup_{n = 1}^{\infty} A_{n}\right) = \mu\left(\bigcup_{n = 1}^{\infty} B_n\right) = \sum_{n = 1}^{\infty} \mu(B_n) \]
    Then, since $\mu(B_n) \leq \mu(A_n)$ as $B_n \subseteq A_n$, we see that
    \[ \mu\left(\bigcup_{n = 1}^{\infty} A_n\right) \leq \sum_{n = 1}^{\infty} \mu(A_n) \]

    \textbf{4)} This is like $5$ so we will do $5$.\\

    \textbf{5)} Let $C_{n} = A_1 \backslash A_n$. By $4$, since $C_{n}$ is increasing, we see that
    \[ \mu\left(\bigcup_{n = 1}^{\infty} C_n\right) = \lim_{n \to \infty} \mu(C_n) = \lim_{n \to \infty} \left( \mu(A_1) - \mu(A_n)\right) = \mu(A_1) - \lim_{n \to \infty}\left(\mu(A_n)\right) \]
    However, we see that
    \[ \bigcup_{n = 1}^{\infty} C_n = A_1 \backslash \bigcap_{n = 1}^{\infty} A_n \]
    so
    \[ \mu\left(\bigcup_{n = 1}^{\infty} C_n\right) = \mu(A_1) - \mu\left(\bigcap_{n = 1}^{\infty} A_n\right) \]
    Thus, we have our desired equality.

    $\qed$
\end{proof}

Note that the condition that $\mu(A_1)$ has finite measure is necessary in $5$. This can be seen by the following counterexample. Take the counting measure on $\R$, and consider the sequence of decreasing sets
\[ A_n = [n, \infty) \]
Then, we see that $\mu(A_n) = \infty$ for each $n$, but $\bigcap_{n} A_n = \emptyset$.

\begin{definition}[\textbf{Definition 4.2:} $\sigma$-algebra Generated by a Collection]
    Let $\cal{C}$ be a collection of subsets of a set $X$. Then, the \emph{$\sigma$-algebra generated by $\cal{C}$} is 
    \[ \sigma(\cal{C}) = \bigcap \{\cal{A} \col \cal{A} \text{ is a } \sigma \text{-algebra containing } \cal{C}\} \]
    Since the intersection of $\sigma$-algebras is a $\sigma$-algebra, $\sigma(\cal{C})$ is indeed a $\sigma$-algebra.
\end{definition}
In a sense, we say that $\sigma(\cal{C})$ is the smallest $\sigma$-algebra containing $\cal{C}$ (since any $\sigma$-algebra containing $\cal{C}$ must contain $\sigma(\cal{C})$). 

\begin{definition}[\textbf{Definition 4.3:} Borel $\sigma$-algebra]
    Let $X$ be a topological space, and let $G$ be it's set of open sets. Then, the \emph{Borel $\sigma$-algebra} is
    \[ \cal{B} := \sigma(G) \]
    We call $A \in \cal{B}$ a \emph{Borel set}.
\end{definition}

\begin{lemma}[\textbf{Proposition 4.4}]
    Consider $\R$ with the normal topology, and let $\cal{B}$ be the Borel $\sigma$-algebra. Then, $\cal{B}$ is generated by each of the following collections of subsets.
    \begin{enumerate}
        \item $\cal{C}_1 = \{(a, b) \col a, b \in \R\}$
        \item $\cal{C}_2 = \{[a, b] \col a, b \in \R\}$
        \item $\cal{C}_3 = \{(a, b] \col a, b \in \R\}$
        \item $\cal{C}_4 = \{(a, \infty) \col a \in \R\}$
    \end{enumerate}
\end{lemma}
\begin{proof}{\textbf{Proof:}}
    We won't do all of them, but we will show a few. The idea is to show that we can either generate the open sets with this collection, or the other way around.\\

    \textbf{1)} First, any open set can be written as $\cup_{n = 1}^{\infty} I_n$ where $I_n = (a, b), (-\infty, a)$ or $(a, \infty)$. Now, we can write
    \[ (a, \infty) = \cup_{n} (a, a + n) \]
    and similarly, 
    \[ (-\infty, a) = \cup_{n} (a - n, a) \]
    Thus, $(a, b), (-\infty, a)$ and $(a, \infty)$ are in $\sigma(\cal{C}_1)$. Thus, we see that $\cal{B} \subseteq \sigma(\cal{C}_1)$. Since $C_1 \subseteq \cal{B}$, we equality.\\
    
    \textbf{2)} The goal is to show that $\sigma(\cal{C}_2) \subseteq \cal{B}$ and that $\cal{B} \subseteq \sigma(\cal{C}_2)$. Since
    \[ [a, b] = \bigcap_{n} (a - 1/n, b + 1/n) \]
    we see that $\cal{C}_2 \subseteq \sigma({\cal{C}_1}) = \cal{B}$. Thus, $\sigma({\cal{C}_2}) \subseteq \cal{B}$. Then, as
    \[ (a, b) = \bigcup_{n} [a + 1/n, b - 1/n] \]
    we see that $\cal{C}_1 \subseteq \sigma(\cal{C}_2)$, so $\cal{B} = \sigma(\cal{C}_1) \subseteq \sigma(\cal{C}_2)$. Thus, we have equality.\\

    \textbf{3)} and \textbf{4)} are very similar, so we leave the proof to the reader.

    $\qed$
\end{proof}

\begin{definition}[\textbf{Definition 4.5:} $\mu^{*}$-Measurability]
    For an outer measure $\mu^{*}$ on a set $X$, we sey that $A \subseteq X$ is \emph{$\mu^{*}$-measurable} if, for any $E \subseteq X$, then
    \[ \mu^{*}(E) = \mu^{*}(E \cap A) + \mu^{*}(E \cap A^{c}) \]
\end{definition}

\tocsection{\textbf{9/4/26 Lecture 5: }More on Measures}

Now, given an outer measure $\mu^{*}$ and the definition of $\mu^{*}$-measurable sets, it turns out we can construct a $\sigma$-algebra and define a measure $\mu$ on it!

\begin{theorem}[\textbf{Theorem 5.1: }Construction of Measures from Outer Measures]
    Given an outer measure $\mu^{*}$, the collection of $\mu^{*}$-measurable sets, denoted $\mathcal{A}$, is a $\sigma$-algebra, and the restriction of $\mu^{*}$ to $\mathcal{A}$ is a measure. Additionally, $\mathcal{A}$ contains all the null sets of $\mu^{*}$
\end{theorem}
\begin{proof}{Proof:}
    First, we show that $\mathcal{A}$ is indeed a $\sigma$-algebra. Firstly, from the symmetry of the definition of $\mu^{*}$-measurability, we see that if $A \in \cal{A}$ then $A^{c} \in \cal{A}$.\\
    
    Now, we want to show that finite unions of sets in $\cal{A}$ are in $\cal{A}$. Firstly, notice that, in order to show that $A \in \cal{A}$, all we need is that for all $E \subseteq X$, $\mu^{*}(E) \supseteq \mu^{*}(E \cap A) + \mu^{*}(E \cap A^{c})$, as the opposite inequality follows from subadditivity.\\

    Let $A, B \in \cal{A}$. We want to show that $A \cup B \in \cal{A}$. First, for $E \subseteq X$, 
    \[ \mu^{*}(E) = \mu^{*}(E \cap A) + \mu^{*}(E \cap A^{c}) \]
    Then, applying the $\mu^{*}$-measurability of $B$ to the set $E \cap A^{c}$, we see taht
    \begin{align*}
        \mu^{*}(E) &= \mu^{*}(E \cap A) + \mu^{*}(E \cap A^{c} \cap B) + \mu^{*}(E \cap A^{c} \cap B^{c})\\
        &= \mu^{*}(E \cap A) + \mu^{*}(E \cap A^{c} \cap B) + \mu^{*}(E \cap (A \cup B)^{c})\\
        &\geq \mu^{*}(E \cap (A \cup B)) + \mu^{*}(E \cap (A \cup B)^{c})
    \end{align*}
    where the last inequality comes from applying subadditivity to the fact that $A \cup (A^{c} \cap B) = A \cup B$. Thus, we see that $A \cup B$ is $\mu^{*}$-measurable, and by induction, we see that all finite unions of $\mu^{*}$-measurable sets is $\mu^{*}$-measurable.\\

    Now, if $A_n \in \cal{A}$, we want to show that $\cup_{n} A_n \in \cal{A}$. Notice that it suffices to assume that $A_n$ are disjoint as we can write
    \[ B_1 = A_1,\quad B_2 = A_1 \backslash A_2,\cdots \quad B_n = A_n \backslash \bigcup_{i = 1}^{n - 1} A_i \]
    and see that the $B_i$ are disjoint, $\cup_{i = 1}^{\infty} B_i = \cup_{i = 1}^{\infty} A_i$, $B_i \in \cal{A}$, and $\mu^{*}(B_i) \leq \mu^{*}(A_i)$. Thus, assuming that the $A_i$ are disjoint, we have
    \begin{align*}
        \mu^{*}(E) &= \mu^{*}(E \cap A_1) + \mu^{*}(E \cap A_1^{c})\\
        &= \mu^{*}(E \cap A_1) + \mu^{*}(E \cap A_1^c \cap A_2) + \mu^{*}(E \cap A_1^{c} \cap A_2^{c})\\
        &= \mu^{*}(E \cap A_1) + \mu^{*}(E \cap A_2) + \mu^{*}(E \cap (A_1 \cup A_2)^{c})\\
        &\vdots\\
        &= \sum_{i = 1}^{N} \mu^{*}(E \cap A_i) + \mu^{*}\left(E \cap \left(\bigcup_{i = 1}^{N} A_i\right)^{c}\right)\\
        &\geq \sum_{i = 1}^{N} \mu^{*}(E \cap A_i) + \mu^{*}\left(E \cap \left(\bigcup_{i = 1}^{\infty} A_i\right)^{c}\right)
    \end{align*}
    for all $N$, where the last step holds from the fact that $\left(\bigcup_{i = 1}^{N} A_i\right)^{c} \supseteq \left(\bigcup_{i = 1}^{\infty} A_i\right)^{c}$. Then, as this holds for all $N$, we can pass to the limit to see that
    \begin{align*}
        \mu^{*}(E) &\geq \sum_{i = 1}^{\infty} \mu^{*}(E \cap A_i) + \mu^{*}\left(E \cap \left(\bigcup_{i = 1}^{\infty} A_i\right)^{c}\right)\\
        &\geq \mu^{*}\left(E \cap \left(\bigcup_{i = 1}^{\infty} A_i\right)\right) + \mu^{*}\left(E \cap \left(\bigcup_{i = 1}^{\infty} A_i\right)^{c}\right)
    \end{align*}
    We see that $\bigcup_{i = 1}^{\infty} A_i \in \cal{A}$. Thus, $\cal{A}$ is a $\sigma$-algebra. Then, using the inequality above, as we see
    \[ \mu^{*}(E) \geq \sum_{i = 1}^{\infty} \mu^{*}(E \cap A_i) + \mu^{*}\left(E \cap \left(\bigcup_{i = 1}^{\infty} A_i\right)\right) \]
    for all $E$, letting $E = \bigcup_{i = 1}^{\infty} A_i$, we see that
    \[ \mu^{*}\left(\bigcup_{i = 1}^{\infty} A_i\right) \geq \sum_{i = 1}^{\infty} \mu^{*}(A_i) \]
    so along with subadditivity, this implies that $\mu^{*}$ when restricted to $\cal{A}$ is countably additive, meaning that $\mu^{*}$ is a measure on $\cal{A}$.\\

    All that is left to show is that $\cal{A}$ contains all of the null sets of $\mu^{*}$ (which are the sets for which $\mu^{*}(A) = 0$). Now, if $\mu^{*}(A) = 0$, then we immediately see
    \[ \mu^{*}(E) \geq \mu^{*}(E \cap A^{c}) = \mu^{*}(E \cap A) + \mu^{*}(E \cap A^{c}) \]
    so thus, $A$ is $\mu^{*}$-measurable!

    $\qed$
\end{proof}

Now, we can define the Lebesgue measure formally using the outer-measure we constructed earlier!

\begin{definition}[\textbf{Definition 5.2: }The Lebesgue Measure]
    The \emph{Lebesgue measure} on $\R$ is the measure given by restricting the outer measure
    \[ \lambda^{*}(A) = \inf\{\sum_{i = 1}^{\infty} \abs{I_{i}} \col A \subseteq \cup I_{i}\} \]
    to the collection of $\lambda^{*}$-measurable sets. The corresponding $\sigma$-algebra $\cal{L}$ is called the \emph{Lebesgue $\sigma$-algebra}.
\end{definition}

Now, we just went through all that work to define a $\sigma$-algebra for an outer measure, but why didn't we just show that the outer measure was a measure on all sets? Well, it turns out we can't always do that, as we will soon show. First, a lemma that will help our construction.
\clearpage

\begin{lemma}[\textbf{Lemma 5.3}]
    If $A \subseteq \R$ and $x \in \R$, then if $A + x = \{a + x \col a \in A\}$, then $\lambda^{*}(A + x) = \lambda^{*}(A)$. Thus, the Lebesgue measure is translation invariant.
\end{lemma}
\begin{proof}{Proof:}
    This is quite obvious from the definition of $\lambda^{*}$.
\end{proof}

\begin{theorem}[\textbf{Theorem 5.4:} Existence of Non-Measurable Sets]
    There exists sets in $\cal{P}(\R)$ that are not Lebesgue measurable.
\end{theorem}
\begin{proof}{Proof:}
    First, define an equivalence relation $\sim$ on $[0, 1]$ by $x \sim y$ if, and only if, $x - y \in \Q$. Now, define $V$ by choosing one representative from each equivalence class of $\sim$.\\

    Now, suppose that $V$ is Lebesgue-measurable. Enumerate $\Q \cap [-1, 1]$ as $q_1, q_2, \cdots$ and let $V_n = V + q_n$. It can be easily shown that all the $V_n$ are disjoint, and by our lemma. $\lambda(V_n) = \lambda(V)$. Thus, we see that
    \[ \lambda\left(\bigcup_{n = 1}^{\infty} V_n\right) = \sum_{n = 1}^{\infty} \lambda(V_n) = \sum_{n = 1}^{\infty} \lambda(V) \]
    Thus, this must be either $0$ or infinity depending on whether $\lambda(V) = 0$ or $\lambda(V) > 0$. We will show that neither can happen.\\

    First, since every $x \in [0, 1]$ must lie in some $V_n$, we see that $[0, 1] \subseteq \bigcup_{n = 1}^{\infty} V_n$, we must have that $\lambda(\cup_{i = 1}^{\infty}V_n) \geq 1$, so $\lambda(V) = 0$ cannot occur. Then, as $V_n \subseteq [-1, 2]$ for all $n$, we must also have that $\lambda(\bigcup_{n = 1}^{\infty} V_n) \leq \lambda([-1, 2]) \leq 3$, so $\lambda(V) > 0$ cannot occur as well. Thus, we have a contradiction, so $V$ must not be measurable.

    $\qed$
\end{proof}

This doesn't look good... What is the point of a function that could possibly fail to give a measure to a set that we want to measure? Well, fortunately, most of the sets we care about are measurable, as we will see next.

\begin{lemma}[\textbf{Proposition 5.5}]
    Every Borel set is Lebesgue measurable.
\end{lemma}
\begin{proof}{Proof:}
    It suffices to show that each interval is Lebesgue measurable. We want to show that, for all $E$, 
    \[ \lambda^{*}(E) \geq \lambda^{*}(E \cap (a, b)) + \lambda^{*}(E \cap (a, b)^{c}) \]
    Let $\epsilon > 0$. By definition of $\lambda^{*}$, there exists a collection $\{I_n\}$ such that $\sum_{n = 1}^{\infty} \abs{I_n} \leq \lambda^*(E) + \epsilon$ and $E \subseteq \bigcup_{n = 1}^{\infty} I_n$. Then, as we have
    \[ \lambda^{*}(E \cap (a, b)) \leq \sum_{i = 1}^{\infty} \lambda^{*}(I_n \cap (a, b)) \leq \sum_{i = 1}^{\infty} \abs{I_n \cap (a, b)} \]
    and
    \[ \lambda^{*}(E \cap (a, b)^{c}) \leq \sum_{i = 1}^{\infty} \abs{I_n \cap [b, \infty)} + \sum_{i = 1}^{\infty} \abs{I_n \cap (-\infty, a]} \]
    we see that
    \begin{align*} 
        \lambda^{*}(E \cap (a, b)) + \lambda^{*}(E \cap (a, b)^{c}) &\leq \sum_{i = 1}^{\infty} \left(\abs{I_n \cap (a, b)} + \abs{I_n \cap (-\infty, a]} + \abs{I_n \cap [b, \infty)}\right)\\
        &= \sum_{i = 1}^{\infty} \abs{I_n}\\
        &\leq \lambda^{*}(E) + \epsilon
    \end{align*}
    as this holds for all $\epsilon > 0$, we have
    \[ \lambda^{*}(E) \geq \lambda^{*}(E \cap (a, b)) + \lambda^{*}(E \cap (a, b)^{c}) \]
    which was to be proved.

    $\qed$
\end{proof}

\begin{lemma}[\textbf{Proposition 5.6}]
    There exists Lebesgue measurable sets that are not Borel. That is
    \[ \cal{B} \subsetneq \cal{L} \subsetneq \P(\R) \]
\end{lemma}
\begin{proof}{Proof Sketch:}
    We won't prove this in it's entirety, but it turns out that $\abs{\cal{B}} = \abs{\R}$, since we can map sequences of $\N$ into $\cal{B}$. Then, it turns out that $\cal{L} > \cal{R}$, so there are quite a bit more Lebesgue measurable sets than Borel measurable sets. 

    $\qed$
\end{proof}
\tocsection{\textbf{9/9/26 Lecture 6: }Properties of the Lebesgue Measure}

\begin{lemma}[\textbf{Proposition 6.1}]
    There exists a set in $\cal{L}$ that is not in $\cal{B}$.
\end{lemma}
\begin{proof}{Proof:}
    First, we assume the fact that $\abs{\cal{B}} = \abs{\R}$.\\

    Now, let $C$ be the middle thirds Cantor set. We know that $\abs{C} = \abs{\R}$. Now, every subset of $C$ is Lebesgue measureable, since $C$ is a null set of $\lambda^{*}$, so $\cal{P}(C) \subseteq \cal{L}$. However, by Cantor's theorem, we see that $\abs{\cal{P}(C)} > \abs{C} = \abs{\R}$, so
    \[ \abs{\cal{L}} \geq \abs{\cal{P}(C)} > \abs{C} = \abs{\R} = \abs{\cal{B}} \]
    so as $\abs{\cal{B}} < \abs{\cal{L}}$, the desired result holds. 

    $\qed$
\end{proof}

It turns out that this construction entirely highlights the issue with $\cal{B}$: it doesn't contain all the null sets of $\lambda$, which we will soon see.

\begin{definition}[\textbf{Definition 6.2: }Complete Metric Spaces]
    A measure space $(X, \cal{A}, \mu)$ is \emph{complete} if $B \in \cal{A}$ and $\mu(B) = 0$ implies that for all $A \subseteq B$, $A \in \cal{A}$. 
\end{definition}

Notice that we can define a measure space $(\R, \cal{B}, \lambda)$ to be a measure space by restricting the Lebesgue measure to $\cal{B}$, but this is not a complete measure space.

\begin{lemma}[\textbf{Proposition 6.3}]
    Let $A \subseteq [0, 1]$ be $\lambda$-measurable. Then,
    \begin{enumerate}
        \item For any $\epsilon > 0$, there exists an open set $\cal{U}$ such that $A \subseteq \cal{U}$ and $\lambda(\cal{U} \backslash A) < \epsilon$.
        \item For any $\epsilon > 0$, there exists a closed set $\cal{C}$ such that $\cal{C} \subseteq A$ and $\lambda{A \backslash \cal{C}} < \epsilon$. 
        \item There are open sets
        \[ \cal{U}_1 \supseteq \cal{U}_2 \supseteq \cal{U}_3 \supseteq \cdots \supseteq A \]
        so that $\lambda\left(\bigcap_{i = 1}^{\infty} \cal{U}_i \backslash A\right) = 0$. 
        \item There are closed sets
        \[ \cal{C}_1 \subseteq \cal{C}_2 \subseteq \cdots \subseteq A \]
        such that $\lambda\left(A \backslash \bigcup_{i = 1}^{\infty} E_{n}\right) = 0$. 
    \end{enumerate}
\end{lemma}
\begin{proof}{Proof:}
    
    \textbf{1)} From the definition of $\lambda$, one can find an open set $\cal{U}$ such that $A \subseteq \cal{U}$ and $\lambda(U) < \lambda(A) + \epsilon$. Then, as $\lambda(\cal{U} \backslash A) = \lambda(\cal{U}) - \lambda(A)$ as $A \subseteq \cal{U}$, we see that $\lambda(\cal{U} \backslash A) < \epsilon$.\\

    \textbf{2)} Now, consider the set $[0, 1] \backslash A$, and using $1$ find an open set $\cal{O} \supseteq [0, 1] \backslash A$ such that $\lambda(\cal{O} \backslash ([0, 1] \backslash A)) < \epsilon$. Then, as 
    \begin{align*}
        \lambda(\cal{O} \backslash ([0, 1] \backslash A)) &= \lambda(\cal{O}) - \left(1 - \lambda(A)\right)\\
        &= \lambda(A) - (1 - \lambda(\cal{O})) \\
        &= \lambda(A) - \lambda([0, 1] \backslash \cal{O})
    \end{align*}
    and we see that $[0, 1] \backslash \cal{O}$ is a closed set contained in $A$, we see that this last expression is equal to $\lambda(A \backslash ([0, 1] \backslash \cal{O})) < \epsilon$.\\

    \textbf{3)} By $1$, there exists open sets $V_n \supseteq A$ such that $\lambda(V_n \backslash A) < 1/n$. Let $U_n = \bigcap_{k = 1}^{n} V_n$. Then, we see that
    \[ U_1 \supseteq U_2 \supseteq \cdots \supseteq A \]
    the $U_n$ are open, and $\lambda\left(\bigcap_{n = 1}^{\infty} V_n \backslash A\right) = \lim_{n \to \infty} \lambda(U_n \backslash A)$, which equals zero as $\lambda(U_n \backslash A) < 1/n$ for each $n$ (as $U_n \subseteq V_n$).\\

    \textbf{4)} Similar proof as \textbf{3)}. 

    $\qed$
\end{proof}

\begin{lemma}[\textbf{Corollary 6.4}]
    $\cal{L}$ is generated by $\cal{B}$ and the null sets of $\lambda$. 
\end{lemma}
\begin{proof}{Proof:}
    Let $A \in \cal{L}$. From the proposition above, there exists a $G_n = \cap_{i = 1}^{\infty} \cal{U}_i \in B$ so that $G_n \supseteq A_n = A \cap [n, n + 1]$ and $\lambda(G_n \backslash A_n) = 0$. Thus, we see that $\bigcup_{i = 1}^{\infty} G_n \in \cal{B}$ and
    \[ A = \left(\bigcup_{i = 1}^{\infty} G_n\right) \backslash \left(\bigcup_{i = 1}^P{\infty} G_n \backslash A_n\right) \]
    so $A$ can be written as the union, complement, and intersection of Borel and null sets. 

    $\qed$
\end{proof}

\begin{definition}[\textbf{Definition 6.5:} $G_{\delta}$ and $F_{\sigma}$ Sets]
    A countable intersection of open sets is a $G_{\delta}$ set. A countable union of closed sets is a $F_{\sigma}$ set. 
\end{definition}

\subsection*{The Lebesgue-Stieljes Measure}

Let $\cal{C}$ be a collection of intervals of the form $(a, b]$ and let $\alpha : \R \to \R$ be increasing and right-continuous. Then, define $\ell((a, b]) = \alpha(b) - \alpha(a)$. Now, for $E \subseteq \R$, define
\[ m^{*}(E) = \inf\{\sum_{n = 1}^{\infty} \ell(A_n) \col A_n \in \cal{C} \text{ and } E \subseteq \bigcup_{n = 1}^{\infty} A_n\} \]
For example, if $\alpha = x$, then we get the Lebesgue outer measure, and if $\alpha$ is zero for all $x < 0$ and $1$ for all $x \geq 0$, we see that $m^{*}$ is the point mass outer measure. 

\begin{lemma}[\textbf{Proposition 6.6}]
    The function $m^{*} : \cal{P} \to [0, \infty]$ is an outer measure. Let $m$ be the measure defined by the restriction of $m^{*}$ to $m^{*}$-measurable sets. Then, every Borel set is $m^{*}$-measurable and $m^{*}((a, b]) = \ell(a, b]$. 
\end{lemma}
\begin{proof}{Proof:}
    This is the same proof for $\lambda^{*}$, but we replace $\abs{(a, b)} = b - a$ with $\alpha(b) - \alpha(a)$. 
\end{proof}
\tocsection{\textbf{9/11/26 Lecture 7:} L-S Measures and More}

\textbf{Remark:} Before we prove that L-S measures satisfy the conclusions of Proposition 6.3, recall that given a L-S outer measure $m^{*}$, all Borel sets are $m^{*}$-measurable, but the $\sigma$-algebra of $m^{*}$-measurable sets is not necessarily the Lebesgue $\sigma$-algebra. 

\begin{lemma}[\textbf{Proposition 7.1}]
    Let $m$ be an L-S measure. Then, if $A \subseteq [0, 1]$, then for all $\epsilon > 0$, there exists an open set $\cal{U} \supseteq A$ such that $m(\cal{U}\backslash A) < \epsilon$. With this, the rest of the conclusions of Proposition 6.3 also follow. 
\end{lemma}
\begin{proof}{Proof:}
    By definition, we can write $A \subseteq \bigcup_{n} (a_n, b_n]$ such that $\sum_{i = 1}^{\infty} \ell(a_n, b_n] \leq m(A) + \frac{\epsilon}{2}$. Now, for each $n$, notice that
    \[ m(a_n, b_n] = \lim_{m \to 0} m(a_n, b_n + \frac{1}{m}) \]
    as we see that $(a_n, b_n + \frac{1}{m}) \downarrow (a_n, b_n]$. Thus, we can find a $c_n$ such that $m(a_n, c_n) \leq m(a_n, b_n] + \frac{\epsilon}{2^{n + 1}}$. Thus, we have that
    \begin{align*}
        m(\bigcup_{n = 1}^{\infty} (a_n, c_n) \backslash A) &= m(\bigcup_{n = 1}^{\infty} (a_n, c_n)) - m(A)\\
        &\leq \sum_{i = 1}^{\infty} m((a_n, c_n)) - m(A)\\
        &\leq \sum_{i = 1}^{\infty} \left(m(a_n, b_n] + \frac{\epsilon}{2^{n + 1}}\right) - m(A)\\
        &= \sum_{i = 1}^{\infty} \ell(a_n, b_n] + \frac{\epsilon}{2} - m(A)\\
        &\leq m(A) + \frac{\epsilon}{2} + \frac{\epsilon}{2} - m(A)\\
        &= \epsilon
    \end{align*}
    which was to be shown as $\bigcup_{n = 1}^{\infty} (a_n, c_n)$ is open. 

    $\qed$
\end{proof}

Now, in these past few lectures we've developed how to go from an outer-measure to a measure on a $\sigma$-algebra, but what if we know a bit more about our ``measuring function''. What if it's already pretty close to a measure? We will explore that now.

\begin{definition}[\textbf{Definition 7.2:} Algebras]
   Given a set $X$, a \emph{algebra} $\cal{A}_0$ is a collection of subsets such that
   \begin{enumerate}
       \item $\emptyset, X \in \cal{A}_0$
       \item $A^{c} \in \cal{A}_{0}$ if $A \in \cal{A}_{0}$
       \item If $A_1, \ldots, A_n \in \cal{A}_0$, then $\bigcup_{i = 1}^{n} A_i \in \cal{A_0}$
   \end{enumerate} 
\end{definition}

\begin{definition}[\textbf{Definition 7.3:} Premeasures]
    A \emph{premeasure} $\ell$ on an algebra $\cal{A}_0$ is a map $\ell : \cal{A}_0 \to [0, \infty]$ such that
    \begin{enumerate}
        \item $\ell(\emptyset) = 0$
        \item If $A_1, \ldots \in \cal{A}_0$ and $\bigcup_{i = 1}^{\infty} A_n \in \cal{A}_0$, and the $A_i$ are pairwise disjoint, then
        \[ \ell\left(\bigcup_{n = 1}^{\infty} A_n\right) = \sum_{n = 1}^{\infty} \ell(A_n) \]
    \end{enumerate}
\end{definition}
Now, one last definition that we will need for our extension theorem.

\begin{definition}[\textbf{Definition 7.4:} Finite and $\sigma$-Finite Measures]
    A measure $\mu$ on $X$ is \emph{finite} if $\mu(X) < \infty$. A measure $\mu$ on $X$ is \emph{$\sigma$-finite} if there exists $X_n \in \cal{A}$ such that $X = \bigcup_{n = 1}^{\infty} X_n$ and $\mu(X_n) < \infty$. 
\end{definition}

Thus, for example $\lambda$ on $\R$ is not finite but it is $\sigma$-finite. Most spaces (and measures) that we will be concerned with are $\sigma$-finite. Now, we move onto our extension theorem.

\begin{theorem}[\textbf{Theorem 7.5:} Caratheodory Extension Theorem]
    Let $\cal{A}_0$ be an algebra on $X$ and let $\ell$ be a premeasure on $\cal{A}_0$. Define
    \[ \mu^{*}(E) = \inf\left\{\sum_{i = 1}^{\infty} \ell(A_n) \col A_n \in \cal{A_0}, E \subseteq \bigcup_{n = 1}^{\infty} A_n \in \cal{A}_0\right\} \]
    Then,
    \begin{enumerate}
        \item $\mu^{*}$ is an outer measure.
        \item $\mu^{*}(A) = \ell(A)$ if $A \in \cal{A}_0$.
        \item Every $\mu^{*}$ null set is $\mu^{*}$-measurable and every $A \in \cal{A}_0$ is $\mu^{*}$-measurable.
        \item If $\ell$ is $\sigma$-finite, then $\mu^{*}$ is the unique extension of $\ell$ to $\sigma(\cal{A}_0)$.
    \end{enumerate}
\end{theorem}
\begin{proof}{Proof:}
    In this lecture we will prove the first three statements, and we will prove the fourth one in the next lecture.\\

    \textbf{1)} This is the same proof as for our other extension theorems.\\

    \textbf{2)} Suppose that $E \in \cal{A}_0$ and $E \subseteq \bigcup_{i =1}^{\infty} A_n$ for $A_n \in \cal{A}_0$. Define $B_n = E \cap \left(A_n \backslash \bigcup_{i = 1}^{n - 1} A_i\right)$. We see that $B_n \in \cal{A}$ for each $n$.\\

    Then, we see that $\ell(E) = \sum_{n} \ell(B_n) \leq \sum_{n} \ell(A_n)$, so taking infimums over such collections, we see that $\ell(E) \leq \mu^{*}(E)$, and as we clearly have that $\mu^{*}(E) \leq \ell(E)$, we have the desired property.\\

    \textbf{3)} Let $A \in \cal{A}_0$ and let $E \subseteq X$. WE want to show that $\mu^{*}(E) \leq \mu^{*}(E \cap A) + \mu^{*}(E \cap A^{c})$. Choose $B_n \in \cal{A}_0$ so that $\bigcup_{n} B_n \in \cal{A_0}$, $E \subseteq \bigcup_{n} B_n$, and
    \[ \sum_{n} \ell(B_n) \leq \mu^{*}(E) + \epsilon \]
    Then, we see that
    \begin{align*}
        \mu^{*}(E) + \epsilon &\geq \sum_{n} \ell(B_n \cap A) + \ell(B_n \cap A^{c})\\
        &= \sum_{n} \ell(B_n \cap A) + \sum_{n} \ell(B_n \cap A^{c})\\
        &\geq \mu^{*}(E \cap A) + \mu^{*}(E \cap A^{c})
    \end{align*}
    where the nontrivial step of the proof is that $\ell(B_n) = \ell(B_n \cap A) + \ell(B_n \cap A^{c})$ as $A \in \cal{A}_0$. Thus, as this holds for all $\epsilon$, we have what was to be shown.

    $\qed$
\end{proof}
\tocsection{\textbf{9/14/26 Lecture 8}: Intro to the Lebesgue Integral}

First, we finish up the proof for Caratheodory's theorem seen above.\\

\begin{proof}{Proof of 4)}
    Let $\nu$ be a measure on $\sigma(\cal{A}_0)$ such that $\nu(A) = \ell(A)$ for all $A \in \cal{A_0}$. WE want to show that $\nu(A) = \mu^{*}(A)$ for all $A \in \sigma(A_n)$.\\

    First, we will consider the case when $\mu^{*}$ is finite. Fix $A \in \sigma(\cal{A}_0)$ and let $A \subseteq \bigcup_{n} A_n$ with $A_n \in \cal{A_0}$. Then, as 
    \[ \nu(A) \leq \sum_{n = 1}^{\infty} \nu(A_n) = \sum_{n = 1}^{\infty} \ell(A) \]
    taking infs over such covers, we see that $\nu(A) \leq \mu^{*}(A)$.\\

    Now, we want to show that $\nu$ and $\mu^{*}$ agree on sets of the form $\bigcup_{n} A_n$ for $A_n \in \cal{A}_0$. Let
    \[ B_n = \bigcup_{j = 1}^{n} A_j \]
    then, we see that
    \[ \mu^{*}(\cup_{i = 1}^{\infty} A_i) = \lim_{n \to \infty} \mu^{*}(B_n) = \lim_{n \to \infty} \nu(B_n) = \nu(\cup_{i = 1}^{\infty} A_i) \]
    Now, choose $A_n$ so that $A \subseteq \bigcup_{i = 1}^{\infty} A_n$ and $A_n \in \cal{A}_0$ such that $\mu^{*}(A) + \epsilon \geq \sum_{n} \mu^{*}(A_n)$. Then, 
    \begin{align*}
        \mu^{*}(A) + \epsilon &\geq \sum_{n} \mu^{*}(A_n)\\
        &\geq \mu^{*}\left(\bigcup_{i = 1}^{\infty} A_n\right)
    \end{align*}
    thus, we see that $\epsilon \geq \mu^{*}\left(\bigcup_{i = 1}^{\infty} A_n\right) - \mu^{*}(A) = \mu^{*}(\bigcup_{i = 1}^{\infty} A_n \backslash A)$ (this step requires that $\mu^{*}$ is finite). Then, we see that
    \begin{align*}
        \mu^{*}(A) \leq \mu^{*}(\bigcup_{i = 1}^{\infty} A_n) &= \nu(\bigcup_{i = 1}^{\infty} A_n)\\
        &= \nu(A) + \nu(\bigcup_{i = 1}^{\infty} A_n \backslash A)\\
        &\leq \nu(A) + \mu^{*}(\bigcup_{i =1}^{\infty} A_n \backslash A)\\
        &< \nu(A) + \epsilon
    \end{align*}
    where the second to last step came from above. Thus, $\mu^{*}(A) \leq \nu(A) + \epsilon$ for all $\epsilon$, so $\mu^{*}(A) \leq \nu(A)$.\\

    Now, if $\mu^{*}$ is $\sigma$-finite, then we can write $X = \bigcup_{n = 1}^{\infty} X_n$ with $X_n \uparrow X$ with $\mu^{*}(X_n) < \infty$. Then, if we consider the measures $\mu^{*}\res{X_n}$ and $\nu\res{X_n}$ (where $\mu\res{X_n}(A) = \mu(X_n \cap A)$), we see that as $\mu^{*}\res{X_n}$ and $\nu\res{X_n}$ are finite measures, they have to agree by the argument above. Then, for $A \in \sigma(\cal{A_0})$, we see that
    \[ \mu^{*}(A) = \lim_{n \to \infty} \mu^{*}(A \cap X_n) = \lim_{n \to \infty} \nu(A \cap X_n) = \nu(A) \]
    which was desired. 

    $\qed$
\end{proof}

Now, lets move on to functions. As a motivation, consider $C[0, 1]$ with the metric
\[ d(f, g) = \int_{0}^{1} \abs{f(x) - g(x)}dx \]
This is nice until you realize that it doesn't make $C[0, 1]$ complete! Consider the functions as described below

\begin{center}
    \begin{tikzpicture}
        \draw[thick] (-2, 0) -- (2, 0);
        \draw[thick] (-2, 0.1) -- (-2, -0.1);
        \draw[thick] (2, 0.1) -- (2, -0.1);


        \node at (-2, -0.5) {$0$};
        \node at (2, -0.5) {$1$};

        \draw[thick, red] (-2, 0) -- (-0.2, 0);
        \draw[thick, red] (-0.2, 0) -- (0, 2);
        \draw[thick, red] (0, 2) -- (2, 2);
        \node at (-0.2, -0.5) {$0.5 - 1/n$};
        \draw[thick] (-0.2, -0.1) -- (-0.2, 0);
    \end{tikzpicture}
\end{center}

These turn out to be Cauchy under this metric, but they converge to the function
\[ f(x) = \chi_{[1/2, 1]} \]
where $\chi_{A}$ is the characteristic function of $A$ which is $1$ on $A$ and $0$ off of $A$. The problem is that $\chi_{[1/2, 1]}$ is not continuous, so it cannot be the limit. This example is somewhat tame in that the limiting function is still Riemann integrable, but there are some Cauchy sequences in $C[0, 1]$ that converge to non Riemann integrable functions, so we want to see how to fix that.\\

First, consider $\chi_{\Q \cap [0, 1]}$. This is not Riemann integrable as $U(\chi_{\Q \cap [0, 1]}, P) = 1$ and $L(\chi_{\Q \cap [0, 1]}, P) = 0$ for any partition. However, maybe we should still try to make sense of it's integral.\\

\textbf{Question:} What should $\int \chi_{\Q \cap [0, 1]} dx$ be?

\textbf{Intuitive Answer} Since $\int$ is the area under a function, $\int \chi_{E}$ should be the measure of $E$ for any set $E$. That is, $\int \chi_{E} = \lambda(E)$. Then, as $\lambda(\Q \cap [0, 1]) = 0$, it would be natural to expect that $\int \chi_{\Q \cap [0, 1]} dx = 0$. How do we formalize this? First, we start by defining the integral on characteristic functions as we did above.\\

\begin{definition}[\textbf{Definition 8.1:} Lebesgue Integral of a Characteristic Function]
    Given a measure space $(X, \cal{A}, \mu)$, define the \emph{Lebesgue integral} of a characteristic function $\chi_{A}$ for $A \in \cal{A}$ to be
    \[ \int \chi_{A} d\mu = \mu(A) \]
\end{definition}

\begin{definition}[\textbf{Definition 8.2:} Simple Functions]
    A \emph{simple function} is a function of the form
    \[ f = \sum_{i = 1}^{\infty} c_n \chi_{A_n} \]
\end{definition}
Then, as we want our integral to be linear as well, we can define
\[ \int \sum_{n = 1}^{\infty} a_n \chi_{A_n} d\mu = \sum_{n = 1}^{\infty} a_n\mu(A_n) \]
for $A_n \in \cal{A}$.\\

Now, heres a definition that seems random right now, but we will soon see it describes exactly what we want.

\begin{definition}[\textbf{Definition 8.3:} Measurable Functions]
    A function $f : X \to \R$ is \emph{measurable} if $\{x \in X \col f(x) > a\} = f^{-1}(a, \infty) \in \cal{A}$ for all $a \in \R$.
\end{definition}

Now, we see, considering some $\chi_{A}$, that $f^{-1}(a, \infty) = \emptyset$ if $a \geq 1$, $f^{-1}(a, \infty) = A$ if $0 \leq a < 1$, and $f^{-1}{a, \infty} = X$ if $a < 0$. Thus, we see that $\chi_{A}$ is measurable if $A$ is a measurable set!
\tocsection{\textbf{9/18/26 Lecture 9: }Measurable Functions}

\begin{lemma}[\textbf{Proposition 9.1}]
    The following are equivalent.
    \begin{enumerate}
        \item $f$ is measurable.
        \item $f^{-1}((-\infty, a]) \in \cal{A}$ for all $a$.
        \item $f^{-1}((-\infty, a)) \in \cal{A}$ for all $a$.
        \item $f^{-1}([a, \infty)) \in \cal{A}$ for all $a$. 
    \end{enumerate}
\end{lemma}
\begin{proof}{Proof:}
    The idea is that $f^{-1}$ preserves complements, unions, and intersections.\\

    $1) \to 2)$ Since $f^{-1}((a, \infty))^{c} = f^{-1}((-\infty, a])$, we see that $f^{-1}((-\infty, a]) \in \cal{A}$.\\

    $2) \to 3)$ Writing $(-\infty, a) = \bigcup_{n = 1}^{\infty} (-\infty, a - 1/n]$, we see that
    \[ f^{-1}((-\infty, a)) = \bigcup_{n = 1}^{\infty} f^{-1}((-\infty, a - 1/n]) \in \cal{A} \]

    $3) \to 4)$ Since $f^{-1}((-\infty, a))^{c} = f^{-1}([a, -\infty))$, we see that $f^{-1}([a, \infty)) \in \cal{A}$.\\

    $4) \to 1)$ Writing $(a, \infty) = \bigcup_{n = 1}^{\infty} [a + 1/n, \infty)$, we see that
    \[ f^{-1}((a, \infty)) = \bigcup_{n = 1}^{\infty} f^{-1}([a + 1/n, \infty)) \in \cal{A} \]

    $\qed$
\end{proof}

\begin{lemma}[\textbf{Proposition 9.2}]
    Suppose that $X$ is a metric space and let $\cal{A}$ be a $\sigma$-algebra on $X$ that contains every open set. Then, every continuous function $f : X \to \R$ is measurable.
\end{lemma}
\begin{proof}{Proof:}
    Since $(a, \infty)$ is open, $f^{-1}((a, \infty))$ is open, so it lies in $\cal{A}$. 

    $\qed$
\end{proof}

\begin{lemma}[\textbf{Proposition 9.3}]
    Suppose that $f, g : X \to \R$ are measurable. Then, so are the following:
    \begin{enumerate}
        \item $f + g$
        \item $cf$ for $c \in \R$
        \item $fg$
        \item $\max(f, g)$
        \item $\min(f, g)$
    \end{enumerate}
\end{lemma}
\begin{proof}{Proof:}
    
    $1)$ We will show that
    \[ (f + g)^{-1}(-\infty, a) = \bigcup_{r \in \Q}(f^{-1}(-\infty, a - r) \cap g^{-1}(-\infty, r]) \]
    which is sufficient as the LHS lies in $\cal{A}$. First, it is quite clear that we have the backwards inclusion. Now, let $x \in (f + g)^{-1}(-\infty, a)$. Then $f(x) + g(x) < a$, so $f(x) < a - g(x)$. Then, by density of $\Q$, we see there exists an $r \in \Q$ for which $f(x) < a - r \leq a - g(x)$, from which we can also include that $g(x) \leq r$. Thus, we see that $x \in f^{-1}(-\infty, a - r) \cap g^{-1}(-\infty, r]$, so we have the forward inclusion.\\

    $2)$ If $c = 0$, this is trivial, so suppose not. Then, we see that
    \[ (cf)^{-1}(-\infty, a) = \{x \col cf(x) < a\} = \left\{\begin{matrix} \{x \col f(x) < \frac{a}{c}\} & c > 0 \\ \{x \col f(x) > \frac{a}{c}\} & c < 0\end{matrix}\right. \]
    which are measurable sets in all cases.\\

    $3)$ First, notice that $f^2$ is measurable since, for $a < 0$, $(f^2)^{-1}(a, \infty) = X$, and for $a \geq 0$, 
    \[ (f^2)^{-1}(a, \infty) = \{x \col f(x) > \sqrt{a}\} \cup \{x \col f(x) < -\sqrt{a}\} \]
    Then, $fg$ is measurable as we can write
    \[ fg = \frac{1}{2}\left[(f + g)^2 - f^2 - g^2\right] \]

    $4)$ We see that
    \[ (\max(f, g))^{-1}(a, \infty) = f^{-1}(a, \infty) \cup g^{-1}(a, \infty) \in \cal{A} \]

    $5)$ We see that
    \[ (\min(f, g))^{-1}(-\infty, a) = f^{-1}(-\infty, a) \cup g^{-1}(-\infty, a) \in \cal{A} \]

    $\qed$

\end{proof}

\begin{lemma}[\textbf{Proposition 9.4}]
  If $f_n$ is a sequence of measurable functions, then $\sup_{n}f_n, \inf_{n}f_n, \limsup_{n}f_n$ and $\liminf_{n}f_n$ are measurable.  
\end{lemma}
\begin{proof}{Proof:}
    We can write
    \[ (\sup_{n} f_n)^{-1}(a, \infty) = \bigcup_{n \geq 1} f_{n}^{-1}(a, \infty) \in \cal{A} \]
    Then, as $\inf_{n}f_n = -\sup_{n}(-f_n)$, we see that $\inf_{n} f_n$ is also measurable. Then, as 
    \[ \limsup_{n} f_n = \inf_{n} \sup_{k \geq n} f_k \quand \liminf_{n} f_n = \sup_{n} \inf_{k \geq n} f_k \]
    we see that $\limsup_{n} f_n$ and $\liminf_{n}f_n$ are measurable. 
\end{proof}

Also, note that if $\lim_{n} f_n$ exists, then it is equal to $\limsup_{n} f_n$, so it is also measurable.\\

\begin{definition}[\textbf{Definition 9.5:} Equality Almost Everywhere]
    We say that two measurable functions $f$ and $g$ are \emph{equal almost everywhere}, denoted $f \overset{\tt{a.e.}}{=} g$, if $\{x \col f(x) \ne g(x)\}$ has measure $0$. 
\end{definition}
Note that the above set is measurable as it is equal to $\{x \col (f - g)(x) \geq 0\} \cap \{x \col (f - g)(x) \leq 0\}$. An example of the above definition is that $\chi_{\Q}$ is equal to $0$ almost everywhere because $\Q$ is the set on which they differ.\\

\begin{definition}[\textbf{Definition 9.6:} Convergence Almost Everywhere] 
    If $f_n$ is a sequence of measurable functions, we say \emph{$f_n$ converges to $f$ almost everywhere}, denoted $f_n \to f \tt{ a.e.}$, if
    \[ \{x \col \lim_{n \to \infty} f_n(x) \ne f(x)\} \]
    has measure zero. 
\end{definition}

\begin{definition}[\textbf{Definition 9.7:} Borel Measurability]
    Let $X$ be a metric space and $\cal{B}(X)$ be it's Borel $\sigma$-algebra. Then, $f : X \to \R$ is \emph{Borel measurable} if it is measurable with respect to $\cal{B}(X)$. That is, $f^{-1}(a, \infty) \in \cal{B}(X)$ for all $a$. 
\end{definition}
\tocsection{\textbf{9/16/26 Lecture 10:} More on Lebesgue Integrals}

\begin{lemma}[\textbf{Proposition 10.1}]
    If $f : \R \to \R$ is monotone, then it is Borel-measurable. 
\end{lemma}
\begin{proof}{Proof:}
	WLOG, assume that $f$ is increasing. We want to show that $f^{-1}(a, \infty)$ is an interval. Now, let
    \[ x_0 = \sup\{x \col f(x) \leq a\} \]
    We have to cases. If $f(x_0) \leq a$, then $f^{-1}(a, \infty) = (x_0, \infty)$. Similarly, if $f(x_0) > a$, then $f^{-1}(a, \infty) = [x_0, \infty]$. 
\end{proof}

\begin{lemma}[\textbf{Proposition 10.2}]
    Let $(X, \cal{A})$ be a measurable space. If $f : X \to \R$ is measurable, then $f^{-1}(A) \in \cal{A}$ whenever $A \in \cal{B}$. 
\end{lemma}
\begin{proof}{Proof:}
    Let $\cal{C} = \{A \subseteq \R \col f^{-1}(A) \in \cal{A}\}$. We see that $\cal{C}$ contains all subsets of the form $(a, \infty)$. Now, we see that $\R$ and $\emptyset \in \cal{C}$, it is closed under complements, and is closed under unions (as $f^{-1}$ distributes across them). Thus, we see that $\sigma(\{(a, \infty)\}) \subseteq \cal{C}$, so $\cal{B} \subseteq \cal{C}$. 

    $\qed$
\end{proof}

\begin{lemma}[\textbf{Proposition 10.3}]
    Let $f \geq 0$ and measurable. Then, there exists a sequence of simple functions $s_n \geq 0$ such that $s_n \uparrow f$. That is, $0 \leq s_1 \leq s_2 \leq \cdots \leq f(x)$ and $\sup_{n} s_n = f$. 
\end{lemma}
\begin{proof}{Proof:}
    For each $n$, define
    \[ I_{j, n} = \left[\frac{j - 1}{2^{n}}, \frac{j}{2^{n}}\right) \]
    now for $j = 1, \ldots, n2^{n}$, let
    \[ s_n = \sum_{j = 1}^{n2^{n}} \frac{j - 1}{2^{n}}\chi_{f^{-1}(I_{j, n})} + n\chi_{f^{-1}[n, \infty)} \]
    Now, it can be checked that $s_n$ is increasing, and that for each $x$, $s_n(x) \uparrow f(x)$. 
\end{proof}
Now, we will work towards proving Lusin's theorem.



\begin{definition}[\textbf{Definition 10.4: }Support of a Function]
    The \emph{support} of a function $f$ is
    \[ \text{supp}(f) = \overline{\{x \col f(x) \ne 0\}} \]
\end{definition}

Now, one way to define continuity is that, for $f : X \to \R$, $f$ is continuous if $f(x_n) \to f(x)$ for a sequence $x_n \to x$.\\

If $Y \subseteq X$, then $f \res{Y} : Y \to \R$ is given by $f\res{Y}(x) = f(x)$. If we have a function $f$ defined only on $Y$, then continuity means that $g(x_n) \to g(x)$ but only for $x_n \to x$ contained in $Y$. 

\begin{example}{\textbf{Example 1}}
    Consider $\chi_{\Q} : \R \to \R$. This is not continuous, but we see that $\chi_{\Q} \res{\Q} : \Q \to \R$ is continuous as it is the constant function $1$. 
\end{example}

\begin{lemma}[\textbf{Lemma 10.5}]
    If $A \subseteq \R$, define the distance from a point to a set as
    \[ d(x, A) := \inf\{d(x, y) \col y \in A\} \]
    This is a continuous function, and if $A$ is closed, then $f(x) = 0$ if, and only if, $x \in A$. 
\end{lemma}

Now, we state and prove Lusin's Theorem.

\begin{theorem}[\textbf{Theorem 10.6: }Lusin's Theorem]
    Let $f : [0, 1] \to \R$ be Lebesgue measurable. Then, for any $\epsilon > 0$, there exists a closed set $E \subseteq [0, 1]$ such that $\lambda([0, 1] \backslash E) < \epsilon$ and $f\res{E}$ is continuous. 
\end{theorem}
\begin{proof}{Proof:}
    First, we start by proving it for characteristic functions. Let $f = \chi_{A}$. Then, we can find $K \subseteq A \subseteq U$ with $K$ closed and $U$ open so that $\lambda(A \backslash K) < \epsilon/2$ and $\lambda(U \backslash A) < \epsilon/2$.\\

    Let $E = K \cup (U^{c}) \cap [0, 1]$. We see that $E$ is closed, so we need to show that $f\res{E}$ is continuous. Let 
    \[ g(x) = \frac{d(x, U^{c})}{d(x, K) + d(x, U^{c})} \]
    This function is continuous because it's components are continuous and $d(x, K) + d(x, U^{c}) = 0 \implies d(x, K) = 0$ and $d(x, U^{c}) = 0$, which isn't possible as $K \cap U^{c} = \emptyset$ and they are both closed.\\

    Then, notice that
    \[ g\res{E}(x) = \left\{\begin{matrix} 0 & x \in U^{c} \\ 1 & x \in K\end{matrix}\right. = f(x) \]
    so as $g\res{E}$ is continuous, we see that $f$ is continuous. Also, notice that 
    \[ \lambda([0, 1] \backslash E) = \lambda(K^{c} \cap U \cap [0, 1]) \leq \lambda(K^{c} \cap U) = \lambda(U \backslash K) < \epsilon \]

    Now, we will prove it for simple functions. Let 
    \[ f = \sum_{n = 1}^{N} a_n\chi_{A_n} \]
    Now, by the previous step, there exist for each $n$, a $K_n$ so that $\lambda([0, 1] \backslash K_n) < \frac{\epsilon}{n}$, and $\chi_{A_n} \res{K_n}$ is continuous. Then, let
    \[ E = \bigcap_{n = 1}^{N} K_n \]
    We see that $E$ is closed, $\mu([0, 1] \backslash E) < \epsilon$, and as each $\chi_{A_n}$ is continuous on $E$, we see that $f\res{E}$ is continuous.\\

    We will finish this proof next time.-
\end{proof}
\tocsection{\textbf{9/21/26 Lecture 11:} More on Integrals}

We first start by finishing the proof of Lusin's theorem. \\

\begin{proof}{Proof of Lusin's Theorem:}
    \vspace{0.2cm}

    Now, we will prove this for positive functions $f$ bounded above by $N$. Let 
    \[ \tilde{s}_n = \sum_{j = 1}^{N2^{N} + 1} \frac{j - 1}{2^{n}}\chi_{f^{-1}[\frac{j - 1}{2^{n}}, \frac{j}{2^{n}})} \]
    This is a similar construction as before, but as $N$ is bounded, we don't need the extra term at the end. Let $f_n = \tilde{s}_{n + 1} - \tilde{s}_{n}$. This is still a simple function, and satisfies
    \[ \abs{f_n(x)} < \frac{1}{2^{n}} \]
    Then, as $f_n$ are simple functions, we can find sets $E_n$ with $\lambda([0, 1] \backslash E_n) < \frac{\epsilon}{2^{n + 1}}$. Now, we see that $f_n \res{E_{n}}$ is continuous, and notice that
    \[ \tilde{s}_0(x) + \sum_{j = 1}^{n} f_{j}(x) = \tilde{s}_{n + 1}(x) \]
    Now, letting $E = \bigcap_{n = 1}^{\infty} E_n$, then we see that each $f_n$ is continuous on $E$, and we see that $(\sum_{j = 1}^{\infty} f_n(x) - \tilde{s}_0(x))$ converges uniformly to $f$ by the Weierstrass $M$-test. Thus, $f$ is continuous on $E$, as each $f_n$ is continuous on $E$. \\

    Since we see that
    \begin{align*}
        \lambda\left([0, 1] \backslash \bigcap_{i = 1}^{n} E_n\right) \leq \sum_{i = 1}^{\infty} \lambda([0, 1] \backslash E_n) < \epsilon
    \end{align*}
    we are done.
    \vspace{0.4cm}

    Now, suppose $f$ is positive, but not necessarily bounded. Let
    \[ B_{N}(x) = \{x \col f(x) \leq N\} \]
    Notice that $[0, 1] = \bigcup_{n = 1}^{\infty} B_N$. Now, we can find an $N$ for which
    \[ \lambda([0, 1] \backslash B_n) < \frac{\epsilon}{3} \]
    and further find a closed $K \subseteq B_N$ so that
    \[ \lambda(B_N \backslash K) < \epsilon/3 \]
    Consider $f\chi_{K}$. This function is positive and bounded, so we can find $E \subseteq [0, 1]$ so that $f\chi_{K} \res{E}$ is continuous and $\lambda([0, 1] \backslash E) < \frac{\epsilon}{3}$. Also, notice that $f \res{K} = f\chi_{K}\res{K}$. \\

    Then, we see that $\lambda([0, 1] \backslash (E \cap K)) < \epsilon$, and $f \res{E \cap K}$ is continuous as it agrees with $f\chi_{K}$ on $E \cap K$.\\

    Finally, let $f$ be any measurable function. Then, write $f = f_{+} - f_{-}$ where $f_{\pm} = (\pm f) \vee 0$. Then, apply the previous step to $f_{+}$ and $f_{-}$, and the intersection of the sets on which they are continuous gives us a set on which $f$ is continuous.
    
    $\qed$
\end{proof}

Here's something that comes directly from this theorem.

\begin{lemma}[\textbf{Corollary 11.1}]
    Let $f : [0, 1] \to \R$ be Lebesgue measurable. For every $\epsilon > 0$, there is a continuous function $g : [0, 1] \to \R$ so that
    \[ \lambda(\{x \col f(x) \ne g(x)\}) < \epsilon \]
\end{lemma}
\begin{proof}{Proof:}
    By Lusin's theorem, we can find a closed set $E \subseteq [0, 1]$ so that $f\res{E}$ is continuous and $\lambda([0, 1] \backslash E) < \epsilon$.\\

    Now, I want to find a $g$ that agrees with $f$ on $E$ that is still continuous on its complement. Now, note that
    \[ [0, 1] \backslash E = \bigcup_{n = 1}^{\infty} (a_n, b_n) \]
    now, on $(a_n, b_n)$, we do linear interpolation. We know that $g(b_n) = f(b_n)$ and $g(a_n) = f(a_n)$. Then, on $(a_n, b_n)$, let $g$ be the line between $(a_n, f(a_n))$ and $(b_n, f(b_n))$. That is, let
    \[ g\res{(a_n, b_n)} = f(a_n) + \frac{f(b_n) - f(a_n)}{b_n - a_n}(x - a_n) \]

    Now, the claim is that $g$ is in fact continuous. We will not prove this.

    $\qed$
\end{proof}

\subsection*{Lebesgue Integral}
\sectionspace

Let $(X, \cal{A}, \mu)$ be a measure space. We will define the Lebesgue integral again, but now with more tools. 

\begin{definition}[\textbf{Definition 11.2: }Lebesuge Integral for Simple Functions]
    For a simple function
    \[ s = \sum_{j = 1}^{n} a_j\chi_{E_j},\ E_j \subseteq \cal{A} \]
    define
    \[ \int s d\mu = \sum_{i = 1}^{\infty} a_j\mu(E_j) \]
\end{definition}

Note that this sum can be infinite if $\mu(E_j) = \infty$, but if $a_j = 0$ and $\mu(E_j) = \infty$, then we take by convention that $0 \cdot \infty = 0$.\\

Now, note that for any simple function $s$, there are many ways to represent it as a sum of characteristic functions. However, luckily it turns out that the above definition is independent of the representation of $s$.\\

\begin{definition}[\textbf{Definition 11.3:} Lebesgue Integral]
    Let $f \geq 0$ be measurable. Then, define
    \[ \int f d\mu = \sup\left\{\int s d\mu \col s \leq f,\ s \tt{ is simple}\right\} \]
    Now, if $f$ is not necessarily positive, then let $f = f^{+} - f^{-}$ and define
    \[ \int f d\mu = \int f^{+}d\mu - \int f^{-}d\mu \]
    but only when not both $\int f^{+}d\mu$ and $\int f^{-}d\mu$ are infinite. 
\end{definition}

\begin{lemma}[\textbf{Proposition 11.4}]
    Assume $f$ and $g$ are measurable. The following are true about Lebesgue integrals
    \begin{enumerate}
        \item If $0 \leq a \leq f \leq b$ for $x \in X$ and $\mu(X) < \infty$, then
        \[ a\mu(X) \leq \int f d\mu \leq b\mu(X) \]
        \item If $0 \leq f(x) \leq g(x)$
        \[ 0 \leq \int fd\mu \leq \int gd\mu \]
        \item For any constant $c$, 
        \[ \int cf d\mu = c \int f d\mu \]
        \item If $\mu(A) = 0$, then
        \[ \int f\chi_{A}d\mu = 0 \]
    \end{enumerate}
\end{lemma}
\begin{proof}{Proof:}
    
    \textbf{(1)} Note that the simple function $a\chi_{X} \leq f(x)$, and as $\int a\chi_{X} = a\mu(X)$, it is immediately clear that $a\mu(X) \leq \int f d\mu$ from the definition. Now, let $s \leq f$ be a simple function with $s = \sum_{j = 1}^{n} a_j\chi_{A_j}$. We want to show that
    \[ \int s d\mu \leq b\mu(X) \]
    Now, we want to make these $A_j$ disjoint. Since $s$ takes on finitely many values, we can write
    \[ s = \sum_{j = 1}^{m} b_j\chi_{B_j} \]
    for disjoint $B_j$. Now, as $b_j \leq b$ for each $j$, then we see that $s \leq f \leq b$ and 
    \[ \int \sum_{j = 1}^{m} b_j\chi_{b} \leq b\sum_{j = 1}^{m}\mu(B_j) \leq b\mu(X) \]

    $\qed$
\end{proof}
\tocsection{\textbf{9/23/26 Lecture 12:} Monotone Convergence Theorem}

\begin{definition}[\textbf{Defintion 12.1: }Integrability]
    We say that $f$ is \emph{integrable} if
    \[ \int \abs{f} d\mu < \infty \]
\end{definition}

\textbf{Question:} Given a sequence of measurable functions $\{f_n\}$ with $\lim_{n \to \infty} f_{n}(x) = f(x)$, when can we write
\[ \lim_{n \to \infty} \int f_n d\mu = \int f d\mu \]
This is not so clear. Consider the function $f_n(x) = \chi_{[n, n + 1]}(x)$. We know that $\lim_{n \to \infty} f_n(x) = 0$, and $\int 0 d\mu = 0$, but for any $n$, we have that $\int f_n d\mu = 1$. Thus, 
\[ \lim_{n \to \infty} f_n d\mu \ne \int \lim_{n \to \infty} f_n \]
Here's another problem. Let $f_{n}(x) = n\chi_{(0, 1/n)}$. We see that $\int f_{n} d\mu = 1$, but $f_n \to 0$.\\

Here's a final example. Let $f_{n}(x) = \frac{1}{n}\chi_{[0, n]}$. Then, $\int f_n d\mu = 1$, but $f_n \to 0$. 

\begin{theorem}[\textbf{Theorem 12.2:} Monotone Convergence Theorem]
    Suppose that $0 \leq f_1 \leq f_2 \leq \cdots$ is a monotone sequence of nonnegative measurable functions. Then, 
    \[ \lim_{n \to \infty} \int f_n d\mu = \int \lim_{n \to \infty} f_n d\mu \]
\end{theorem}
\begin{proof}{Proof:}
    Let $f(x) = \lim_{n \to \infty} f_n(x)$. Then, $f_{n}(x) \leq f(x)$ for all $x \in X$ and $n$, so we know that
    \[ \int f_n d\mu \leq \int f d\mu \implies \lim_{n \to \infty} \int f_n d\mu \leq \int f d\mu \]
    Thus, what is left to show is that $\int f d\mu \leq \lim_{n \to \infty} \int f_n d\mu$, which is to show that if $s = \sum_{j = 1}^{N} a_j\chi_{A_j} \leq f$, then $\int s d\mu \leq \lim_{n \to \infty} \int f_n d\mu$.\\

    Take the $A_j$ to be disjoint and let $0 < \epsilon < 1$. Then, for every $x \in A_j$, we see that $f_{n}(x) > (1 - \epsilon)a_j$ for sufficiently large $n$. Now, let
    \[ A_{j, n} = \{x \in A_j \col f_{n}(x) > (1 - \epsilon)a_j\} \]
    Then, notice that $A_{j, n} \uparrow A_{j}$. This implies that
    \[ \mu(A_j) = \lim_{n \to \infty} \mu(A_{j, n}) \]
    and as $f_n \geq \sum_{j = 1}^{N} (1 - \epsilon)a_j\chi_{A_{j, n}}$ for all $n$,  we see that
    \[ \int f_n d\mu \geq \sum_{j = 1}^{N} (1 - \epsilon)a_j\mu(A_{j, n}) \]
    Taking limits as $n \to \infty$, we see that
    \[ \lim_{n \to \infty} f_n d\mu \geq \sum_{j = 1}^{N} (1 - \epsilon)a_j\mu(A_j) = (1 - \epsilon)\int s d\mu \]
    Thus, taking $\epsilon \to 0$, we have that
    \[ \lim_{n \to \infty} \int f_n d\mu \geq \int s d\mu \]
    which was to be shown.

    $\qed$
\end{proof}

\begin{lemma}[\textbf{Corollary 12.3}]
    If $f$ and $g$ are integrable, then $f + g$ is integrable and
    \[ \int (f + g) d\mu = \int f d\mu + \int g d\mu \]
    Additionally, if $f$ and $g$ are non-negative but not necessarily integrable, this equality still holds. 
\end{lemma}
\begin{proof}{Proof:}
    First, suppose $f$ and $g$ are non-negative simple functions $f = \sum_{j = 1}^{n} a_j \chi_{A_j}$ and $g = \sum_{j = 1}^{m} b_j\chi_{B_j}$. Now, $f + g$ is a simple function with
    \[ \int (f + g)d\mu = \sum_{j = 1}^{n}a_j\mu(A_j) + \sum_{j = 1}^{m} b_j\mu(B_j) = \int f d\mu + \int g d\mu \]

    Now, suppose $f$ and $g$ are non-negative. Then, we can find sequences of non-negative simple functions $s_n, t_n$ with $s_n \uparrow f$ and $t_n \uparrow g$. Then, 
    \begin{align*}
        \int f + g d\mu &= \int \lim_{n \to \infty} s_n + t_n d\mu\\
        &= \lim_{n \to \infty} \int s_n + t_n d\mu && (\tt{MCT})\\
        &= \lim_{n \to \infty} \int s_n d\mu + \lim_{n \to \infty} \int t_n d\mu && (\tt{need step 1!})\\
        &= \int \lim_{n \to \infty} s_n d\mu + \int \lim_{n \to \infty} t_n d\mu\\
        &= \int f d\mu + \int g d\mu
    \end{align*}

    Finally, let $f$ and $g$ be integrable but not necessarily nonnegative. To see that $f + g$ is integrable, notice that
    \[ \int \abs{f + g}d\mu \leq \int \abs{f} + \abs{g}d\mu = \int \abs{f} d\mu + \int \abs{g}d\mu \]
    so the integral must be finite. Now, write
    \[ (f + g)^{+} - (f + g)^{-} = (f + g) = (f^{+} - f^{-})  + (g^{+} - g^{-}) \]
    Thus, we see that
    \[ (f + g)^{+} + f^{-} + g^{-} = (f + g)^{-} + f^{+} + g^{+} \]
    These are non-negative functions, so by step 2,
    \[ \int (f + g)^{+}d\mu + \int f^{-}d\mu + \int g^{-}d\mu = \int (f + g)^{-}d\mu + \int f^{+}d\mu + \int g^{+}d\mu \]
    Thus,
    \[ \int (f + g)^{+}d\mu - \int (f + g)^{-}d\mu = \left(\int f^{+}d\mu - \int f^{-}d\mu\right) + \left(\int g^{+}d\mu - \int g^{-}d\mu\right) \]
    which implies that
    \[ \int (f + g)d\mu = \int f d\mu + \int g d\mu \]

    $\qed$
\end{proof}

\begin{lemma}[\textbf{Corollary 12.4}]
   If $\{f_n\}$ is a sequence of non-negative measurable functions, then
   \[ \int \sum_{n = 1}^{\infty} f_n d\mu = \sum_{n = 1}^{\infty} \int f_n d\mu \] 
\end{lemma}
\begin{proof}{Proof:}
    Let $F_N = \sum_{n = 1}^{N} f_n$. We see that $0 \leq F_1 \leq F_2 \leq \cdots$, so
    \[ \lim_{N \to \infty} \int F_{N}d\mu = \int \lim_{N \to \infty} F_N d\mu = \int \sum_{n = 1}^{\infty} f_nd\mu \]
    Then, for each $N$, as we have that
    \[ \int F_{N}d\mu = \sum_{n = 1}^{N}\int f_nd\mu \]
    by Corollary 12.3, we see that
    \[ \sum_{n = 1}^{\infty} \int f_n d\mu = \int \sum_{n = 1}^{\infty} f_n d\mu \]

    $\qed$
\end{proof}

Finally, we will state a lemma which we will prove next time.

\begin{lemma}[\textbf{Lemma 12.5:} Fatou's Lemma]
    Given a sequence $\{f_n\}$ of non-negative measurable functions, we have that
    \[ \int \liminf_{n \to \infty} f_nd\mu \leq \liminf_{n \to \infty} \int f_nd\mu \]
\end{lemma}

This lemma actually lets us understand our previous examples. Let $f_n = \chi_{[n, n + 1]}$. We see that $\liminf_{n \to \infty} f_n = 0$, so clearly
\[ \int \liminf_{n \to \infty} f_nd\mu \leq 1 = \liminf_{n \to \infty} \int f_nd\mu \]
\tocsection{\textbf{9/25/26 Lecture 13: }Dominated Convergence Theorem}

First, we will give a proof of Fatou's Lemma, as promised last section.

\begin{theorem}[\textbf{Lemma 13.1:} Fatou's Lemma]
    For any sequence $\{f_n\}_{n = 1}^{\infty}$ of non-negative measurable functions, we have that
    \[ \int \liminf_{n \to \infty} f_n d\mu \leq \liminf_{n \to \infty} \int f_n d\mu \]
\end{theorem}
\begin{proof}{Proof:}
    Consider the sequence $g_n(x) = \inf_{j \geq n} f_{j}(x)$. Notice it is a sequence of non-negative functions such that $g_{n} \uparrow \liminf_{n \to \infty} f_n$. Now, as $g_n \leq f_{j}$ for all $j \geq n$, we see that
    \[ \int g_n d\mu \leq \inf_{j \neq n} \int f_{j}d\mu \]
    Thus, taking limits, we see that
    \[ \liminf_{n \to \infty} \inf f_{n}d\mu \geq \lim_{n \to \infty} \int g_n d\mu = \int \liminf_{n \to \infty} f_n d\mu \]
    by the MCT.

    $\qed$
\end{proof}

Note that the non-negative requirement is important: If we let $f_n =- 1/n\chi_{[n, 2n]}$ then $\liminf_{n \to \infty} f_n = 0$, but $\int f_nd\mu = -1$ for all $n$.\\

Now, we will prove the Dominated Convergence Theorem, which tells us that we can interchange limits when our sequence of functions is dominated by an integrable function.

\begin{theorem}[\textbf{Theorem 13.2:} Dominated Convergence Theorem]
    Let $\{f_n\}_{n = 1}^{\infty}$ be a sequence of measurable functions, and let $f_n(x) \to f$ for all $x \in X$. Suppose that there is some non-negative integrable function $g$ for which $\abs{f_n(x)} \leq g(x)$ for all $x$. Then, 
    \[ \int f d\mu = \lim{n \to \infty} \int f_nd\mu \]
\end{theorem}
\begin{proof}{Proof:}
    The idea is that $f_n + g \geq 0$ and $g - f_n \geq 0$ are both positive functions. By Fatou's lemma, we see that
    \[ \int (f + g)d\mu \leq \liminf_{n \to \infty} \int (f_n + g) d\mu = \liminf_{n \to \infty} \int f_n d\mu + \int g d\mu \]
    Since $g$ is integrable, we can subtract it to get
    \[ \int f d\mu \leq \liminf_{n \to \infty} \int f_n d\mu \]
    We also have by Fatou's lemma that
    \[ \int g - f d\mu \leq \liminf_{n \to \infty} \int g - f_n d\mu = \int g + \liminf_{n \to \infty} \int -f_n d\mu \]
    Subtracting $\int g$ again, we have
    \[ -\int f d\mu \leq \liminf_{n \to \infty} -\int f_nd\mu \]
    Now, recall that $\liminf_{n \to \infty} -\int f_n d\mu = -\limsup_{n \to \infty} \int f_n d\mu$, so dividing by $-1$, we get
    \[ \int f d\mu \geq \limsup_{n \to \infty} \int f_n d\mu \]
    THus, we see that
    \[ \int f d\mu \leq \liminf_{n \to \infty} \int f_n d\mu \leq \limsup_{n \to \infty} \int f_n d\mu \leq \int f d \mu \]
    so equality holds and thus
    \[ \int f d\mu = \lim_{n \to \infty} \int f_n d\mu \]

    $\qed$
\end{proof}

Now, here's a corollary of the MCT and DCT.

\begin{lemma}[\textbf{Corollary 13.3}]
    If $f_n$ is a sequence of functions satisfying the hypotheses of the monotone or dominated convergence theorems, but suppose that $f_{n}(x) \to f(x) \tt{ a.e.}$. Then, we still have that
    \[ \int \lim_{n \to \infty} d\mu = \int \lim_{n \to \infty} f_n d\mu \]
\end{lemma}
\begin{proof}{Proof:}
    Let $A$ be the set of points for which $f_n(x)$ converges to $f(x)$. Notice that for any measurable function, we must have that
    \[ \int g d\mu = \int g\chi_{A} d\mu + \int g\chi_{A^{c}}d\mu = \int g \chi_{A}d\mu \]
    because the integral of a function only supported on a measure 0 set is 0.\\

    Then, applying this fact to $f_{n}\chi_{A}$, we see that
    \[ \lim_{n \to \infty} \int f_nd\mu = \lim_{n \to \infty} \int f_n\chi_{A}d\mu = \int f \chi_{A}d\mu = \int f d\mu \]

    $\qed$
\end{proof}

Now, some notation that we will use in the next proposition, as well as for the rest of this class. For $A \subseteq X$, define
\[ \int_{A} f d\mu := \int f \chi_{A} d\mu \]
Also, if we have a Lebesgue measurable function $f : \R \to \R$, write
\[ \int f(x) dx := \int f d\lambda \]
where $\lambda$ is the Lebesgue measure. Finally, we define
\[ \int_{a}^{b} f d\mu = \left\{\begin{matrix} \int_{[a, b]}fd\mu & a < b \\ -\int_{[b, a]} f d\mu & b < a \end{matrix}\right. \]

\begin{lemma}[\textbf{Proposition 13.4}]
    The following are true about integrals.
    \begin{enumerate}
        \item If $f \geq 0$ is measurable and $\int f d\mu = 0$, the $f = 0$ almost everywhere.
        \item If $f$ is integrable and $\int_{A} d\mu = 0$ for all measurable $A$, then $f = 0$ almost everywhere.
        \item If $f : \R \to \R$ is Lebesgue integrable and, for some $a$, $\int_{a}^{x} f(y)dy = 0$ for all $x$, then $f = 0$ almost everywhere. 
    \end{enumerate}
\end{lemma}
\begin{proof}{Proof:}
    
    \textbf{(1)} Notice that
    \[ f \geq \frac{1}{n}\chi_{\{x \col f(x) > 1/n\}} \]
    Thus, we see that
    \[ \int f d\mu \geq \frac{1}{n} \mu(\{x \col f(x) > 1/n\}) \]
    which implies that $\mu(\{x \col f > 1/n\}) = 0$. Thus, as $\bigcup_{n = 1}^{N} \{x \col f(x) > 1/n\} \uparrow \{x \col f(x) > 0\}$, we see that
    \[ \mu({x \col f(x) > 0}) = 0 \]
    Thus, as $f$ is non-negative, this implies that $f$ is $0$ almost everywhere.\\

    \textbf{(2)} Apply (1) to the sets $f\chi_{\{x \col f(x) > 0\}}$ and $-f\chi_{\{x \col f(x) > 0\}}$, which are both non-negative.\\

    \textbf{(3)} By linearity, we see that
    \[ \int_{b}^{c} f(y) dy = \int_{a}^{c} f(y)dy - \int_{a}^{b}f(y)dy = 0 \]
    for all intervals $[b, c]$. Then, this implies that for any open set $U = \bigcup_{n}(a_n, b_n)$ (take intervals disjoint), we have that
    \[ \int_{U} f d\mu = \int \sum_{n} f(y)\chi_{(a_n, b_n)}(y)dy = \sum_{n} \int_{a_n}^{b_n} f(y)dy = 0 \]
    by the dominated convergence theorem applied to the sums (the bounding function is $f$). 
\end{proof}
\tocsection{\textbf{9/30/26 Lecture 14:} Riemann and Lebesgue Integrals}
\vspace{0.2cm}

\begin{proof}{Proof of 13.4 cont.}
    Recall that, if $a \in [b, c]$ we had
    \[ \int_{b}^{c}f(y)dy = \int_{b}^{a} f(y)dy + \int_{a}^{c} f(y) dy = 0 \]
    for all $b, c$ by our assumption. If $a \notin [b, c]$, then we get a similar identity that shows $\int_{b}^{c} f(y)dy = 0$.\\
 
    Now, if $U$ is open, then we can written as $U = \bigcup_{i= 1}^{\infty} (b_n, c_n)$ with the intervals taken to be disjoint, which tells us that
    \begin{align*}
        \int_{U} f dy &= \int \sum_{n} f(y)\chi_{b_n, c_n}dy\\
        &= \sum_{n} \int_{b_n}^{c_n} f(y)dy\\
        &= 0
    \end{align*}

    Then, for any measurable $A$, then we can write $A = \left(\bigcap_{n = 1}^{\infty} U_n \right) \cup E$ where $U_n$ is open and $E$ is a null set. Thus, we have
    \[ \int_{A} f d\mu = \int_{\bigcap_{n = 1}^{\infty} U_n} f(y) dy + \int_{E} f dy = 0\]
    as $E$ is a null set, and as $\bigcap_{n = 1}^{\infty} U_n$ is in an open set. Thus, we see that $f$ is zero almost everywhere by $2)$. 

    $\qed$
\end{proof}

\begin{lemma}[\textbf{Proposition 14.1}]
    For any Lebesgue integrable function $f$, for all $\epsilon > 0$, there exists a continuous $g$ with a compact support so that
    \[ \int \abs{f - g}dy < \epsilon \]
\end{lemma}
\begin{proof}{Proof:}
    Let's start with simple functions. Consider $\chi_{A}$, where $A$ is a bounded set. By inner and outer regularity, we can find compact $K$ and open $U$ with $K \subseteq A \subseteq U$ so that
    \[ \lambda(A \setminus K) < \epsilon/2 \quand \lambda(U \setminus A) < \epsilon/2 \]
    Then, let
    \[ g(x) = \frac{d(x, U^{c})}{d(x, K) + d(x, U^{c})} \]
    This function is continuous and is $1$ on $K$ and $0$ on $U^{c}$.\\

    Now, we see that $\abs{\chi_{A} - g} \leq 1$, and is supported on $(A \setminus K) \cup (U \setminus A)$, which is a set with measure zero, so
    \[ \int \abs{\chi_{A} - g}dx \leq \epsilon \]
    Additionally, we see that the support of $g$ is the closure of $U$, which is what we want.\\

    Now, let's do simple functions. Let $f$ be a simple function with compact support
    \[ f = \sum_{j = 1}^{n} a_j\chi_{A_j} \]
    with $A_j$ bounded. We can also take $a_j \ne 0$. For each $j$, by the previous step, we can find $g_j$ so that
    \[ \int \abs{\chi_{A_j} - g_j} < \frac{\epsilon}{na_j} \]
    Now, consider
    \[ g = \sum_{j = 1}^{n} a_jg_j \]
    then, we see that
    \begin{align*}
        \int \abs{\sum_{j} a_j\chi_{A_j} - g}dx &\leq \sum_{j} \int \abs{a_j\chi_{j} - a_jg_j} dx\\
        &\leq \sum_{j} a_j\left(\frac{\epsilon}{na_j}\right)dx\\
        &\leq \epsilon
    \end{align*}

    Now, let $f$ be nonnegative with compact support. By the definition of the Lebesgue integral, we can find a simple $s$ such that $s \leq f$ and
    \[ \int f \leq \int s + \epsilon/2 \]
    Then, applying the above step to $s$, we can find $g$ so that $\int\abs{s - g} dx < \epsilon/2$, which gives us that
    \begin{align*}
        \int \abs{f - g}dx &\leq \int \abs{f - s}dx + \int \abs{s - g}dx\\
        &= \int (f - s)dx + \int \abs{s - g}dx\\
        &< \epsilon 
    \end{align*}

    Now, let $f$ be nonnegative but with a support that is not necessarily bounded. Now, we know that
    \[ f_n = f \chi_{[-n, n]} \uparrow f \]
    by the MCT, we have that $\int f_n dx \uparrow \int f$. Thus, we can pick $n$ large enough so that
    \[ \int f dx - \int f_n dx < \epsilon / 2 \]
    then we can apply the previous step to the $f_n$ (which now is nonnegative with compact support).\\

    Finally, consider a general Lebesgue integrable $f$. Then, we can apply step 4 to $f^{+}$ and $f^{-}$ and use the triangle inequality on $f = f^{+} - f^{-}$. 

    $\qed$
\end{proof}

Now that we've defined a new notion of integration, let's go back and see if it coincides with the Riemann integral. Consider a bounded function $f : [a, b] \to \R$. If $P$ is a partition of $[a, b]$,
\[ \{x_0, \ldots, x_n\} \tt{ with } a = x_0 < x_1 < x_2 < \cdots < x_n = b \]
we have
\[ U(P, f) = \sum_{j = 1}^{n} \sup_{x \in [x_{j - 1}, x_j]}f(x) \cdot \lambda[x_{j - 1}, x_{j}] \]
\[ L(P, f) = \sum_{j = 1}^{n} \inf_{x \in [x_{j - 1}, x_j]}f(x) \cdot \lambda[x_{j - 1}, x_{j}] \]
When $\inf_{P} U(P, f) = \sup_{P} L(P, f)$ we say that $f$ is Riemann integrable with the integral being either of these values. Consider the simple function
\[ s_{P} = \sum_{i = 1}^{n} (\sup_{x \in [x_{j - 1}, x_{j}]}f(x) )\chi_{[x_{j - 1}, x_{j}]}(x) \]
\fig{0.8}{Lecture14Fig1}{}

Notice that $\int s_{P}(x)dx = U(P, f)$ (where we are taking the Lebesgue integral).\\

Thus, Riemann integration is like Lebesgue integration, except for the fact that we can only use simple functions over intervals (so strictly weaker). Here are some interesting things to note that will lead us to the following theorem.

\begin{example}{Example}
    Consider $\chi_{\Q \cap [0, 1]}$. This is not Riemann integrable, but it isn't continuous anywhere (importantly, not continuous on a full measure set).\\

    Now, consider the Thomae's function, which is $0$ on the irrationals, and equal to the denominator on any rational (in lowest terms). This is Riemann integrable, and is continuous on $[0, 1] \setminus \Q$, a full measure set. 
\end{example}

\begin{theorem}[\textbf{Theorem 14.2: }Equivalence of Riemann and Lebesgue Integrals]
    Let $f : [a, b] \to \R$ be bounded. 
    \begin{enumerate}
        \item If $f$ is Riemann integrable then it is continuous Lebesgue almost everywhere and its Riemann integrable is equal to it's Lebesgue integrable.
        \item If $f$ is measurable and bounded almost everywhere, then it is Riemann integrable.
    \end{enumerate}
\end{theorem}
\begin{proof}{Proof:}
    Let $P$ be any partition and define
    \begin{align*}
        T_{P} &= \sum_{j = 1}^{n} \left(\sup_{x \in [x_{j - 1}, x_{j}]} f(x)\right) \chi_{[x_{j - 1}, x_{j}]}(x) \\
        S_{P} &= \sum_{j = 1}^{n} \left(\inf_{x \in [x_{j - 1}, x_{j}]} f(x)\right) \chi_{[x_{j - 1}, x_{j}]}(x)
    \end{align*}
    Let $R(f)$ denote the Riemann integral of $f$. By it's definition, we can find partitions $Q_j$ and $Q_j'$ so that
    \[ U(Q_j, f) \downarrow R(f) \quand L(Q', f) \uparrow R(f) \]
    Now, recall that if $P \subseteq P'$, then
    \[ U(P', f) \leq U(P, f) \quand L(P', f) \geq L(P, f) \]
    So if we define
    \[ P_{j} := \bigcup_{i \leq j} (Q_i \cup Q_{i}') \]
    then we still have that
    \[ U(P_j, f) \downarrow R(f) \quand L(P_j, f) \downarrow R(f) \]
\end{proof}
\tocsection{\textbf{10/2/26 Lecture 15:} Modes of Convergence}

\vspace{0.2cm}

\begin{proof}{\textbf{Proof of Theorem 14.2 (cont.)}}
    Now, recall that we have a sequence $P_1 \subseteq P_2 \subseteq \cdots$ of partitions so that $T_{P_j}$ is decreasing and $S_{P_j}$ is increasing. Also, we have that $T_{P_{j}} \geq f \geq S_{P_{j}}$, and by the monotone convergence theorem, there must be some $T$ and $S$ so that $T_{P_{j}} \to T$ and $S_{P_{j}} \to S$. Also, recall these functions have the property that
    \[ \int T_{P_{j}} dx = U(P_j, f) \quand \int S_{P_{j}} dx = L(P_{j}, f) \]

    Now, $\abs{T_{P_{j}}} \leq \sup_{x \in [0, 1]}\abs{f(x)} < \infty$ ($f$ bounded) and the same inequality holds for $\abs{S_{P_{j}}}$. Because of this, we can apply the DCT to the following:
    \begin{align*}
        \int T - S dx &= \lim_{j \to \infty} \int T_{j} - S_{j} dx \\
        &= \lim_{j \to \infty} U(P_{j}, f) - L(P_{j}, f)\\
        &= 0
    \end{align*}
    as $f$ is Riemann integrable. Then, because $\int (T - S) dx = 0$, and because $T - S \geq 0$, we see that $T = S$ almost everywhere. Thus, we must have that $T = f = S$ almost everywhere.\\

    Now, we want to show that $f$ is continuous a.e., so what we want to show that if $T(x) = S(x)$ and $x \notin P_{j}$ for any $j$, then $f$ is continuous at $x$. Let $\epsilon > 0$. Then, as $T$ and $S$ agree at $x$, find a $j$ so that
    \[ T_{P_{j}}(x) - S_{P_{j}}(x) < \epsilon \]
    and choose $\delta$ so that $x \in (x - \delta, x + \delta) \subseteq [x_{j - 1}, x_{j}]$. On this interval, note that $T_{P_{j}}$ and $S_{P_{j}}$ are constant. Thus, $T_{P_{j}}(y) - S_{P_{j}}(y) < \epsilon$ for all $y \in (x - \delta, x + \delta)$. Notice this implies that $T_{P_{k}}(y) - S_{P_{k}}(y) < \epsilon$ for all $k \geq j$ on $(x - \delta, x + \delta)$, which implies that $f$ is continuous at $x$ (bound the variance of $f$ by the constant values of $T_{P_{j}}$ and $S_{P_{j}}$).\\

    Finally, we see that
    \[ \int f = \lim_{j \to \infty} \int T_{P_{j}} = \lim_{j \to \infty} U(P_{j}, f) = R(f) \]
    so the Lebesgue and Riemann integrals of $f$ agree, which was to be shown.\\

    \textbf{(2)} Let $P_1 \subseteq P_2 \subseteq \cdots$ be some increasing sequence of partitions with $\tt{mesh}(P_{j}) \to 0$. For instance, take
    \[ P_{j} = \{a + n(b - a)2^{-j} \col n = 0, \ldots, 2^{j}\} \]
    Suppose that $x$ is a point of continuity of $f$ that is not in any of the $P_{j}$. For each $j$, pick $[y_{j}, z_{j}]$ (with endpoints in $P_{j}$) containing $x$. Now $y_{j} \to x$ and $z_{j} \to x$. Then, we see that
    \[ \sup_{y \in [y_{j}, z_{j}]} f(y) \to f(x) \]
    because $f$ is continuous as $x$. Thus, we see that $T_{P_{j}}(x) \to f(x)$ and by the same argument, $S_{P_{j}}(x) \to f(x)$ for almost every $x \in [a, b]$. Then, 
    \[ U(P_{j}, f) = \int T_{P_{j}} dx \to \int f dx\]
    and
    \[ L(P_{j}, f) = \int S_{P_{j}} dx \to \int f dx \]
    Thus, we see that $f$ is Riemann integrable.

\end{proof}

\subsection*{Modes of Convergence}

Let $(X, \cal{A}, \mu)$ be a measure space. 

\begin{definition}[\textbf{Definition 15.1: }Pointwise Almost Everywhere Convergence]
    Given $f_n : X \to \R$ measurable, we say that $f_n \to f$ \emph{pointwise almost everywhere} if 
    \[ \mu\{x \col f_{n}(x) \not \to f(x)\} = 0 \]
    We usually denote this $f_n \to f$ a.e.
\end{definition}

\begin{definition}[\textbf{Definition 15.2: }Convergence in Measure]
    Given $f_n : X \to \R$ measureable, we say that $f_n \to f$ \emph{in measure} if, for any $\epsilon > 0$, 
    \[ \mu\{x \col \abs{f_{n}(x) - f(x)} > \epsilon\} \to 0 \]
    as $n \to \infty$. 
\end{definition}

It turns out convergence in measure does not imply convergence pointwise almost everywhere. Consider the following function:\\

\fig{0.6}{Lecture15Fig1}{}


Note that if we go far enough, $\mu(\abs{f_n} > 0)$ is $1/n$, so $f_n$ converges in measure to $0$. However, notice that $f_n(x) = 1$ infinitely many times, so $f_n$ never converges pointwise to $0$. \\


It also turns out that, when $\mu(X) = \infty$, convergence pointwise does not imply convergence in measure. Consider the functions
\[ f_n = \chi_{[n, n + 1]} \]
We see that $f_n \to 0$ almost everywhere (pointwise), but $\mu(\abs{f_n}(x) > 0) = 1$. Convergence pointwise seems stronger, but we have problems when mass escapes to infinity. Thus, if we assume this isn't possible, it turns out we can have one implication.

\begin{lemma}[\textbf{Proposition 15.3}]
    Let $f_n : X \to \R$ be measurable.
    \begin{enumerate}
        \item If $\mu$ is a finite measure, then $f_n \to f$ almost everywhere implies that $f_n \to f$ in measure.
        \item If $\mu$ is any measure, and if $f_n \to f$ in measure, then there is a subsequence $f_{n_{j}} \to f$ almost everywhere. 
    \end{enumerate}
\end{lemma}
\begin{proof}{Proof}
    \textbf{(1)} Let
    \[ A_n = \{x \col \abs{f_{n}(x) - f(x)} > \epsilon\} \]
    Then, as $f_n(x) \to f(x)$ almost everywhere, this implies that $\chi_{A_n}(x) \to 0$ almost everywhere. Then, we see that
    \[ \lim_{n \to \infty} \mu(A_n) = \int \lim_{n \to \infty} \chi_{A_n}(x) dx = \int 0 = 0 \]
    where the dominated convergence theorem was needed for the above equality, using the fact that $\abs{\chi_{A_n}} \leq 1$ and $1$ is integrable when $\mu(X) < \infty$. 
\end{proof}
\tocsection{\textbf{10/5/26 Lecture 16:} $L^p$ Convergence and Product Measures}
\vspace{0.2cm}

\begin{proof}{Proof of (2):}
    Let $n_1 = 1$. Now, pick $n_{j} > n_{j - 1}$ so that
    \[ \mu(A_{j}) < 2^{-j} \]
    where $A_{j} = \{x \col \abs{f_{n_{j}}(x) - f(x)} > 1/j\}$. Let
    \[ A = \bigcap_{k = 1}^{\infty} \bigcup_{j \geq k} A_{j} \]
    so we see 
    \begin{align*}
        \mu(A) &= \lim_{k \to \infty} \mu\left(\bigcup_{j \geq k} A_{j}\right)\\
        &= \lim_{k \to \infty} \sum_{j = k}^{\infty} 2^{-j} = 0
    \end{align*}
    We claim that if $x \notin A$, then $f_{n_{j}}(x) \to f(x)$. Note that points in $A$ are points that are in infinitely many $A_{j}$. Thus, for sufficiently large $j$, we see that $\abs{f_{n_{j}}(x) - f(x)} < \frac{1}{j}$, so $f_{n_{j}}(x) \to f(x)$. Thus, $f_{n_{j}}$ converges almost everywhere to $f$. 

    $\qed$
\end{proof}

Now, we give one more definition of convergence. 

\begin{definition}[\textbf{Definition 16.1:} Convergesnce in $L^{p}$]
    Let $1 \leq p < \infty$. A sequence of measurable functions $f_n$ \emph{converges to $f$ in $L^{p}$} if
    \[ \int \abs{f_n - f}^{p}d\mu \to 0 \]
    as $n \to \infty$. 
\end{definition}

Why didn't we just choose $p = 0$? Here is why:

\begin{example}{Example}
    Let $f_n = n^{1/2}\chi_{[0, 1/n]}$. We see that
    \[ \int \abs{f_n - 0}dx = \frac{1}{n^{1/2}} \to 0 \]
    so $f_n \to 0$ in $L^{1}$. However, notice that
    \[ \int \abs{f_n - 0}^2 dx = \int n\chi_{[0, 1/n]}^2 dx = \int 1 dx = 1 \]
    Heuristically, in $L^1$ area under the curve is width times height, but in $L^2$ area is width times height squared. 
\end{example}

\begin{theorem}[\textbf{Lemma 16.2:} Chebyshev's Inequality]
    For any integrable $f$ and some $1 \leq p < \infty$, then for any $\epsilon > 0$, 
    \[ \mu(\{x \col \abs{f(x)} \geq a\}) \leq \frac{\int \abs{f}^{p} d\mu}{a^{p}} \]
\end{theorem}
\begin{proof}{Proof:}
    Notice that
    \[ a^{p}\chi_{\{x \col \abs{f(x)} \geq a\}} \leq \abs{f(x)}^{p} \]
    on $A$. Integrating, we see that
    \[ a^{p}\mu(\{x \col \abs{f(x)} \geq a\}) \leq \int_{A} \abs{f}^{p} \leq \int \abs{f}^{p} \]
    which gives us our desired inequality.

    $\qed$
\end{proof}

\begin{lemma}[\textbf{Proposition 16.3}]
    Convergence in $L^{p}$ implies convergence in measure, which implies a subsequence converges almost everywhere. 
\end{lemma}
\begin{proof}{Proof:}
    Fix $\epsilon > 0$. By Chebyshev's inequality, we have that
    \[ \mu(\{x \col \abs{f_{n}(x) - f(x)} \geq \epsilon\}) \leq \epsilon^{p} \int \abs{f_n(x) - f(x)}^{p} d\mu \]
    which goes to zero as $n \to \infty$. 

    $\qed$
\end{proof}

Now, recall that $f_{n}(x) \to f$ uniformly on $A$ if
\[ \lim_{n \to \infty} \sup_{x \in A} \abs{f_{n}(x) - f(x)} = 0 \]
This brings us to a very useful theorem.

\begin{theorem}[\textbf{Theorem 16.4:} Egorov's Theorem]
    Let $\mu$ be a finite measure. Then, if $f_n \to f$ almost everywhere, for any $\epsilon > 0$, there exists a set $A$ with $\mu(A) < \epsilon$ such that $f_n \to f$ uniformly on $A^{c}$.
\end{theorem}
\begin{proof}{Proof:}
    Define the sets
    \[ A_{n, k} = \bigcup_{m \geq n} \{x \col \abs{f_m(x) - f(x)} > 1/k\} \]
    Fixing $k$, we see that $A_{n, k}$ is decreasing in $n$. Now, notice that $\bigcap_{n} A_{n, k} \subseteq \{x \col f_m(x) \not \to f(x)\}$, which was measure zero. We also see that
    \[ \lim_{n \to \infty} \mu(A_{n, k}) = \mu\left(\bigcap_{n} A_{n, k}\right) \]
    which goes to zero, so choose $n_{k}$ so that $\mu(A_{n_{k}, k}) < \epsilon 2^{-k}$. Now, let
    \[ A = \bigcup_{k = 1}^{\infty} A_{n_{k}, k} \]
    by subadditivity, $\mu(A) < \epsilon$. Now, if $x \notin A$, then for all $k$, 
    \[ \abs{f_{n}(x) - f(x)} < 1/k \]
    for $n \geq n_{k}$. As this holds for all $x \in A^{c}$, we see that
    \[ \sup_{x \in A} \abs{f_{n}(x) - f(x)} \leq 1/k \]
    for all $n \geq n_{k}$, so uniform convergence follows taking $k \to 0$.

    $\qed$
\end{proof}

\textbf{Remark:} This can be used to give a pretty quick proof of Lusin's theorem. 

\subsection*{Product Measures}
\sectionspace

Recall the Lebesgue measure on $\R$. The way we constructed it was starting with a premeasure on algebra of intervals, moving this to the Lebesgue outer measure, and then going to the measure. Now what we want is the Lebesgue measure on $\R^n$.\\

Lets first consider $\R^2$. Define
\[ \lambda^{2}([a, b] \times [c, d]) = (b - a) \cdot (d - c) = \lambda([a, b]) \cdot \lambda([c, d]) \]
Thus, we are thinking of $\R^2$ as $\R \times \R$, and saying that the measures of rectangles is length times height!\\

Fix two measure spaces $(X, \cal{A}, \mu)$ and $(Y, \cal{B}, \nu)$.
\begin{definition}[\textbf{Definition 16.5:} Measurable Rectangles and Product $\sigma$-Algebras]
    A \emph{measurable rectangle} in $X \times Y$ is a set of the form $A \times B$ where $A \in \cal{A}$ and $B \in \cal{B}$. The \emph{product $\sigma$-algebra} is the $\sigma$-algebra generated by measurable rectangles. 
\end{definition}

Now, we want to define a premeasure on an algebra that we can then extend. Let $\cal{C}_{0}$ be the collection of finite unions of measurable rectangles. We claim that this is an algebra. It is clear this is closed under finite unions, so we need to check complements. We see that
\[ (A \times B)^{c} = (X \times B^{c}) \cup (A^{c} \times B) \]
Now, for $A \times B$, let $\ell(A \times B) = \mu(A) \cdot \nu(B)$. 

\begin{lemma}[\textbf{Proposition 16.6}]
    Let $A \times B = \bigcup_{n = 1}^{\infty} A_n \times B_n$ where $A_n \in \cal{A}$, $B_n \in \cal{B}$, and they are all disjoint. Then, 
    \[ \ell(A \times B) = \sum_{n} \ell(A_{n} \times B_{n}) \]
    which is to show that
    \[ \mu(A)\nu(B) = \sum_{n} \mu(A_n) \nu(B_n) \]
\end{lemma}
\begin{proof}{Proof:}
    First, write
    \[ \chi_{A}(x)\chi_{B}(y) = \chi_{A \times B}(x, y) = \sum_{n} \chi_{A_n \times B_{n}}(x, y) = \sum_{n} \chi_{A_n}(x)\chi_{B_n}(y) \]
    fixing $y$, we get an equality of functions of $x$, so integrating
    \begin{align*}
        \mu(A)\chi_{B}(y) &= \sum_{n} \int \chi_{A_n}(x)\chi_{B_n}(y)d\mu(x)\\
        &= \sum_{n} \mu(A_n)\chi_{B_n}(x)
    \end{align*}
    Then, as this gives us a function of $y$, we see that integrating this again, 
    \[ \mu(A)\mu(B) = \sum_{n} \mu(A_n) \int \chi_{B_n}(y)d\nu = \sum_{n} \mu(A_n)\nu(B_n) \]

    $\qed$
\end{proof}

By the Caratheodory extension theorem, we get a measure on $(X \times Y, \cal{A} \times \cal{B})$ with the property that
\[ (\mu \times \nu)(A \times B) = \mu(A) \nu(B),\quad A \in \cal{A}, B \in \cal{B} \]
\tocsection{\textbf{10/7/26 Lecture 17:} Product Measures (cont.)}

Recall from last time we defined a measure on $(X \times, Y, \cal{A} \times \cal{B})$ with the property that
\[ (\mu \times \nu)(A \times B) = \mu(A) \times \nu(B) \]
and by Caratheodory, this is unique. Now, we need to define how to integrate with respect to this measure.\\

Let $E \subseteq X \times Y$, and, for fixed $x, y$, define
\[ E_{x} = \{y \in Y \col (x, y) \in E\} \]
and
\[ E^{y} = \{x \in X \col (x, y) \in E\} \]

Similarly, given a function $f : X \times Y \to \R$, for fixed $x, y$, define
\[ f_{x} : Y \to \R \tt{ by } f_{x}(y) = f(x, y) \]
and
\[ f^{y} : X \to \R \tt{ by } f^{y}(x) = f(x, y) \]

\begin{lemma}[\textbf{Proposition 17.1}]
    \vspace{-0.6cm}
    \begin{enumerate}
        \item If $E \in A \times B$, then $E_{x} \in \cal{B}$ for any $x \in X$ and $E^{y} \in \cal{A}$ for any $y \in Y$. 
        \item If $f$ is $\mu \times \nu$ measurable, then $f_{x}$ and $f^{y}$ are measurable with respect to $\mu$ and $\nu$.
    \end{enumerate}
\end{lemma}
\begin{proof}{Proof:}
    
    (1) Let $\cal{C}$ be the collection of subsets $E \subseteq X \times Y$ such that $E_{x} \in \cal{B}$ for all $x$ and $E^{y} \in \cal{A}$ for all $y$. We want to show that $\cal{C}$ contains measurable rectangles and is a $\sigma$-algebra.\\

    First, we see that if $E = A \times B$, then 
    \[ E_{x} = \left\{\begin{matrix} B & x \in A \\ \emptyset & x \notin A\end{matrix}\right. \]
    and
    \[ E^{y} = \left\{\begin{matrix} A & y \in B \\ \emptyset & y \notin B \end{matrix}\right. \]
    so all rectangles are in $\cal{C}$. Clearly, $\cal{C}$ is a $\sigma$-algebra as it is clear that it is closed under unions and complements (pass conditions to the restrictions). Thus, we see that $\cal{A} \times \cal{B} \in \cal{C}$, which was to be shown.\\

    (2) We see that for $U \subseteq \R$, so as
    \[ f_{x}^{-1}(U) = (f^{-1}(U))_{x}\]
    so if $f^{-1}(U) \in \cal{A} \times \cal{B}$, we see that $f_{x}^{-1}(U) \in \cal{B}$ by (1).

    $\qed$
\end{proof}

\begin{definition}[\textbf{Definition 17.2:} Monotone Classes]
    We say a set $\cal{M} \subseteq \cal{P}(X)$ is a \emph{monotone class} if
    \begin{enumerate}
        \item $A_1 \subseteq A_2 \subseteq \cdots,\quad A_n \in \cal{M} \implies \bigcup_{n} A_n \in \cal{M}$
        \item $A_1 \supseteq A_2 \supseteq \cdots,\quad A_n \in \cal{M} \implies \bigcap_{n} A_n \in \cal{M}$
    \end{enumerate} 
\end{definition}

Given a collection $\cal{C} \in \cal{P}(X)$, the monotone class generated by $\cal{C}$ is the smallest monotone class containing $\cal{C}$; that is, it is the intersection of all monotone classes containing $\cal{C}$. 

\begin{theorem}[\textbf{Theorem 17.3: }Monotone Class Theorem]
    Let $\cal{A} \subseteq \cal{P}(X)$ be an algebra. Let $\cal{M}$ be the monotone class generated by $\cal{A}$. Then, $M = \sigma(\cal{A})$. 
\end{theorem}
\begin{proof}{Proof:}
    Firstly, it is clear that $\cal{M} \subseteq \sigma(\cal{A})$ as $\sigma$-algebras are monotone classes. Now, we want to show that $\sigma(\cal{A}) \subseteq \cal{M}$, which is to show that $\cal{M}$ is a $\sigma$-algebra.\\
    
    Firstly, we need to show that if $A \in \cal{M}$, then $A^{c} \in \cal{M}$. Let $\cal{N}_1 = \{A \in \cal{M} \col A^{c} \in \cal{M}\}$. We see that $\cal{N}_1$ contains $\cal{A}$. Now, if $A_1 \subseteq A_2 \subseteq \cdots$ with $A_n \in \cal{N}_1$, since $\bigcup_{n} A_n \in \cal{M}$, and
    \[ \left(\bigcup_{n} A_n\right)^{c} = \left(\bigcap_{n} A_n^{c}\right) \]
    then as $A_1^{c} \supseteq A_2^{c} \supseteq \cdots$, we see that $\bigcap_{n} A_n^{c} = \left(\bigcup_{n} A_n\right)^{c} \in \cal{M}$. Thus, $\bigcup_{n} A_n \in \cal{N}_1$. By the same argument, we see that if $A_1 \supseteq A_2 \supseteq \cdots$, then $\bigcap_{n} A_n \in \cal{N}_1$, so $\cal{N}_1$ is a monotone class. Thus, $\cal{M} \subseteq \cal{N}_1$, so $\cal{M}$ is closed under complements.\\

    Now, we want to show that it is closed under intersection. Start by letting $\cal{N}_2 = \{B \in \cal{M} \col A \cap B \in \cal{M} \text{ for all } A \in \cal{A}\}$. We see that $\cal{N}_2$ contains $\cal{A}$. Now, suppose that $B_1 \subseteq B_2 \subseteq \cdots$ with $B_n \in \cal{N}_2$. Then, for any $A \in \cal{A}$, we see that as $\bigcup_{n} B_n \in \cal{M}$,
    \[ \bigcup_{n} B_n \cap A = \bigcup_{n} (B_n \cap A) \]
    and as $B_1 \cap A \subseteq B_2 \cap A \cdots$ and $B_n \cap A \in \cal{M}$, we see that $\bigcup_{n} B_n \cap A \in \cal{M}$, so $\bigcup_{n} B_n \in \cal{N}_2$. Since a similar proof for intersections follows, we see that $\cal{M} \subseteq \cal{N}_2$.\\

    For the next step, let $\cal{N}_3 = \{B \in \cal{M} \col A \cap B \in \cal{M} \tt{ for all } A \in \cal{M}\}$. Now, $\cal{N}_3$ contains $\cal{A}$ because of the last step. By the same argument as above, we see that $\cal{N}_3$ is a monotone class and so $\cal{M} \subseteq \cal{N}_3$.\\
    
    The last thing to check is countable intersections. Let $\{A_n\} \subseteq \cal{M}$. Now, since we know that $\bigcap_{i = 1}^{n} A_i \in \cal{M}$ for each $n$, and as
    \[ \bigcap_{i = 1}^{1} A_i \supseteq \bigcap_{i = 1}^{2} A_i \supseteq \cdots \]
    we see that
    \[ \bigcap_{n = 1}^{\infty} \bigcap_{i = 1}^{n} A_i = \bigcap_{i = 1}^{\infty} A_n \in \cal{M} \]
    Thus, $\cal{M}$ is a $\sigma$-algebra, which was to be proved.

    $\qed$
\end{proof}

\begin{lemma}[\textbf{Proposition 17.4}]
    Suppose that $\mu$ and $\nu$ are $\sigma$-finite. If $E \in \cal{A} \times \cal{B}$, then $x \mapsto \nu(E_{x})$ and $y \mapsto \mu(E^{y})$ are measurable on $X$ and $Y$ respectviely, and
    \[ (\mu \times \nu)(E) = \int_{X} \nu(E_{x})d\mu = \int_{Y} \mu(E^{y})dy \]
\end{lemma}
\begin{proof}{Proof:}
    First, assume that $\mu$ and $\nu$ are finite. Then, let $\cal{M}$ be the collection of the sets so that the conclusions of the theorem holds. We want to show that $\cal{M}$ contains finite unions of measurable rectangles and is a monotone class.\\

    First, let $E = A \times B$. Then, as
    \[ E_{x} = \left\{\begin{matrix} \emptyset & x \notin A \\ B & x \in A \end{matrix}\right. \]
    so $\nu(E_{x}) = \chi_{A}(x) \cdot \nu(B)$, which is measurable, and similarly for $E^{y}$. Similarly, we also see that
    \[ \int_{X} \nu(E_{x})dx = \mu(A)\nu(B) = (\mu \times \nu)(E) \]
    and similarly for the other integral. Thus, measurable rectangles are in $\cal{M}$.\\

    Now, let $E_1 \subseteq E_2 \subseteq \cdots$ are in $\cal{M}$ and let $E = \bigcup_{n} E_n$. Let $f_{n}(y) = \mu(E_{n}^{y})$. Each of these are measurable, and they increase pointwise to $f(y) = \mu(E^{y})$. Thus, $f$ is measurable as it is the limit of measurable functions. Then, by the monotone convergence theorem
    \begin{align*}
        \int_{Y} f(y) d\nu &= \lim_{n \to \infty} \int f_{n}(y)d\nu\\
        &= \lim_{n \to \infty} (\mu \times \nu)(E_n)\\
        &= (\mu \times \nu)(E)
    \end{align*}
    To see that $E_1 \supseteq E_2 \supseteq \cdots$ implies $E = \bigcap_{n} E_n$, it is the same proof but we use DCT using the fact ath $\mu \times \nu$ is $\sigma$-finite. 

    $\qed$
\end{proof}
\tocsection{\textbf{10/9/26 Lecture 18:} Fubini-Tonelli}

Consider $(X, \cal{A}, \mu)$ and $(Y, \cal{B}, \nu)$, and consider $\cal{C}_0$, the algebra of rectangles $\bigcup_{i = 1}^{N} A_i \times B_i$. Recall that for $E \subseteq X \times Y$,
\[ E_{x} = \{y \col (x, y) \in E\} \quand E^{y} = \{x \col (x, y) \in E\} \]
and for $f : X \times Y \to \R$,
\[ f_{x}(y) = f(x, y) \quand f^{y}(x) = f(x, y) \]

\begin{theorem}[\textbf{Theorem 18.1:} Fubini-Tonelli]
    Suppose $\mu$ and $\nu$ are $\sigma$-finite and $f$ is $\mu \times \nu$ measurable. Then, if $f$ is nonnegative or integrable, then
    \[ g(x) := \int f(x, y) d\nu(y) \quand h(y) := \int f(x, y) d\mu(x) \]
    are measurable and
    \begin{align*}
        \int f(x, y) d(\mu \times \nu)(x, y) &= \iint f(x, y) d\nu(y) d\mu(x)\\
        &= \iint f(x, y) d\mu(x) d\nu(y)
    \end{align*}
\end{theorem}
\begin{proof}{Proof:}
    By Proposition 17.4 from last time, the theorem holds for simple functions. Now, if $f$ is non-negative, let $0 \leq f_n \uparrow f$ be a sequence of simple functions increasing to $f$. Let
    \[ g_n(x) := \int f_n(x, y) d\nu(y) \quand h_n(y) := \int f_n(x, y) d\mu(x) \]
    We see that as $f_n(x, y) \uparrow f(x, y)$, $g_n(x) \uparrow g(x)$ and $h_n(y) \uparrow h(y)$ by the MCT. Thus, $g$ and $h$ are measurable.\\

    Now,
    \[ \int g(x) d\mu(x) = \lim_{n \to \infty} \int g_n(x) d\mu(x) = \lim_{n \to \infty} \int f_n(x, y) d(\mu \times \nu)(x, y) = \int f(x, y) d(\mu \times \nu)(x, y) \]
    by the MCT. The analogous result for $h$ holds.\\

    Now, if $f$ is integrable (this is so that $\int f d\mu = \int f^{+} d\mu - \int f^{-} d\mu$), write $f = f^{+} - f^{-}$ and apply the nonnegative result. 

    $\qed$
\end{proof}

\begin{remark}{Remark}
    The utility of the theorem is that we can switch the order of integration:
    \[ \int g(x) d\mu(x) = \int h(y) d\nu(y) \]
\end{remark}

\begin{example}{\textbf{Example}}
    Let $A \subseteq \R$ have finite Lebesgue measure and let $f(t) = \lambda(A \cap (A + t))$. We want to show that
    \[ \int f(t) dt = (\lambda(A))^2 \]
\end{example}
\begin{solution}{Proof:}
    We can write
    \[ \int f(t) dt = \int_{\R} \left(\int_{\R} \chi_{A}(x)\chi_{A + t}(x) dx\right) dt \]
    Since $\chi_{A}(x)\chi_{A + t}(x) \geq 0$, we can switch the order of integration to get
    \begin{align*}
        \int_{\R} \left(\int_{\R} \chi_{A}(x)\chi_{A + t}(x) dx\right) dt &= \int_{\R} \left(\int_{\R} \chi_{A}(x)\chi_{A + t}(x) dt\right) dx\\
        &= \int_{\R} \chi_{A}(x)\left(\int_{\R} \chi_{x - A}(t) dt\right) dx\\
        &= \int_{\R} \chi_{A}(x)\lambda(A) dx\\
        &= \lambda(A)^2
    \end{align*}
    where we used that $\chi_{A + t}(x) = \chi_{x - A}(t)$. 
\end{solution}

\begin{example}{\textbf{Example:} Necessity of Integrability}
    Let $f : [0, \infty)^2 \to \R$ be given by
    \[ f(x, y) = \left\{\begin{matrix} 1 & x \leq y < x + 1 \\ -1 & x + 1 \leq y \leq x + 2 \\ 0 & \tt{else} \end{matrix}\right. \]
    \fig{0.5}{Lecture18Fig1}{}
    Then,
    \[ \int_{0}^{\infty} f(x, y) dy = 0 \quand \int_{0}^{\infty} f(x, y) dx = \left\{\begin{matrix} 0 & y \geq 2 \\ 2 - y & 1 \leq y \leq 2 \\ y & 0 \leq y \leq 1 \end{matrix}\right. \]
    so
    \[ \int_{0}^{\infty}\int_{0}^{\infty} f(x, y) dy dx = 0 \quand \int_{0}^{\infty}\int_{0}^{\infty} f(x, y) dx dy = 1 \]
    Luckily, $f$ is not integrable, as $\int f^{+} = \infty$ and $\int f^{-} = \infty$.
\end{example}

\tocsection{\textbf{10/12/26 Lecture 19:} Completions and Signed Measures}

Recall the product measure from last time. It is quite clear how to extend it to finitely many measure spaces by
\[ (\mu_1 \times \cdots \times \mu_n)(A_1 \times \cdots \times A_n) = \prod_{i = 1}^{n} \mu_i(A_i) \]
Unfortunately, product measures are not complete. Let $A \in \cal{A}_1$ with $A \neq \emptyset$ and $\mu_1(A) = 0$. Take a non-measurable $E \in \P(A_2)$. 
Now, $A \times E$ is non-measurable (otherwise, $E$ would be measurable by taking restrictions), but $A \times E \subseteq A \times A_2$, and $A \times A_2$ has measure $0$. Thus, $\cal{A}_1 \times \cal{A}_2$ is not complete.\\ 

Recall that, given a measure space $(X, \cal{A}, \mu)$ that is not complete, we can consider its completion $(X, \overline{\cal{A}}, \overline{\mu})$.

\begin{definition}[\textbf{Definition 19.1:} $d$-Dimensional Lebesgue Measure]
    The \emph{$d$-dimensional Lebesgue measure} is the completion of $(\R^d, \cal{A}^d, \lambda^d)$.
\end{definition}

\textbf{Note:} We can also define the Lebesgue measure on $\R^d$ using an outer measure with rectangles, and it agrees!

\subsection*{Signed Measures}
\sectionspace

\textbf{Motivation:} If $f$ is smooth and compactly supported, then
\[ f(b) - f(a) = \int_{a}^{b} f'(x) dx \]
Now, as seen in the homework, $\nu_{f'}(E) = \int_{E} f'(x) dx$ is a measure when $f' \geq 0$, so $f(b) - f(a) = \nu_{f'}([a, b])$. This is the Lebesgue-Stieltjes measure for $f$! Thus, the measure acts like $f'$ in the sense that if we integrate, we get $f$ (even if $f'$ doesn't exist!)\\

\textbf{Issue:} This only works for increasing, right-continuous functions, so if we want to be more general, we want to allow signed measures.

\begin{definition}[\textbf{Definition 19.2:} Signed Measures]
    Given a $\sigma$-algebra $\cal{A}$, a \emph{signed measure} $\mu : \cal{A} \to (-\infty, \infty]$ is a function such that
    \begin{enumerate}
        \item $\mu(\emptyset) = 0$
        \item If $\{A_i\}_{i = 1}^{\infty}$ are disjoint, then
        \[ \mu\left(\bigcup_{i = 1}^{\infty} A_i\right) = \sum_{i = 1}^{\infty} \mu(A_i) \]
        and the sum converges absolutely when the LHS is finite.
    \end{enumerate}
\end{definition}

\begin{lemma}[\textbf{Proposition 19.3}]
    Let $\mu$ be a signed measure on $(X, \cal{A})$. Then,
    \begin{enumerate}
        \item If $A_1 \subseteq A_2 \subseteq \cdots$ with $A_n \in \cal{A}$, then $\mu\left(\bigcup_{n} A_n\right) = \lim_{n} \mu(A_n)$.
        \item If $A_1 \supseteq A_2 \supseteq \cdots$ with $A_n \in \cal{A}$ and $\mu(A_1) < \infty$, then $\mu\left(\bigcap_{n} A_n\right) = \lim_{n} \mu(A_n)$.
    \end{enumerate}
\end{lemma}

\begin{example}{\textbf{Examples:} Signed Measures}
    \vspace{-0.5cm}
    \begin{itemize}
        \item If $\mu_1$ and $\mu_2$ are measures and $\mu_2$ is finite, then $\nu = \mu_1 - \mu_2$ is a signed measure.
        \item If $f$ is integrable, then $\nu(E) = \int_{E} f dx$ is a signed measure. Note that this is the same as above if we let $\mu_1(E) = \int_{E} f^{+} dx$ and $\mu_2(E) = \int_{E} f^{-} dx$.
    \end{itemize}
\end{example}

\begin{definition}[\textbf{Definition 19.4:} Positive, Negative, and Null Sets]
    Let $\mu$ be a signed measure. A set $A \in \cal{A}$ is \emph{positive} if $\mu(B) \geq 0$ for any measurable $B \subseteq A$. Similarly, $A$ is \emph{negative} if $\mu(B) \leq 0$ for any measurable $B \subseteq A$. Finally, $A$ is \emph{null} if $\mu(B) = 0$ for any measurable $B \subseteq A$.
\end{definition}

\begin{lemma}[\textbf{Theorem 19.5:} Hahn Decomposition Theorem]
    Let $\mu$ be a signed measure on $(X, \cal{A})$.
    \begin{enumerate}
        \item There exist disjoint sets $A, B \in \cal{A}$ such that $A$ is positive, $B$ is negative, and $A \cup B = X$. 
        \item If $A'$ and $B'$ are another pair of sets satisfying 1, then $A \triangle A'$ and $B \triangle B'$ are null sets.
        \item If $\mu$ is not a positive measure, then $\mu(B) < 0$, and if $-\mu$ is not positive, then $\mu(A) > 0$.
    \end{enumerate}
\end{lemma}
\begin{proof}{Proof:}
    We need a lemma first.
\end{proof}

\begin{lemma}[\textbf{Lemma 19.6}]
    Let $\mu$ be a signed measure. If $\mu(B) < 0$, then we can find $B' \subseteq B$ that is negative such that $\mu(B') < 0$. 
\end{lemma}
\begin{proof}{Proof:}
    Suppose $B$ is not negative (the result is obvious if it is). The idea is that we pull positive pieces off of $B$.\\

    As $B$ is not negative, find the smallest integer $n_1$ such that there exists $A_1 \subseteq B$ with $\mu(A_1) \geq 1/n_1$. Now, let
    \[ B_1 = B \setminus A_1 \]
    If $B_1$ is negative, we are done, as $\mu(B_1) = \mu(B) - \mu(A_1) < 0$. If not, iterate. Given $B_j$ not negative, find $A_{j + 1} \subseteq B_j$ such that $\mu(A_{j + 1}) \geq 1/n_{j + 1}$, where $n_{j + 1}$ is the smallest such integer. Then, define % ?? page has \mu(A_{j+1}) = 1/n_{j+1}; I wrote \geq to match the first step
    \[ B_{j + 1} = B_j \setminus A_{j + 1} \]
    Now, let $B' = \bigcap_{j} B_j$. Then, $B = B' \cup \left(\bigcup_{j = 1}^{\infty} A_j\right)$. Since this is a disjoint union,
    \[ \mu(B) = \mu(B') + \sum_{i = 1}^{\infty} \mu(A_i) \]
    Since $\mu(B) < 0$, we must have $\mu(B') < 0$. Also, as $\mu(B') < 0$, $\mu(A_j) \to 0$, which implies $n_j \to \infty$.\\

    Now, we need to check that $B'$ is a negative set.
\end{proof}

\tocsection{\textbf{10/14/26 Lecture 20:} The Hahn and Jordan Decompositions}

\begin{proof}{Proof of Lemma 19.6 (cont.):}
    Now, we need to check that $B'$ is a negative set. Suppose not. Then, there exists $P \subseteq B'$ such that $\mu(P) > 0$. Then, there exists $N$ such that $0 \leq 1/N < \mu(P)$. Now, since $n_j \to \infty$, there exists $J$ such that $1/n_J < 1/N$. However, this is a contradiction, as $n_J$ was chosen to be the smallest such integer for a set in $B_{J - 1}$, but there exists a set in $B_{J - 1}$ with measure greater than $1/N$.

    $\qed$
\end{proof}
Now, we can prove the Hahn Decomposition theorem.
\vspace{0.5cm}

\begin{theorem}[\textbf{Theorem 20.1:} Hahn Decomposition Theorem]
    Let $\mu$ be a signed measure on $(X, \cal{A})$.
    \begin{enumerate}
        \item There exist disjoint sets $A, B \in \cal{A}$ such that $A$ is positive, $B$ is negative, and $A \cup B = X$. 
        \item If $A'$ and $B'$ are another pair of sets satisfying 1, then $A \triangle A'$ and $B \triangle B'$ are null sets.
        \item If $\mu$ is not a positive measure, then $\mu(B) < 0$, and if $-\mu$ is not positive, then $\mu(A) > 0$.
    \end{enumerate}
\end{theorem}
\begin{proof}{Proof:}
    \textbf{(1)} Let
    \[ L = \inf\{\mu(B) \col B \tt{ is a negative set}\} \]
    Choose $B_n$ so that $\mu(B_n) \downarrow L$. Now, let $B_1' = B_1$ and $B_n' = B_n \setminus \bigcup_{j = 1}^{n - 1} B_j'$. Now, $B = \bigcup_{n} B_n = \bigsqcup_{n} B_n'$. Then, for $E \subseteq B$,
  .
    \[ \mu(E) = \lim_{n \to \infty} \mu \left(E \cap \bigcup_{j = 1}^{n} B_j'\right) = \lim_{n \to \infty} \sum_{j = 1}^{n} \mu(E \cap B_j') \leq 0 \]
    so $B$ is negative. Now, because $B$ is negative,
    \[ \mu(B) = \mu(B_n) + \mu(B \setminus B_n) \leq \mu(B_n) \]
    so taking the limit, we see that $\mu(B) \leq L$, so as $\mu(B) \geq L$, we have equality. Note that this also shows that this shows $L > -\infty$ as $\mu(B) > -\infty$.\\

    Let $A = B^{c}$. We want to show that $A$ is positive. If not, there exists $E \subseteq A$ such that $\mu(E) < 0$, so we can also find a negative $E' \subseteq E$ such that $\mu(E') < 0$. But then
    \[ \mu(B \cup E') = \mu(B) + \mu(E') < L \]
    which is a contradiction, as $B \cup E'$ is negative.\\

    \textbf{(2)} Let $E \subseteq A \setminus A' = B' \setminus B$. Then, as $A$ is positive and $B'$ is negative, $\mu(E) \geq 0$ and $\mu(E) \leq 0$. A similar fact holds for $E \subseteq A' \setminus A = B \setminus B'$. Thus, $A \triangle A'$ and $B \triangle B'$ are null.

    \textbf{(3)} Suppose $\mu$ is not positive. Then, there exists $E$ such that $\mu(E) < 0$. Then,
    \begin{align*}
        \mu(E) &= \mu(E \cap A) + \mu(E \cap B)\\
        &\geq \mu(E \cap B)\\
        &\geq \mu(B) && (B \tt{ is negative})
    \end{align*}
    Thus, $0 > \mu(E) \geq \mu(B)$, so $\mu(B) < 0$. Similarly, if $\mu$ is not negative, there exists $E$ such that $\mu(E) > 0$, so
    \begin{align*}
        \mu(E) &= \mu(E \cap A) + \mu(E \cap B)\\
        &\leq \mu(E \cap A) \leq \mu(A)
    \end{align*}
    so $0 < \mu(E) \leq \mu(A)$.

    $\qed$
\end{proof}

\begin{definition}[\textbf{Definition 20.2:} Mutually Singular Measures]
    Two measures $\mu$ and $\nu$, both on $(X, \cal{A})$, are \emph{mutually singular} if there exist disjoint $A, B$ such that $A \cup B = X$ and $\mu(B) = \nu(A) = 0$. We denote this by $\mu \perp \nu$.
\end{definition}

\begin{example}{\textbf{Example 1}}
    On $\R$, consider the Lebesgue measure and the point mass at $0$. Then, $\R = (\R \setminus \{0\}) \sqcup \{0\}$ and
    \[ \delta_0(\R \setminus \{0\}) = \lambda(\{0\}) = 0 \]
\end{example}

\begin{example}{\textbf{Example 2}}
    Recall the Cantor function $f$:
    \fig{0.4}{Lecture20Fig1}{}
    Let $m$ be the Lebesgue-Stieltjes measure of $f$. We claim that $\lambda$ and $m$ are mutually singular. The idea is that $m$ only assigns measure to the Cantor set $C$. Then,
    \[ m(C^{c}) = \lambda(C) = 0 \]
\end{example}

\begin{theorem}[\textbf{Theorem 20.3:} Jordan Decomposition Theorem]
    If $\mu$ is a signed measure, then there are positive measures $\mu^{+}$ and $\mu^{-}$, with $\mu^{-}$ finite, so that $\mu = \mu^{+} - \mu^{-}$ and $\mu^{+} \perp \mu^{-}$, and these measures are unique.
\end{theorem}
\begin{proof}{Proof:}
    Apply the Hahn decomposition to get a positive set $A$ and a negative set $B$. Then, let
    \[ \mu^{+}(E) = \mu(A \cap E) \quand \mu^{-}(E) = -\mu(E \cap B) \]
    These are positive measures, and
    \[ \mu(E) = \mu(A \cap E) + \mu(E \cap B) = \mu^{+}(E) - \mu^{-}(E) \]
    so $\mu = \mu^{+} - \mu^{-}$. It is clear that $\mu^{+} \perp \mu^{-}$ as $\mu^{-}(A) = \mu^{+}(B) = 0$ and $A \cup B = X$. We also see that $\mu^{-}$ is finite as $\mu^{-}(X) = -\mu(B) < \infty$. \\

    Now, we check uniqueness. Suppose $\mu = \nu^{+} - \nu^{-}$. We want to show that $\nu^{+} = \mu^{+}$ and $\nu^{-} = \mu^{-}$. We know $\nu^{+} \perp \nu^{-}$, so pick $A'$ and $B'$ so that $\nu^{+}(B') = \nu^{-}(A') = 0$. Then, for $E \subseteq A'$,
    \[ \mu(E) = \nu^{+}(E) \geq 0 \]
    so $A'$ is positive for $\mu$, and similarly $B'$ is negative for $\mu$. Thus, $A'$ and $B'$ form a Hahn decomposition for $\mu$, so $A \triangle A'$ is null. Now,
    \begin{align*}
        \nu^{+}(E) = \nu^{+}(E \cap A') &= \nu^{+}(E \cap A') - \nu^{-}(E \cap A')\\
        &= \mu(E \cap A')\\
        &= \mu(E \cap A)\\
        &= \mu^{*}(E)
    \end{align*}
    so $\mu^{+} = \nu^{+}$. Then, as $\nu^{-} = \nu^{+} - \mu = \mu^{+} - \mu = \mu^{-}$, we are done.

    $\qed$
\end{proof}

\tocsection{\textbf{10/16/26 Lecture 21:} The Radon-Nikodym Theorem}

If $f$ is integrable with respect to $\mu$, we can define
\[ \nu(A) = \int_{A} f d\mu \]
In a sense, ``$f$ is the derivative of $\nu$ with respect to $\mu$''.\\

\textbf{Question:} When does $\nu$ have a derivative with respect to $\mu$?\\

\textbf{Answer:} Note that $\mu(A) = 0 \implies \nu(A) = 0$. This is necessary, but it turns out it is sufficient as well.

\begin{definition}[\textbf{Definition 21.1:} Absolute Continuity]
    We say that $\nu$ is \emph{absolutely continuous with respect to $\mu$} if $\nu(A) = 0$ whenever $\mu(A) = 0$, denoted $\nu \ll \mu$.
\end{definition}

\begin{lemma}[\textbf{Proposition 21.2}]
    If $\nu \ll \mu$ and $\nu \perp \mu$, then $\nu = 0$.
\end{lemma}
\begin{proof}{Proof:}
    Find $A, B$ such that $\nu(A) = \mu(B) = 0$. Then, $\nu(B) = 0$. Thus, as $\nu(A) + \nu(B) = \nu(X)$, we see that $\nu(X) = 0$.

    $\qed$
\end{proof}

\begin{lemma}[\textbf{Proposition 21.3}]
    Let $\nu$ be a finite measure. Then, $\nu \ll \mu$ if, and only if, for any $\epsilon > 0$, there exists $\delta > 0$ so that
    \[ \mu(A) < \delta \implies \nu(A) < \epsilon \]
\end{lemma}
\begin{proof}{Proof:}
    $(\Leftarrow)$ If $\mu(A) = 0$, then $\mu(A) < \delta$ for any $\delta$, so $\nu(A) < \epsilon$ for any $\epsilon$, so $\nu(A) = 0$.\\

    $(\Rightarrow)$ Suppose not. Then, for some $\epsilon > 0$ and every $k$, there exists $A_k$ such that $\mu(A_k) < 2^{-k}$ but $\nu(A_k) > \epsilon$. Let
    \[ B = \bigcap_{n = 1}^{\infty} \bigcup_{k = n}^{\infty} A_k \]
    Then,
    \[ \mu(B) = \lim_{n \to \infty} \mu\left(\bigcup_{k = n}^{\infty} A_k\right) \leq \lim_{n \to \infty} \sum_{k = n}^{\infty} 2^{-k} = 0 \]
    Now, $\mu(B) = 0$, but as $\nu$ is finite,
    \[ \nu(B) = \lim_{n \to \infty} \nu\left(\bigcup_{k = n}^{\infty} A_k\right) \geq \epsilon \]
    This contradicts absolute continuity. Note that we used the finiteness of $\nu$ to make sure these descending limits exist. 

    $\qed$
\end{proof}


\begin{theorem}[\textbf{Theorem 21.4:} Radon-Nikodym Theorem]
    Let $\mu$ be $\sigma$-finite and $\nu$ be a finite measure. Suppose $\nu \ll \mu$. Then, there is an integrable $f \geq 0$ so that
    \[ \nu(A) = \int_{A} f d\mu \]
    and $f$ is $\mu$-a.e. unique. We call $f$ the \emph{Radon-Nikodym derivative} of $\nu$, and denote it $\frac{d\nu}{d\mu}$.
\end{theorem}
\begin{proof}{Proof:}
    We need a quick lemma.
\end{proof}

\begin{lemma}[\textbf{Lemma 21.5}]
    Let $\mu$ and $\nu$ be finite measures. If $\mu$ and $\nu$ are not mutually singular, then there is some $\epsilon > 0$ and some set $A$ so that $\mu(A) > 0$ and $A$ is a positive set for $\nu - \epsilon\mu$. That is, $\nu(B) \geq \epsilon\mu(B)$ for $B \subseteq A$.
\end{lemma}
\begin{proof}{Proof:}
    Apply the Hahn decomposition to $\nu - \frac{1}{n}\mu$ to get positive sets $F_n$ and negative sets $E_n$. Let $E = \bigcap_{n} E_n$. For each $n$,
    \[ \nu(E) \leq \nu(E_n) \leq \frac{1}{n}\mu(E_n) \leq \frac{1}{n}\mu(X) \]
    Taking $n \to \infty$, we see that $\nu(E) = 0$.\\

    Consider $\mu(E^{c})$. We can't have $\mu(E^{c}) = 0$, as this would imply that $\mu$ and $\nu$ are mutually singular. Thus, $\mu(E^{c}) > 0$. We see $E^{c} = \bigcup_{n} E_n^{c}$, so for some $n$, $\mu(E_n^{c}) > 0$. Let $A = E_n^{c} = F_n$ and let $\epsilon = \frac{1}{n}$. Then, $\nu - \frac{1}{n}\mu$ is positive on $A$ and $\mu(A) > 0$.

    $\qed$
\end{proof}
\vspace{0.5cm}

\begin{proof}{Proof of Theorem 21.4 (Radon-Nikodym):}
    Start by assuming $\mu$ is finite. Let
    \[ F = \left\{g : X \to [0, \infty) \col \int_{A} g d\mu \leq \nu(A)\ \forall A \in \cal{A}\right\} \]
    $F$ is nonempty because it contains $0$.\\

    Now, let $L = \sup\{\int g d\mu \col g \in F\}$. Choose $f_n$ so that $\int f_n d\mu \to L$, and let $g_n = f_1 \vee \cdots \vee f_n$. We want to show that $g_n \in F$. If $h, g \in F$, then
    \begin{align*}
        \int_{A} (h \vee g) d\mu &= \int_{A \cap \{h \geq g\}} (h \vee g) d\mu + \int_{A \cap \{h < g\}} (h \vee g) d\mu\\
        &\leq \int_{A \cap \{h \geq g\}} h d\mu + \int_{A \cap \{h < g\}} g d\mu\\
        &\leq \nu(A \cap \{h \geq g\}) + \nu(A \cap \{h < g\})\\
        &= \nu(A)
    \end{align*}
    so $h \vee g \in F$. Thus, by induction, $g_n \in F$ for all $n$. Now, let $f = \lim_{n \to \infty} g_n$. By the MCT,
    \[ \int_{A} f d\mu = \lim_{n \to \infty} \int_{A} g_n d\mu \leq \nu(A) \]
    so $f \in F$. Also, $\int_{X} f d\mu = L$, by the MCT.\\

    Now, we want to show that $\nu(A) = \int_{A} f d\mu$. Let
    \[ \tilde{\nu}(A) = \nu(A) - \int_{A} f d\mu \]
    We see that $\tilde{\nu}$ is a measure, so we want to show it is zero.\\

    We know $\tilde{\nu} \ll \mu$, so we want to show $\tilde{\nu} \perp \mu$. Suppose not. By our lemma, we can find $A$ and $\epsilon$ so that $A$ is positive for $\tilde{\nu} - \epsilon\mu$ and $\mu(A) > 0$. Then, for $B \subseteq A$,
    \[ \nu(B) \geq \int_{B} f d\mu + \epsilon\mu(B) \]
    Because of this, we know that $f + \epsilon\chi_{A} \in F$, as
    \begin{align*}
        \int_{B} (f + \epsilon\chi_{A}) d\mu &= \int_{B} f d\mu + \epsilon\mu(A \cap B)\\
        &\leq \int_{B} f d\mu + \nu(A \cap B) - \int_{A \cap B} f d\mu\\
        &= \int_{B \setminus A} f d\mu + \nu(A \cap B)\\
        &\leq \nu(B \setminus A) + \nu(A \cap B) = \nu(B)
    \end{align*}
    However,
    \[ \int (f + \epsilon\chi_{A}) d\mu = \int f d\mu + \epsilon\mu(A) = L + \epsilon\mu(A) > L \]
    This contradicts the maximality of $L$. Thus, $\tilde{\nu} \perp \mu$, so $\tilde{\nu} = 0$ by Proposition 21.2. Thus,
    \[ \nu(A) = \int_{A} f d\mu \]

    To show uniqueness, suppose $\nu(A) = \int_{A} g d\mu$. Then,
    \[ \int_{A} (f - g) d\mu = 0 \]
    for all $A$. Then, by a previous result (Proposition 13.4), we know $f - g = 0$ a.e., so $f = g$ a.e.\\

    To extend to $\sigma$-finite $\mu$, write $X = \bigsqcup_{n} X_n$ so that each $X_n$ has finite measure. Then, apply the above to $\mu\res{X_n}$. Thus, we get a function $f_n$ on each $X_n$ with $\nu(A) = \int_{A} f_n d\mu$ for $A \subseteq X_n$. Now, defining $f = \sum_{n} f_n \chi_{X_n}$, we have that 
    \[ \nu(A) = \int_{A} f d\mu \]
    for any $A$ by countable additivity and the MCT, integrability holds as $\nu$ is finite, and uniqueness one ach $X_n$ extends to uniqueness everywhere. 

    $\qed$
\end{proof}

\tocsection{\textbf{10/19/26 Lecture 22:} Lebesgue Decomposition and Differentiation}

\begin{theorem}[\textbf{Theorem 22.1:} Lebesgue Decomposition Theorem]
    Let $\mu$ be positive and $\sigma$-finite and $\nu$ positive and finite. Then, there exists $\rho$ positive satisfying $\rho \perp \mu$ and $f \geq 0$ integrable such that
    \[ \nu(A) = \int_{A} f d\mu + \rho(A) \]
\end{theorem}
\begin{proof}{Proof:}
    Assume $\mu$ is finite and let $f$ be as in the proof of the Radon-Nikodym theorem; that is, such that $\int f d\mu = L = \sup\{\int g d\mu \col g \in F\}$, where
    \[ F = \left\{g : X \to [0, \infty) \col \int_{A} g\ d\mu \leq \nu(A)\ \forall A \in \cal{A}\right\} \]
    Define $\rho(A) = \nu(A) - \int_{A} f d\mu$. We see that $\rho(A)$ is positive from the fact $f \in F$.\\

    We want $\rho \perp \mu$. If not, by Lemma 21.5 we can find a set $A$ and $\epsilon > 0$ such that $\mu(A) > 0$ and $A$ is positive for $\rho - \epsilon\mu$. Then, for $B \subseteq A$,
    \[ \epsilon\mu(B) \leq \nu(B) - \int_{B} f d\mu \]
    Now, consider $f + \epsilon\chi_A$. Given $B$, we see that
    \begin{align*}
        \int_{B} f + \epsilon\chi_A d\mu &= \int_{B \setminus A} f d\mu + \int_{B \cap A} f + \epsilon\chi_A d\mu\\
        &\leq \nu(B \setminus A) + \nu(B \cap A)\\
        &= \nu(B)
    \end{align*}
    where for the second term we used that $\nu(B) \geq \int_{B} f d\mu + \epsilon\mu(B)$ for $B \subseteq A$.\\

    Thus, $f + \epsilon \chi_{A} \in F$, which contradicts that $\int f d\mu = L$ is the sup over integrals in $F$. Thus, we must have $\rho \perp \mu$, which was to be shown.

    $\qed$
\end{proof}

\begin{example}{\textbf{Example:} Necessity of Absolute Continuity}
    We need absolute continuity to have R-N derivatives. For instance, if we wanted $\frac{d\delta_0}{d\mu}$, where $\mu$ is the Lebesgue measure, we would want an $f$ such that
    \[ \delta_0(\{0\}) = \int_{\{0\}} f d\mu = 1 \]
    but no $f$ exists since $\{0\}$ has measure $0$.
\end{example}

\begin{lemma}[\textbf{Corollary 22.2}]
    Let $\nu$ be a finite measure and $\mu$ $\sigma$-finite with $\nu \ll \mu$. If $f$ is integrable with respect to $\nu$, then $f \frac{d\nu}{d\mu}$ integrable also with respect to $\mu$ with
    \[ \int f d\nu = \int f \frac{d\nu}{d\mu} d\mu \]
\end{lemma}
\begin{proof}{Proof:}
    For a simple function $\chi_A$,
    \[ \int \chi_A \frac{d\nu}{d\mu} d\mu = \int_{A} \frac{d\nu}{d\mu} d\mu = \nu(A) = \int \chi_A d\nu \]
    By linearity, we get simple functions and by monotone convergence, we get non-negative and integrable functions.

    $\qed$
\end{proof}

\begin{definition}[\textbf{Definition 22.3:} Lipschitz Functions]
    $f : \R \to \R$ is \emph{Lipschitz} if there exists $L$ such that $\abs{f(x) - f(y)} \leq L\abs{x - y}$ for $x, y \in \R$.
\end{definition}

\begin{example}{\textbf{Example:} $C^1$ Functions}
    If $f \in C^1$ and $\abs{f'} \leq L$, then $f$ is $L$-Lipschitz. The converse is not true. For example, $\abs{x}$ is $1$-Lipschitz.
\end{example}

\begin{lemma}[\textbf{Proposition 22.4}]
    Let $f : \R \to \R$ be Lipschitz. Then, $f$ is differentiable a.e.
\end{lemma}
\begin{proof}{Proof:}
    We will soon prove the following fact that we will use here: If $f(x) = \int_{a}^{x} g(y) dy$, then $f$ is differentiable a.e. and $f' = g$.\\

    Let $L$ be such that $\abs{f(x) - f(y)} \leq L\abs{x - y}$. First, suppose $f$ is increasing and let $m$ be its L-S measure. Then
    \begin{align*}
        m((a, b]) &= f(b) - f(a)\\
        &\leq L(b - a)
    \end{align*}
    This implies $m \ll \lambda$. Then, $m(A) = \int_{A} g(y) dy$ for some $g$. Thus,
    \[ m((a, x]) = \int_{a}^{x} g(y) dy = f(x) - f(a) \]
    so $f(x) - f(a)$ is differentiable a.e.\\

    If $f$ is not increasing, then $f + Lx$ must be. This is because if $x > y$,
    \begin{align*}
        (f(x) + Lx) - (f(y) + Ly) &= (f(x) - f(y)) + L(x - y)\\
        &\geq 0 \quad \tt{as } f(x) - f(y) \geq -L(x - y)
    \end{align*}
    Thus, we see $f + Lx$ is differentiable a.e. by above, so $f$ must be as well.

    $\qed$
\end{proof}

\subsection*{Differentiation}
\sectionspace

In this section, we want to show that the derivative of $\int_{a}^{x} f(y) d\mu(y)$ is equal to $f(x)$ $\mu$-a.e. Even stronger, we want to study
\[ f_r(x) = \frac{1}{\lambda(B_r(x))} \int_{B_r(x)} f(y) dy \]
where $B_r(x) = \{y \col d(x, y) < r\}$.\\

\textbf{Q:} Does $f_r \to f$ as $r \to 0$? In 1d, we have
\[ f_r = \frac{1}{2r} \int_{x - r}^{x + r} f(y) dy \]

First, suppose $g$ is continuous. Then,
\begin{align*}
    \abs{g_r(x) - g(x)} &= \abs{\frac{1}{\lambda(B_r(x))} \int_{B_r(x)} g(y) - g(x) dy}\\
    &\leq \frac{1}{\lambda(B_r(x))} \int_{B_r(x)} \abs{g(y) - g(x)} dy
\end{align*}
Since $g$ is continuous, we see that for any $\epsilon > 0$, there exists $r$ such that $y \in B_r(x) \implies \abs{g(y) - g(x)} < \epsilon$. Then, for such $r$,
\[ \abs{g_r(x) - g(x)} \leq \frac{1}{\lambda(B_r(x))} \int_{B_r(x)} \abs{g(y) - g(x)} dy < \epsilon \]
so taking $\epsilon \to 0$, we see $g_r(x) \to g(x)$. Thus, if $g$ is continuous, we are done.\\

For general integrable $f$, we want to pick $g$ that is continuous and compactly supported such that
\[ \int \abs{f - g} dy < \epsilon \]
Thus,
\[ \abs{f_r(x) - f(x)} \leq \abs{f_r(x) - g_r(x)} + \abs{g_r(x) - g(x)} + \abs{g(x) - f(x)} \]
We want to study the measure of the set where $\abs{f_r(x) - f(x)} > \frac{1}{n}$. 
Note that we can control $\abs{g(x) - f(x)}$ by Chebyshev's and $\abs{g_r(x) - g(x)}$ by above, so we really want to study $\abs{f_r(x) - g_r(x)}$, which is at most
\[ \frac{1}{\lambda(B_r(x))} \int_{B_r(x)} \abs{f(y) - g(y)} dy \]
We want to show that this is large on a small set using the fact that $\abs{f(y) - g(y)}$ is small.

\begin{definition}[\textbf{Definition 22.5:} Hardy-Littlewood Maximal Function]
    We define
    \[ M(h(x)) = \sup_{r > 0} \frac{1}{\lambda(B_r(x))} \int_{B_r(x)} \abs{h(y)} dy \]
\end{definition}

We want to show that $Mh$ is controlled by $\int \abs{h}$.

\tocsection{\textbf{10/21/26 Lecture 23:} The Maximal Function and Lebesgue Differentiation}

\textbf{Recall:} For $f$ integrable, $f : \R^d \to \R$,
\[ f_r(x) = \frac{1}{\lambda(B_r(x))} \int_{B_r(x)} f(y) dy \]
Does $f_r(x) \to f(x)$ as $r \to 0$?\\

\textbf{In general:} No. Let $f(x) = \chi_{C}$ for $C$ the Cantor set. Then $f_r(x) = 0$. So $f(x) \neq f_r(x)$, but note $f(x) = 0$ a.e.

\begin{definition}[\textbf{Definition 23.1:} Locally Integrable Functions]
    We say $f$ measurable is \emph{locally integrable} if $\int_{K} \abs{f} < \infty$ for any compact set $K$, and we write $f \in L^1_{\tt{loc}}$.
\end{definition}

\begin{lemma}[\textbf{Theorem 23.2:} Lebesgue Differentiation Theorem]
    If $f \in L^1_{\tt{loc}}$, then $f_r(x) \to f(x)$ a.e. as $r \to 0$.
\end{lemma}

We will work toward proving this, but not yet.

\begin{lemma}[\textbf{Proposition 23.3}]
    If $f \in L^1_{\tt{loc}}$, then $Mf(x)$ is measurable.
\end{lemma}
\begin{proof}{Proof:}
    Given sets $A$ and $B$,
    \begin{align*}
        \abs{\int_{A} f - \int_{B} f} &= \abs{\int_{A} f - \left(\int_{A \cap B} f + \int_{B \setminus A} f\right)}\\
        &= \abs{\int_{A \setminus B} f - \int_{B \setminus A} f}\\
        &\leq \int_{A \setminus B} \abs{f} + \int_{B \setminus A} \abs{f}\\
        &= \int_{A \triangle B} \abs{f}
    \end{align*}
    Thus, for $r$ fixed,
    \[ \abs{\frac{1}{\lambda(B_r(x))} \int_{B_r(x)} f(y) dy - \frac{1}{\lambda(B_r(x'))} \int_{B_r(x')} f(y) dy} \leq \frac{1}{\lambda(B_r(0))} \int_{B_r(x) \triangle B_r(x')} \abs{f} dy \]
    As $x' \to x$, $\lambda(B_r(x) \triangle B_r(x')) \downarrow 0$, so the integral goes to $0$ by DCT. Thus,
    \[ \frac{1}{\lambda(B_r(x))} \int_{B_r(x)} \abs{f} dy \]
    is continuous in $x$. Thus, it is measurable, and as the sup of measurable functions is measurable, $Mf(x)$ is measurable.

    $\qed$
\end{proof}

\begin{lemma}[\textbf{Lemma 23.4:} Wiener Covering Lemma]
    Let $\{B_{\alpha}\}_{\alpha \in I}$ be a collection of balls in $\R^d$ with $\lambda\left(\bigcup_{\alpha} B_{\alpha}\right) < \infty$. Then, for every $\epsilon > 0$, there exist finitely many disjoint $B_1, \ldots, B_n$ such that
    \[ \sum_{j = 1}^{n} \lambda(B_j) \geq 3^{-d}(1 - \epsilon)\lambda\left(\bigcup_{\alpha} B_{\alpha}\right) \]
\end{lemma}
\begin{proof}{Proof:}
    By inner regularity, we can find $K$ compact so that $\lambda(K) > (1 - \epsilon)\lambda\left(\bigcup_{\alpha} B_{\alpha}\right)$ and $K \subseteq \bigcup_{\alpha} B_{\alpha}$. Then, by compactness, we can find $\tilde{B}_1, \ldots, \tilde{B}_m$ that cover $K$.\\

    Let $B_1$ be the $\tilde{B}_j$ with the largest radius. Then, let $B_j$ be the $\tilde{B}_j$ with the largest radius that does not intersect the $B_1, \ldots, B_{j - 1}$. Do this until we can no longer pick $B_j$. Let $B_1, \ldots, B_n$ be our collection.\\

    Now, for every $\tilde{B}_j$, we have that it is in our collection or intersects a ball in our collection. If it does intersect another ball in our collection, let $B_k$ be the smallest index such that $\tilde{B}_j \cap B_k \neq \emptyset$. This implies $\tilde{B}_j$ does not intersect $B_1, \ldots, B_{k - 1}$.\\

    Thus, we must have $\tt{rad}(\tilde{B}_j) \leq \tt{rad}(B_k)$ by our method of selection. Thus, writing $B_k = B_{r_k}(x)$ and $B_k^{*} = B_{3r_k}(x)$,\\

    \textbf{Claim:} $\tilde{B}_j \subseteq B_k^{*}$.
.
    \fig{0.42}{Lecture23Fig1}{}
    Thus, we see
    \[ (1 - \epsilon)\lambda\left(\bigcup_{\alpha} B_{\alpha}\right) < \lambda(K) \leq \sum_{j = 1}^{n} \lambda(B_j^{*}) = \sum_{j = 1}^{n} 3^d\lambda(B_j) \]
    which gives us the inequality.

    $\qed$
\end{proof}

\begin{theorem}[\textbf{Theorem 23.5:} Weak 1-1 Inequality for $M$]
    There is a constant $C > 0$ so that for all $f \in L^1$ and any $\alpha > 0$,
    \[ \lambda(\{x \col Mf > \alpha\}) \leq \frac{C}{\alpha}\int \abs{f} dx \]
\end{theorem}
\begin{proof}{Proof:}
    Let $A_{\alpha} = \{x \col Mf > \alpha\}$. Then, for each $x \in A_{\alpha}$, we can find a ball $B_{r_x}(x)$ such that
    \[ \frac{1}{\lambda(B_{r_x}(x))} \int_{B_{r_x}(x)} \abs{f} dy > \alpha \]

    By our lemma, we can find finitely many disjoint balls $B_1, \ldots, B_n$ such that
    \[ 3^{-d}(1 - \epsilon)\lambda(A_{\alpha}) \leq \sum_{j = 1}^{n} \lambda(B_j) \]
    Thus,
    \begin{align*}
        (1 - \epsilon)\lambda(A_{\alpha}) \leq 3^d\sum_{j = 1}^{n} \lambda(B_j) &\leq \frac{3^d}{\alpha}\sum_{j = 1}^{n} \int_{B_j} \abs{f} dy\\
        &\leq \frac{3^d}{\alpha}\int \abs{f} dy
    \end{align*}
    because the $B_j$ are disjoint. Then, taking $\epsilon \to 0$, we get
    \[ \lambda(A_{\alpha}) \leq \frac{3^d}{\alpha}\int \abs{f} dy \]
    which was to be shown!

    $\qed$
\end{proof}

\textbf{Remark:} Suppose we knew $\int Mf dx \leq C\int \abs{f} dx$; this is a strong 1-1 bound. Then, by Chebyshev,

\[ \lambda\{Mf > \alpha\} \leq \frac{\int Mf}{\alpha} \leq \frac{C}{\alpha}\int \abs{f} \]
which is our weak bound. However, this is typically not true.\\

Consider $d = 1$ and let $f = \chi_{[0, 1]}(x)$.
\fig{0.38}{Lecture23Fig2}{}

then, we see that for $x > 1$,
\[ \sup_{r > 0} \frac{1}{2r}\int_{B_r(x)} \chi_{[0, 1]}(y) dy \geq \frac{1}{2x} \]
so for $x > 1$, $Mf \geq \frac{1}{2x}$, which is not integrable. Thus, the strong 1-1 inequality can't hold.\\

Now, we go back to the theorem we want to prove.
\vspace{0.2cm}

\begin{proof}{Proof of Theorem 23.2:}
    It suffices to consider $f\chi_{\abs{x} \leq N}$, since we only care about the limit behavior of $f_r$. Thus, assume $f$ is integrable. Given $\epsilon > 0$, find $g$ continuous and compactly supported so that $\int \abs{f - g} < \epsilon$. Then,
    \begin{align*}
        \limsup_{r \to 0} \abs{f_r(x) - f(x)} &\leq \limsup_{r \to 0} \abs{(f - g)_r(x) + g_r(x) - g(x) + g(x) - f(x)}\\
        &\leq \limsup_{r \to 0} \left(\abs{(f - g)_r(x)} + \abs{g_r(x) - g(x)} + \abs{g(x) - f(x)}\right)\\
        &\leq M(f - g)(x) + 0 + \limsup_{r \to 0} \abs{g(x) - f(x)}
    \end{align*}

    Let $E_{\alpha} = \{\limsup_{r \to 0} \abs{f_r(x) - f(x)} > \alpha\}$ and $F_{\alpha} = \{\abs{f(x) - g(x)} > \alpha\}$. Then,
    \[ E_{\alpha} \subseteq F_{\alpha/2} \cup \{M(f - g) > \alpha/2\} \]
    We will finish next lecture using the fact that $F_{\alpha}$ is controlled by Chebyshev and $\{M(f - g) > \alpha/2\}$ is controlled by weak 1-1.
\end{proof}

\tocsection{\textbf{10/23/26 Lecture 24:} Lebesgue Differentiation and Its Consequences}

\begin{theorem}[\textbf{Theorem 24.1:} Lebesgue Differentiation Theorem]
    If $f \in L^1_{\tt{loc}}$, then $\lim_{r \to 0} f_r(x) = f(x)$ a.e.
\end{theorem}
\begin{proof}{Proof:}
    Recall that it suffices to take $f \in L^1$. We see
    \[ \{x \col f_r(x) \not\to f(x)\} \subseteq \bigcup_{n} A_n \]
    where $A_n = \{x \col \limsup_{r \to 0} \abs{f_r(x) - f(x)} > 1/n\}$. We want to show $\lambda(A_n) = 0$ for all $n$. Given $\epsilon > 0$, there exists $g$ continuous and integrable such that
    \[ \int \abs{f - g} dx < \epsilon \]
    Then,
    \begin{align*}
        \limsup_{r \to 0} \abs{f_r(x) - f(x)} &\leq \limsup_{r \to 0} \left(\abs{f_r(x) - g_r(x)} + \abs{g_r(x) - g(x)} + \abs{g(x) - f(x)}\right)\\
        &\leq M\abs{f - g}(x) + 0 + \abs{f(x) - g(x)}
    \end{align*}
    Then,
    \[ A_n \subseteq \left\{x \col M\abs{f - g}(x) > \frac{1}{2n}\right\} \cup \left\{x \col \abs{f(x) - g(x)} > \frac{1}{2n}\right\} \]
    By the weak 1-1 inequality,
    \[ \lambda\left\{x \col M\abs{f - g}(x) > \frac{1}{2n}\right\} \leq \frac{C}{1/(2n)}\int \abs{f - g} dx \leq 2nC\epsilon \]
    and by Chebyshev,
    \[ \lambda\left\{x \col \abs{f(x) - g(x)} > \frac{1}{2n}\right\} \leq \frac{1}{1/(2n)}\int \abs{f - g} dx \leq 2n\epsilon \]
    Then, as $\epsilon$ is independent of $n$ and $C$, we see that by taking $\epsilon \to 0$, $\lambda(A_n) = 0$.

    $\qed$
\end{proof}

Now, we derive some interesting theorems.\\

Recall that, from the homework, if $A \subseteq \R$ such that $\lambda(A) > 0$, then for all $\alpha \in (0, 1)$ there exists an interval $I$ such that $\lambda(A \cap I) > \alpha\lambda(I)$. It turns out this comes for free!

\begin{lemma}[\textbf{Corollary 24.2:} Lebesgue Density Theorem]
    Suppose $A \subseteq \R$ is measurable such that $\lambda(A) > 0$. Then, for a.e. $x \in A$,
    \[ \lim_{r \to 0} \frac{\lambda(A \cap B_r(x))}{\lambda(B_r(x))} = 1 \]
    and for a.e. $x \notin A$,
    \[ \lim_{r \to 0} \frac{\lambda(A \cap B_r(x))}{\lambda(B_r(x))} = 0 \]
\end{lemma}
\begin{proof}{Proof:}
    Apply the theorem to $f = \chi_A$, as
    \[ \int_{B_r(x)} \chi_A = \lambda(A \cap B_r(x)) \]

    $\qed$
\end{proof}

From our theorem, we also see
\[ \lim_{r \to 0} \frac{1}{\lambda(B_r(x))} \int_{B_r(x)} f(y) - f(x) dy = 0 \]

\begin{lemma}[\textbf{Proposition 24.3}]
    If $f \in L^1_{\tt{loc}}$, then
    \[ \lim_{r \to 0} \frac{1}{\lambda(B_r(x))} \int_{B_r(x)} \abs{f(y) - f(x)} dy = 0 \quad \tt{a.e. } x \]
\end{lemma}
\begin{proof}{Proof:}
    For $q \in \Q$, apply the theorem to $\abs{f(x) - q}$ to see
    \[ \frac{1}{\lambda(B_r(x))} \int_{B_r(x)} \abs{f(y) - q} dy \to \abs{f(x) - q} \]
    for all $x \notin A_q$ (for a null set $A_q$).\\


    If $x \notin A_q$ for any $q$, choose $q_*$ so that $\abs{f(x) - q_*} < \epsilon$. Then,
    \begin{align*}
        \limsup_{r \to 0} \frac{1}{\lambda(B_r(x))} \int_{B_r(x)} \abs{f(y) - f(x)} dy &\leq \limsup_{r \to 0} \frac{1}{\lambda(B_r(x))} \int_{B_r(x)} \abs{f(y) - q_*} dy\\
        &\quad + \limsup_{r \to 0} \frac{1}{\lambda(B_r(x))} \int_{B_r(x)} \abs{q_* - f(x)} dy\\
        &\leq 2\epsilon
    \end{align*}

    $\qed$
\end{proof}

\begin{lemma}[\textbf{Corollary 24.4}]
    Let $f : \R \to \R$ be integrable and fix some $a \in \R$. Then, let
    \[ F(x) = \int_{a}^{x} f(y) dy \]
    then $F$ is differentiable almost everywhere and $F'(x) = f(x)$ a.e.
\end{lemma}
\begin{proof}{Proof:}

    Let $h > 0$. Then,
    \[ F(x + h) - F(x) = \int_{x}^{x + h} f(y) dy \]
    so
    \begin{align*}
        \abs{\frac{F(x + h) - F(x)}{h} - f(x)} &= \abs{\frac{1}{h}\int_{x}^{x + h} f(y) - f(x) dy}\\
        &\leq \frac{2}{\lambda(B_h(x))}\int_{x - h}^{x + h} \abs{f(y) - f(x)} dy
    \end{align*}
    Then, as $h \to 0$, the RHS goes to $0$, so
    \[ \lim_{h \to 0} \frac{F(x + h) - F(x)}{h} = f(x) \]
    almost everywhere.

    $\qed$
\end{proof}

\begin{definition}[\textbf{Definition 24.5:} Regular Borel Measures]
    We say that a Borel measure $\nu$ on $\R^d$ is \emph{regular} if
    \begin{enumerate}
        \item $\nu(K) < \infty$ for any compact set $K$.
        \item $\nu(A) = \inf\{\nu(\cal{U}) \col \cal{U} \tt{ open and } A \subseteq \cal{U}\}$.
    \end{enumerate}
\end{definition}

\begin{example}{\textbf{Example:} Regular Measures}
    Counting measure is not regular! Lebesgue and Lebesgue-Stieltjes measures are.
\end{example}

\begin{lemma}[\textbf{Proposition 24.6}]
    Let $\nu$ be a finite positive regular Borel measure. Let
    \[ \nu(A) = \int_{A} f d\lambda + \rho(A) \]
    be its Lebesgue decomposition with respect to Lebesgue measure. Then, for $\lambda$-a.e. $x \in \R^d$,
    \[ \lim_{r \to 0} \frac{\nu(B_r(x))}{\lambda(B_r(x))} = f(x) \]
\end{lemma}
\begin{proof}{Proof:}
    By Lebesgue differentiation,
    \[ \lim_{r \to 0} \frac{1}{\lambda(B_r(x))} \int_{B_r(x)} f d\lambda = f(x) \quad \lambda\tt{-a.e.} \]
    Now, we want to show
    \[ \lim_{r \to 0} \frac{\rho(B_r(x))}{\lambda(B_r(x))} = 0 \]
    Choose $A$ so that $\rho(A) = \lambda(A^{c}) = 0$. Notice that $A^{c}$ is a $\lambda$-null set, so we just need to show that this limit holds on $A$. \\

    Let $A_n = \left\{x \in A \col \limsup_{r \to 0} \frac{\rho(B_r(x))}{\lambda(B_r(x))} > 1/n\right\}$. We want to show $\lambda(A_n) = 0$ for all $n$. By regularity, find $\cal{U} \supseteq A$ ($A$ is a $\rho$-null set) such that
    \[ \rho(\cal{U}) < \epsilon \]
    If $x \in A_n$, we can find $B_{r_x}(x) \subseteq \cal{U}$ so that $\rho(B_{r_x}(x)) > \frac{1}{n}\lambda(B_{r_x}(x))$. Applying the Wiener covering lemma to $\bigcup_{x} B_{r_x}(x)$, we get disjoint $B_1, \ldots, B_N$ such that
    \begin{align*}
        (1 - \delta)\lambda\left(\bigcup_{x} B_{r_x}(x)\right) &\leq 3^d\sum_{j = 1}^{N} \lambda(B_j)\\
        &\leq 3^d\sum_{j = 1}^{N} n\rho(B_j)\\
        &\leq 3^d n\rho(\cal{U})\\
        &\leq 3^d n\epsilon
    \end{align*}
    Taking $\delta = \frac{1}{2}$ and $\epsilon \to 0$, we see $\lambda\left(\bigcup_{x} B_{r_x}(x)\right) \to 0$, so $\lambda(A_n) = 0$.

    $\qed$
\end{proof}

\tocsection{\textbf{10/26/26 Lecture 25:} Differentiating Monotone Functions and Bounded Variation}

\begin{lemma}[\textbf{Proposition 25.1}]
    Let $F : \R \to \R$ be increasing and right-continuous. Then, $F'(x)$ exists a.e. and
    \[ \int_{a}^{b} F'(x) dx \leq F(b) - F(a) \]
\end{lemma}
\begin{proof}{Proof:}
    Because differentiation is local, it suffices to prove the claim for a function that agrees with $F$ on $[0, 1]$.\\

    Redefine $F$ so that $F(x) = \lim_{x \to 0^{+}} F(x)$ for $x \leq 0$ and $F(x) = F(1)$ for $x \geq 1$. Then,
    \fig{0.5}{Lecture25Fig1}{}
    $F$ is still increasing and right continuous, but its LS measure is finite. Let $\nu$ be the LS measure of this $F$. Then,
    \[ \nu(A) = \int_{A} f d\lambda + \rho(A) \quad \tt{for } \rho \perp \lambda \]
    Then,
    \begin{align*}
        \rho((a, b]) &= \nu((a, b]) - \int_{a}^{b} f(y) dy\\
        &= F(b) - F(a) - \int_{a}^{b} f(y) dy
    \end{align*}
    Now, $\rho$ is a Lebesgue-Stieltjes measure, and
    \[ \frac{F(x + h) - F(x)}{h} = \frac{1}{h}\int_{x}^{x + h} f(y) dy + \frac{\rho((x, x + h])}{h} \]
    so by the corollary from before, the first term goes to $f$ almost everywhere, and by a proposition from last time,
    \[ \frac{\rho((x, x + h])}{h} \leq \frac{\rho(x - 2h, x + 2h)}{h} = 4\frac{\rho(x - 2h, x + 2h)}{\lambda(x - 2h, x + 2h)} \to 0 \]
    because $\rho$ is mutually singular with $\lambda$. Thus,
    \[ \int_{a}^{b} F'(x) dx = \int_{a}^{b} f(x) dx \leq \nu((a, b]) = F(b) - F(a) \]

    $\qed$
\end{proof}

\begin{example}{\textbf{Example:} Cantor Function}
    Let $F$ be the Cantor function. $F$ has the property that it is constant outside the Cantor set, so $F'(x) = 0$ a.e. Thus, the L-S measure of $F$ is mutually singular with $\lambda$. Thus, we see
    \[ \int_{a}^{b} F'(x) dx = 0 \leq F(b) - F(a) \]
    Thus, we don't always achieve equality.
\end{example}

\textbf{Note:} In our proof above, we see $\int_{a}^{b} F'(x) dx = F(b) - F(a)$ when the LS measure $\nu$ is absolutely continuous with respect to $\lambda$.

\begin{lemma}[\textbf{Theorem 25.2}]
    Let $F : \R \to \R$ be increasing and let $G(x) = F(x^{+}) = \lim_{y \to x^{+}} F(y)$. Then, the set of points where $F$ is not continuous is countable, and $F$ and $G$ are differentiable with
    \[ F' = G' \tt{ a.e.} \quand \int_{a}^{b} F'(y) dy \leq F(b) - F(a) \]
\end{lemma}
\begin{proof}{Proof:}
    $G$ is increasing and right continuous.\\

    Now, the points where $F$ is discontinuous are points such that $F(x^{+}) \neq F(x^{-})$ (and namely $F(x^{-}) < F(x^{+})$).\\

    At any such point, we can find a rational in $(F(x^{-}), F(x^{+}))$, but as $F$ is increasing, all such intervals are disjoint, and thus, all such rationals are unique. Thus, the number of points $F$ is discontinuous is countable.\\

    Now, note that $G$ agrees with $F$ where $F$ is continuous. Since $G'$ exists a.e. from above, we can find $x$ such that $G'(x)$ exists, $F(x) = G(x)$, and $F$ is continuous at $x$.\\

    Fix $\epsilon > 0$. For any $h \geq 0$,
    \[ F(x + h) \leq G(x + (1 + \epsilon)h) \]
    This is because in $(x + h, x + (1 + \epsilon)h)$, we can find a point $x_h$ where $F$ and $G$ agree, and $G(x + (1 + \epsilon)h) \geq G(x_h)$. Thus,
    \begin{align*}
        \limsup_{h \to 0^{+}} \frac{F(x + h) - F(x)}{h} &\leq \limsup_{h \to 0^{+}} \frac{G(x + (1 + \epsilon)h) - G(x)}{h}\\
        &= (1 + \epsilon)\limsup_{h \to 0^{+}} \frac{G(x + (1 + \epsilon)h) - G(x)}{(1 + \epsilon)h}\\
        &= (1 + \epsilon)G'(x)
    \end{align*}
    Thus, $\limsup_{h \to 0^{+}} \frac{F(x + h) - F(x)}{h} \leq (1 + \epsilon)G'(x)$. By a similar argument, $\liminf_{h \to 0^{+}} \frac{F(x + h) - F(x)}{h} \geq (1 - \epsilon)G'(x)$. Then, taking $\epsilon \to 0$, we get $F'(x) = G'(x)$.\\

    Finally, to show $\int_{a}^{b} F'(x) dx \leq F(b) - F(a)$, given $a < b$, choose $a_n \downarrow a$ and $b_n \uparrow b$ so that $F$ and $G$ agree at these points. Then,
    \begin{align*}
        \int_{a_n}^{b_n} F'(x) dx &= \int_{a_n}^{b_n} G'(x) dx\\
        &\leq G(b_n) - G(a_n)\\
        &= F(b_n) - F(a_n)\\
        &\leq F(b) - F(a)
    \end{align*}
    where we apply MCT to get the inequality as $F'$ is non-negative almost everywhere. 

    $\qed$
\end{proof}

Let $F : \R \to \R$. Define
\[ T_F(x) = \sup\left\{\sum_{j = 1}^{n} \abs{F(x_j) - F(x_{j - 1})} \col n \in \N,\ x_0 < x_1 < \cdots < x_n = x\right\} \]
If $F$ is smooth, $T_F(x) = \int_{-\infty}^{x} \abs{F'(t)} dt$.

\begin{definition}[\textbf{Definition 25.3:} Bounded Variation]
    We say $F$ is of \emph{bounded variation}, denoted BV, if $\lim_{x \to \infty} T_F(x) = T_F(\infty) < \infty$.
\end{definition}

\begin{example}{\textbf{Example:} Bounded Increasing Functions}
    If $F$ is bounded and increasing, then $F$ is BV.
\end{example}

\begin{example}{\textbf{Example:} $x\sin(1/x)$}
    Monotonicity is important. Consider $F(x) = x\sin\left(\frac{1}{x}\right)$.
    \fig{0.3}{Lecture25Fig2}{}
    This function is bounded near 0, but not BV on any interval containing 0. 
\end{example}

\tocsection{\textbf{10/28/26 Lecture 26:} Functions of Bounded Variation}

\textbf{Recall:} If $F : \R \to \R$ is increasing, $F$ is differentiable a.e.\\

For $F : \R \to \R$, consider
\[ T_F(x) = \sup\left\{\sum_{j = 1}^{n} \abs{F(x_j) - F(x_{j - 1})} \col n \in \N,\ x_0 < x_1 < \cdots < x_n = x\right\} \]
and if $T_F(\infty) < \infty$, $F \in \tt{BV}$.\\

Now, if $x_{j - 1} < x' < x_j$,
\[ \abs{F(x_j) - F(x_{j - 1})} \leq \abs{F(x_j) - F(x')} + \abs{F(x') - F(x_{j - 1})} \]
and if $a < b$,
\[ T_F(b) = T_F(a) + \sup\left\{\sum \abs{F(x_j) - F(x_{j - 1})} \col a = x_0 < \cdots < x_n = b\right\} \]

\begin{definition}[\textbf{Definition 26.1:} Bounded Variation on an Interval]
    We say $F$ is of \emph{bounded variation on $[a, b]$} if
    \[ T_F(b) - T_F(a) < \infty \]
    and we write $F \in \tt{BV}([a, b])$.
\end{definition}

\begin{lemma}[\textbf{Proposition 26.2}]
    If $F \in \tt{BV}$, then $T_F + F$ and $T_F - F$ are increasing.
\end{lemma}
\begin{proof}{Proof:}
    Fix $x < y$. By definition, choose $x_0 < \cdots < x_n = x$ so that $\sum_{j} \abs{F(x_j) - F(x_{j - 1})} \geq T_F(x) - \epsilon$. Then,
    \begin{align*}
        T_F(y) + F(y) &\geq \sum_{j = 1}^{n} \abs{F(x_j) - F(x_{j - 1})} + \abs{F(y) - F(x)} + F(y)\\
        &\geq T_F(x) - \epsilon + F(x)
    \end{align*}
    so taking $\epsilon \to 0$, we get what we want. A similar argument works for $T_{F} - F$. 

    $\qed$
\end{proof}

\begin{theorem}[\textbf{Theorem 26.3:} Classification of BV]
    A function $F : \R \to \R$ is BV iff it may be written as a difference of bounded increasing functions. In particular, $F \in \tt{BV} \implies F$ is differentiable almost everywhere and $F$ is bounded.
\end{theorem}
\begin{proof}{Proof:}
    $(\Rightarrow)$ By the proposition, we can write
    \[ F = \frac{1}{2}(T_F + F) - \frac{1}{2}(T_F - F) \]
    and both are increasing. To show boundedness, we see that, for $x < y$,

    \[ F(x) - F(y) \leq T_F(y) - T_F(x) \leq T_F(\infty) - T_F(-\infty) < \infty \]
    and similarly
    \[ F(y) - F(x) \leq T_F(y) - T_F(x) \leq T_F(\infty) - T_F(-\infty) < \infty \]
    so we see $F$ is bounded, so $T_F + F$ and $T_F - F$ are bounded.\\


    $(\Leftarrow)$ Bounded increasing functions have bounded variation, and as $T_{H - G} \leq T_{H} + T_{G}$ by the triangle inequality, we see that the difference also has bounded variation.

    $\qed$
\end{proof}

Recall if $F$ is increasing and right-continuous, we get the LS measure $m_F$ by $m_F((a, b]) = F(b) - F(a)$.

\begin{lemma}[\textbf{Proposition 26.4}]
    Let $m$ be a positive Borel measure with the property that $m(A) < \infty$ if $A$ is bounded. Then, if
    \[ F(x) = \left\{\begin{matrix} m((0, x]) & x > 0 \\ 0 & x = 0 \\ -m((x, 0]) & x < 0 \end{matrix}\right. \]
    then $F$ is increasing, right continuous, and $m = m_F$.
\end{lemma}
\begin{proof}{Proof:}

    Increasing follows from $A \subseteq B \implies m(A) \leq m(B)$, and right continuity comes from
    \[ m((0, x]) = \lim_{h \to 0^{+}} m((0, x + h]) \]
    It is then clear that $m = m_F$.

    $\qed$
\end{proof}

\begin{lemma}[\textbf{Corollary 26.5}]
    If $m$ is a Borel measure with $m(A) < \infty$ for bounded $A$, then $m$ is regular!
\end{lemma}

This gives us a correspondence $m_F \longleftrightarrow F$, but not quite, as if $G = F + C$ for a constant $C$, then $m_F = m_G$.\\

To fix this, we will normalize as follows: let
\[ \tt{NBV} = \{F \in \tt{BV} \col F \tt{ is right continuous and } F(-\infty) = 0\} \]

\begin{theorem}[\textbf{Theorem 26.6:} NBV Functions and Signed Borel Measures]
    If $\mu$ is a finite signed measure on $\R$, then $F(x) = \mu((-\infty, x])$ is NBV. Conversely, if $F \in \tt{NBV}$, then there is a unique signed Borel measure $\mu$ with $\mu((a, b]) = F(b) - F(a)$.
\end{theorem}

Before we prove this, we need a lemma.

\begin{lemma}[\textbf{Lemma 26.7}]
    If $F \in \tt{BV}$ and is right-continuous, then so is $T_F(x)$. \
\end{lemma}
\begin{proof}{Proof:}
    Fix $x$ and let $\alpha = T_F(x^{+}) - T_F(x)$. We want to show $\alpha = 0$. Given $\epsilon > 0$, choose $\delta$ so that
    \[ \abs{F(x + h) - F(x)} < \epsilon \quand T_F(x + h) - T_F(x^{+}) < \epsilon \]
    for $0 < h < \delta$ by right-continuity and the definition of limits. Find a partition $x = x_0 < \cdots < x_n = x + h$ so that
    \begin{align*}
        \sum_{j = 1}^{n} \abs{F(x_j) - F(x_{j - 1})} &\geq \frac{3}{4}\abs{T_F(x + h) - T_F(x)}\\
        &\geq \frac{3}{4}\alpha
    \end{align*}
    \[ \sum_{j = 2}^{n} \abs{F(x_j) - F(x_{j - 1})} \geq \frac{3}{4}\alpha - \abs{F(x_1) - F(x)} \geq \frac{3}{4}\alpha - \epsilon \]

    Now, find $x = y_0 < \cdots < y_m = x_1$ so that
    \[ \sum_{j = 1}^{m} \abs{F(y_j) - F(y_{j - 1})} \geq \frac{3}{4}\abs{T_F(x_1) - T_F(x)} \geq \frac{3}{4}\alpha \]
    Thus,
    \[ \sum_{j = 1}^{m} \abs{F(y_j) - F(y_{j - 1})} + \sum_{j = 2}^{n} \abs{F(x_j) - F(x_{j - 1})} \geq \frac{3}{2}\alpha - \epsilon \]
    But as this is just a partition (of $[x, x + h]$),
    \begin{align*}
        \sum_{j = 1}^{m} \abs{F(y_j) - F(y_{j - 1})} + \sum_{j = 2}^{n} \abs{F(x_j) - F(x_{j - 1})} &\leq T_F(x + h) - T_F(x)\\
        &= (T_F(x + h) - T_F(x^{+})) + (T_F(x^{+}) - T_F(x))\\
        &< \epsilon + \alpha
    \end{align*}
    so
    \[ \frac{3}{2}\alpha - \epsilon \leq T_F(x + h) - T_F(x) < \epsilon + \alpha \]
    so taking $\epsilon \to 0$, we must see $\alpha = 0$. Since $T_F(x)$ is increasing and bounded, $T_F \in \tt{BV}$.

    $\qed$
\end{proof}

\begin{proof}{Proof of Theorem 26.6:}
    $(\Rightarrow)$ If $\mu$ is a finite signed Borel measure, write $\mu = \mu^{+} - \mu^{-}$ and let
    \[ F^{\pm}(x) = \mu^{\pm}((-\infty, x]) \]
    Each of these are BV and $F^{\pm}(-\infty) = 0$ by continuity. They are also right continuous by continuity.\\

    $(\Leftarrow)$ Write $F = F^{+} - F^{-} = \frac{1}{2}(T_F + F) - \frac{1}{2}(T_F - F)$. $F^{\pm}$ are both increasing and right continuous from above, so we get the LS measures $m_{F^{+}}$ and $m_{F^{-}}$, and it can be checked that
    \[ m = m_{F^{+}} - m_{F^{-}} \]
    has the property we want.

    $\qed$
\end{proof}

\tocsection{\textbf{10/30/26 Lecture 27:} Absolute Continuity and the Lebesgue FTC}

\textbf{Recall:} A function $F : \R \to \R$ is BV iff it can be written as a difference of bounded increasing functions (Theorem 26.3), and
\[ \tt{NBV} = \{F \in \tt{BV} \col F \tt{ is right continuous and } F(-\infty) = 0\} \]
Finite signed Borel measures on $\R$ correspond to NBV functions (Theorem 26.6).

\begin{lemma}[\textbf{Proposition 27.1}]
    If $F \in \tt{NBV}$, then $F' \in L^1$. Moreover, the LS measure $m_F$ satisfies
    \[ m_F \perp \lambda \iff F' = 0 \tt{ a.e.} \]
    and
    \[ m_F \ll \lambda \iff F(x) = \int_{-\infty}^{x} F'(y)\, dy \]
\end{lemma}

\begin{definition}[\textbf{Definition 27.2:} Absolute Continuity]
    A function $F : [a, b] \to \R$ is \emph{absolutely continuous} if for every $\epsilon > 0$ there is $\delta > 0$ such that, given finitely many disjoint intervals $(a_1, b_1), \ldots, (a_N, b_N)$ in $[a, b]$ with
    \[ \sum_{j = 1}^{N} (b_j - a_j) < \delta \]
    then
    \[ \sum_{j = 1}^{N} \abs{F(b_j) - F(a_j)} < \epsilon \]
\end{definition}

\begin{example}{\textbf{Example:} Bounded Derivative}
    If $F$ is differentiable and $\abs{F'} \leq M$, then by the mean value theorem,
    \[ \abs{F(b_j) - F(a_j)} \leq \abs{F'(c_j)}(b_j - a_j) \leq M(b_j - a_j) \]
    for some $c_j \in (a_j, b_j)$. Given $\epsilon$, make $\delta = \frac{\epsilon}{M}$.
\end{example}

\begin{example}{\textbf{Example:} The Cantor Function}
    An NBV function that is not AC is the Cantor function $F_c$, since $F_c' = 0$ a.e.
\end{example}

\begin{lemma}[\textbf{Proposition 27.3}]
    If $F \in \tt{NBV}$, then $F$ is absolutely continuous on $\R$ iff $m_F \ll \lambda$.
\end{lemma}
\begin{proof}{Proof:}
    If $m_F \ll \lambda$, for every $\epsilon$ we can find $\delta$ so that if $\lambda(A) < \delta$, then $\abs{m_F}(A) < \epsilon$. So if $\lambda\left(\bigcup_j (a_j, b_j)\right) < \delta$,
    \[ \sum_j \abs{F(b_j) - F(a_j)} = \sum_j \abs{m_F((a_j, b_j])} \leq \abs{m_F}\left(\bigcup_j (a_j, b_j]\right) < \epsilon \]

    Conversely, let $A$ such that $\lambda(A) = 0$. Let $\epsilon > 0$, and pick the $\delta > 0$ that comes from absolute continuity of $F$. By outer regularity, we can find open sets $U_1 \supseteq U_2 \supseteq \cdots \supseteq A$ so that $\abs{m_F}(U_k) \to \abs{m_F}(A)$.
    Also, we can find $V$ open so that $\lambda(V) < \delta$ and $A \subseteq V$, so $\lambda(U_n \cap V) < \delta$.\\

    Write $U_n \cap V = \bigcup_{j = 1}^{\infty} (a_{nj}, b_{nj})$ as a disjoint union of open intervals. By absolute continuity, for every finite $N$,
    \[ \sum_{j = 1}^{N} \abs{m_F((a_{nj}, b_{nj}))} = \sum_{j = 1}^{N} \abs{F(b_{nj}) - F(a_{nj})} < \epsilon \]
    Taking $N \to \infty$, $\abs{m_F(U_n \cap V)} \leq \epsilon$, so letting $n \to \infty$, $\abs{m_F(A)} \leq \epsilon$. Taking $\epsilon \to 0$, we get $m_F(A) = 0$.

    $\qed$
\end{proof}

\begin{lemma}[\textbf{Corollary 27.4}]
    If $f \in L^1$, then $F(x) = \int_{-\infty}^{x} f(y)\, dy$ is NBV and absolutely continuous, with $F' = f$ a.e.\\

    Conversely, if $F \in \tt{NBV}$ and is AC, then $F' \in L^1$ and $F(x) = \int_{-\infty}^{x} f(y)\, dy$ where $f = F'$ a.e.
    % NOTE (check): no proof on the page. First part: F(x) = mu((-infinity, x]) for d mu = f d lambda, a finite signed measure << lambda, so F is NBV (Theorem 26.6), AC (Prop 27.3), and F' = f a.e. by the Lebesgue differentiation theorem. Converse: Prop 27.3 gives m_F << lambda, then Prop 27.1.
\end{lemma}

\begin{lemma}[\textbf{Proposition 27.5}]
    If $F$ is AC on $[a, b]$, then it is BV on $[a, b]$.
\end{lemma}
\begin{proof}{Proof:}
    Choose $\delta$ so that if $\sum_j (b_j - a_j) < \delta$, then
    \[ \sum_j \abs{F(b_j) - F(a_j)} < 1 \quad (\epsilon = 1) \]
    Divide $[a, b]$ into $N$ pieces $I_1, \ldots, I_N$ with width less than $\delta$.\\

    We want to show that, given $a = x_0 < \cdots < x_n = b$,
    \[ \sum_{j = 1}^{n} \abs{F(x_j) - F(x_{j - 1})} < L \]
    for any partition.
    % NOTE (check): adding points only increases the sum (triangle inequality), so it's enough to bound partitions that contain the endpoints of the I_k.
    We can always add the endpoints of the $\delta$-intervals to the partition.\\

    % FIXED: page calls the pieces I_n, which clashes with the n in x_n; renamed them I_k.
    Group the set of points into intervals: since $\abs{I_k} < \delta$, set $N$ groups, and for every group $G$,
    % FIXED: page has |F(x_j) - F(x_j)|; changed to F(x_j) - F(x_{j-1}).
    \[ \sum_{j \in G} \abs{F(x_j) - F(x_{j - 1})} < 1 \]
    so in total,
    % NOTE (check): so L = N works. The point is that N depends only on delta and b - a, not on the partition, so the sup over partitions is at most N.
    \[ \sum_{j = 1}^{n} \abs{F(x_j) - F(x_{j - 1})} < N \]

    $\qed$
\end{proof}

\begin{theorem}[\textbf{Theorem 27.6:} Fundamental Theorem of Calculus for Lebesgue Integrals]
    For $F : [a, b] \to \R$ (on a finite interval), the following are equivalent:
    \begin{enumerate}
        \item $F$ is AC.
        \item $F(x) - F(a) = \int_{a}^{x} f(y)\, dy$ for some $f \in L^1([a, b])$.
        \item $F$ is differentiable a.e.\ with $F' \in L^1([a, b])$ and $F(x) - F(a) = \int_{a}^{x} F'(y)\, dy$.
    \end{enumerate}
\end{theorem}
\begin{proof}{Proof:}
    $(1 \Rightarrow 3)$ Let $G : \R \to \R$,
    \[ G(x) = \left\{\begin{matrix} F(b) - F(a) & x \geq b \\ F(x) - F(a) & x \in [a, b) \\ 0 & x < a \end{matrix}\right. \]
    % NOTE (check): G is NBV because F is BV on [a, b] (Prop 27.5) and continuous (AC implies continuous), so G is right continuous with G(-infinity) = 0.
    % NOTE (check): the first equality below is the converse in Corollary 27.4, which needs G to be AC on R, not just NBV (the Cantor function is NBV). G is AC on R: it's constant off [a, b] and AC on [a, b] (split any interval that straddles a or b). For the second equality, G' = F' a.e. on (a, b) and G' = 0 off [a, b]; this also gives F' in L^1.
    $G$ is NBV. This implies it is differentiable a.e., and
    % FIXED: page has G'(x) dy; changed to G'(y) dy.
    \[ F(x) - F(a) = \int_{-\infty}^{x} G'(y)\, dy = \int_{a}^{x} F'(y)\, dy \]
    which gives 3.\\

    $(3 \Rightarrow 2)$ Immediate.\\

    % NOTE (check): with f extended by 0, Corollary 27.4 says H(x) = int_{-infinity}^x f is AC on R, and F(x) = F(a) + H(x) on [a, b], so F is AC on [a, b].
    $(2 \Rightarrow 1)$ Extend $f$ by $0$ outside $[a, b]$ and apply Corollary 27.4.

    $\qed$
\end{proof}

\begin{definition}[\textbf{Definition 27.7:} Discrete and Continuous Measures]
    % NOTE (check): for the sum to define a finite signed measure you need sum |a_j| < infinity (for a positive measure, a_j >= 0).
    A Borel measure is \emph{discrete} if it can be written as
    \[ \sum_{j = 1}^{\infty} a_j \delta_{x_j}, \quad x_j \in \R^d,\ a_j \in \R \]
    A measure $\mu$ is \emph{continuous} if $\mu(\{x\}) = 0$ for all $x$.
\end{definition}

\begin{lemma}[\textbf{Proposition 27.8}]
    % NOTE (check): "AC" means mu_ac << lambda, and "singular continuous" means mu_sc is continuous and mu_sc is singular to lambda.
    Let $\mu$ be a finite positive Borel measure. Then
    \[ \mu = \underbrace{\mu_d}_{\tt{discrete}} + \underbrace{\mu_{ac}}_{\tt{AC}} + \underbrace{\mu_{sc}}_{\tt{singular continuous}} \]
\end{lemma}
\begin{proof}{Proof:}
    % FIXED: page names the set X on the first line and A afterwards; called it A throughout.
    Let $A = \{x \col \mu(\{x\}) > 0\}$. Then
    \[ \{x \col \mu(\{x\}) > 1/k\} \tt{ is finite for every } k \]
    % FIXED: page says "A is continuous"; changed to countable.
    % NOTE (check): each of these sets has at most k mu(R^d) points, and A is their union over k.
    because $\mu$ is finite, so $A$ is countable. Then $\mu|_A$ is discrete and $\mu|_{A^c}$ is continuous.\\

    % NOTE (check): mu|_A = sum_{x in A} mu({x}) delta_x by countable additivity. The Lebesgue decomposition mu|_{A^c} = mu_ac + mu_sc has both parts positive and <= mu|_{A^c}, so both are continuous; that's why mu_sc is singular *continuous*.
    Now, apply the Lebesgue decomposition to $\mu|_{A^c}$.

    $\qed$
\end{proof}

\tocsection{\textbf{11/4/26 Lecture 28:} Zorn's Lemma and Topological Spaces}

\subsection*{Preliminaries from Logic}
\sectionspace

\begin{definition}[\textbf{Definition 28.1:} Cartesian Product]
    Let $X_{\alpha}$, $\alpha \in I$ be a collection of sets. The \emph{Cartesian product} of the $X_{\alpha}$ is
    \[ \prod_{\alpha \in I} X_{\alpha} = \left\{f : I \to \bigsqcup_{\alpha} X_{\alpha} \col f(\alpha) \in X_{\alpha}\right\} \]
\end{definition}

\textbf{Question:} If every $X_{\alpha}$ is nonempty, can $\prod_{\alpha \in I} X_{\alpha}$ be empty?\\

\textbf{Answer:} Whether or not it is, is independent of ZF (you need AC).

\begin{theorem}[\textbf{Axiom 28.2:} Axiom of Choice]
    Given any collection of nonempty sets $X_{\alpha}$, $\alpha \in I$, $\prod_{\alpha \in I} X_{\alpha}$ is nonempty.
\end{theorem}

\begin{definition}[\textbf{Definition 28.3:} Partially Ordered Sets]
    A \emph{partially ordered set (poset)} $(X, \preceq)$ is a set $X$ with a binary relation $\preceq$ that is reflexive, antisymmetric, and transitive. A poset is \emph{totally ordered} if for every $x, y \in X$, $x \preceq y$ or $y \preceq x$.
\end{definition}

\begin{example}{\textbf{Example:} Power Set}
    For any $X$ with at least two elements, $(\P(X), \subseteq)$ is a poset that is not totally ordered.
\end{example}

\begin{example}{\textbf{Example:} Trees}
    A tree $T$ can be made a poset by picking a ``root'' and saying that $x \preceq y$ if the unique path from the root to $y$ contains $x$.
\end{example}

\begin{definition}[\textbf{Definition 28.4:} Chains, Upper Bounds, and Maximal Elements]
    If $(X, \preceq)$ is a poset, a subset $Y \subseteq X$ is called a \emph{chain} if it is totally ordered. An \emph{upper bound} for $Y \subseteq X$ is an element $x \in X$ such that $y \preceq x$ for all $y \in Y$. An element $x \in X$ is \emph{maximal} if $x \preceq z \implies x = z$.
\end{definition}

\begin{lemma}[\textbf{Lemma 28.5:} Zorn's Lemma]
    Let $(X, \preceq)$ be a poset. Suppose every chain in $X$ has an upper bound. Then, $X$ has a maximal element.
\end{lemma}

\textbf{Fact:} Zorn's Lemma $\iff$ Axiom of Choice.

\begin{theorem}[\textbf{Theorem 28.6:} Basis Existance Theorem]
    Every vector space has a basis.
\end{theorem}
\begin{proof}{Proof:}
    Let $X = \{A \subseteq V \col A \tt{ linearly independent}\}$ with $\subseteq$. We need to show every chain has an upper bound. Let $Y$ be a chain in $X$. We claim
    \[ U = \bigcup_{B \in Y} B \]
    is an upper bound, which is to show it is linearly independent.\\


    Let $v_1, \ldots, v_r \in U$ be distinct. We want to show these are LI. Well, $v_i \in B_i \in Y$, so as $Y$ is totally ordered, all of them lie in some $B$. However, since $B$ is LI, $v_1, \ldots, v_r$ are LI.\\

    Thus, every chain has an upper bound, so by Zorn's Lemma, we have a maximal element $\tilde{A}$. Suppose $\tilde{A}$ is not a basis. Then, there exists $v \notin \Span(\tilde{A})$. Then, $\{v\} \cup \tilde{A}$ is LI, which contradicts that $\tilde{A}$ is maximal.

    $\qed$
\end{proof}

\begin{definition}[\textbf{Definition 28.7:} Topological Spaces]
    A \emph{topological space} $(X, \cal{T})$ is a set $X$ and a collection $\cal{T} \subseteq \P(X)$ called a \emph{topology} such that
    \begin{enumerate}[label = (\roman*)]
        \item $\emptyset, X \in \cal{T}$
        \item If $U_1, \ldots, U_n \in \cal{T}$, then $\bigcap_{i} U_i \in \cal{T}$
        \item If $U_{\alpha} \in \cal{T}$, then $\bigcup_{\alpha} U_{\alpha} \in \cal{T}$
    \end{enumerate}
    Elements of $\cal{T}$ are called \emph{open sets}.
\end{definition}

\begin{example}{\textbf{Example:} Discrete and Trivial Topologies}
    If $X$ is any set, $(X, \P(X))$ and $(X, \{\emptyset, X\})$ are topological spaces.
\end{example}

\begin{example}{\textbf{Example:} Metric Spaces}
    Given a metric space $(X, d)$, let $B_r(x) = \{y \col d(x, y) < r\}$. Then, say $U$ is open if for all $x \in U$, there exists $r > 0$ such that $x \in B_r(x) \subseteq U$.
\end{example}

\begin{definition}[\textbf{Definition 28.8:} Closed Sets]
    We say $A$ is \emph{closed} if $A^{c}$ is open.
\end{definition}

By De Morgan's laws, closed sets are preserved under finite unions and arbitrary intersections.

\begin{definition}[\textbf{Definition 28.9:} Closure, Interior, and Boundary]
    Given $A \subseteq X$, the \emph{closure} of $A$, $\overline{A}$, is the smallest closed set containing $A$. Similarly, the \emph{interior} of $A$, $A^{\circ}$, is the largest open set contained in $A$. The \emph{boundary} of $A$, $\partial A$, is $\overline{A} \setminus A^{\circ}$.
\end{definition}

\tocsection{\textbf{11/6/26 Lecture 29:} Continuity, Bases, and Subspaces}

\begin{definition}[\textbf{Definition 29.1:} Continuity]
    Let $(X, \cal{T})$, $(Y, \cal{U})$ be topological spaces. A map $f : X \to Y$ is \emph{continuous} if $f^{-1}(U) \in \cal{T}$ for all $U \in \cal{U}$.
\end{definition}

\begin{definition}[\textbf{Definition 29.2:} Convergence]
    Let $(X, \cal{T})$ be a topological space. Then, $x_n$ \emph{converges} to $x$ if every open set $U \ni x$ contains all but finitely many points of the sequence.
\end{definition}

\begin{lemma}[\textbf{Proposition 29.3}]
    Let $f : X \to Y$ be continuous. Then, if $x_n \to x$, then $f(x_n) \to f(x)$. Furthermore, if $X$ is a metric space and for all $x$, $x_n \to x \implies f(x_n) \to f(x)$, then $f$ is continuous.
\end{lemma}
\begin{proof}{Proof:}
    Let $U$ be an open set containing $f(x)$. $f^{-1}(U)$ is an open set containing $x$, so it contains all but finitely many of the $x_n$. Then, all but finitely many of the $f(x_n)$ are contained in $U$.\\

    Now, suppose that $f$ is not continuous (in the topological sense). Then, there exists an open $U$ such that $f^{-1}(U)$ is not open. Thus, there exists $x \in f^{-1}(U)$ such that $B_r(x) \not\subseteq f^{-1}(U)$ for all $r > 0$. Picking $x_n$ so that $x_n \in B_{1/n}(x) \setminus f^{-1}(U)$, we see that $x_n \to x$, but as $f(x_n) \notin U$ for all $n$, $f(x_n) \not\to f(x)$, which is a contradiction.

    $\qed$
\end{proof}

\begin{example}{\textbf{Example:} The Co-Countable Topology}
    Consider $f : (\R, \cal{T}) \to (\R, \abs{\cdot})$ by $f(x) = x$, where $(\R, \abs{\cdot})$ has the usual metric topology. Let
    \[ \cal{T} = \{U \subseteq \R \mid U^{c} \tt{ is countable}\} \cup \{\emptyset\} \]
    We need to show $x_n \to x \implies f(x_n) \to f(x)$. Suppose $x_n \to x$. Let $A = \{x_n \mid x_n \neq x\}$. Now, $A^{c}$ is open and contains $x$, which implies $A^{c}$ contains all but finitely many of the $x_n$. Thus, for all but finitely many $x_n$, $x_n = x$.\\

    Thus, $f(x_n) \to f(x)$, since every open set containing $f(x)$ contains all $f(x_n)$ with $x_n = x$.\\

    However, $f$ is not continuous, since the preimage of nearly every open set is not open in $\cal{T}$. Thus, the metric topology requirement in the proposition above is important.
\end{example}

\begin{definition}[\textbf{Definition 29.4:} Sub-Bases and Bases]
    Let $X$ be a set and $\cal{B} \subseteq \P(X)$. We say that $\cal{B}$ is a \emph{sub-basis} if every $x \in X$ is contained in some $B \in \cal{B}$. We say a sub-basis $\cal{B}$ is a \emph{basis} if whenever $x \in B_1 \cap B_2$ for $B_1, B_2 \in \cal{B}$, there exists $B_3 \in \cal{B}$ such that $x \in B_3 \subseteq B_1 \cap B_2$.
\end{definition}

\begin{lemma}[\textbf{Proposition 29.5}]
    Let $X$ be a set and $\cal{B}$ a sub-basis.
    \begin{enumerate}
        \item Let $\cal{T}$ denote the collection of subsets of $X$ of the form
        \[ O = \bigcup U_{\alpha} \]
        where $U_{\alpha}$ is the intersection of finitely many elements of $\cal{B}$. Then, $\cal{T}$ is a topology.
        \item If $\cal{B}$ is a basis, and $\cal{U}$ is the collection of arbitrary unions of sets in $\cal{B}$, then $\cal{U}$ is a topology.
    \end{enumerate}
\end{lemma}
\begin{proof}{Proof:}
    \textbf{(1)} Since $\cal{B}$ is a sub-basis, $X \in \cal{T}$, and taking an empty union, $\emptyset \in \cal{T}$. Arbitrary unions are immediate. Then, for
    \[ O_1 = \bigcup U_{\alpha} \qquad O_2 = \bigcup U_{\beta} \]
    we see
    \[ O_1 \cap O_2 = \bigcup_{\alpha, \beta} U_{\alpha} \cap U_{\beta} \]
    which is of the desired form. Thus, $\cal{T}$ is a topology.\\

    \textbf{(2)} As before, we just need to check finite intersections. Let $O_1, O_2 \in \cal{U}$ and let $x \in O_1 \cap O_2$. Then, $x \in B_1 \subseteq O_1$ and $x \in B_2 \subseteq O_2$. Thus, $x \in B_x \subseteq B_1 \cap B_2 \subseteq O_1 \cap O_2$. Then, we can write
    \[ O_1 \cap O_2 = \bigcup_{x \in O_1 \cap O_2} B_x \]
    which is in $\cal{U}$.

    $\qed$
\end{proof}

We call the topologies above the \emph{topology generated by $\cal{B}$}.

\begin{example}{\textbf{Example:} Countable Bases for $\R^2$}
    A countable basis for $\R^2$ with the metric topology is given by
    \[ \{(a, b) \times (c, d) \col a, b, c, d \in \Q\} \]
    An example of a countable sub-basis for $\R^2$ is given by
    \[ \{(a, b) \times \R \col a, b \in \Q\} \cup \{\R \times (c, d) \col c, d \in \Q\} \]
\end{example}

Given a topological space $(X, \cal{T})$ and $Y \subseteq X$, we can give this a topology by letting $Z \subseteq Y$ be open if and only if it is given by $U \cap Y$ for $U \in \cal{T}$. We call this the \emph{subspace topology} on $Y$.

\begin{example}{\textbf{Example:} The Circle}
    Let $X = \R^2$ and $S^1 = \{x \in \R^2 \col \abs{x} = 1\}$.
    \fig{0.7}{Lecture29Fig1}{}
    Note that the arc alone is not open in $\R^2$, but is in $S^1$.
\end{example}

\begin{definition}[\textbf{Definition 29.6:} Homeomorphisms]
    A map $f : X \to Y$ is a \emph{homeomorphism} if it is a continuous bijection and its inverse is continuous.
\end{definition}

\begin{example}{\textbf{Example:} $1/x$}
    The map $f : (0, 1) \to (1, \infty)$ by $f(x) = \frac{1}{x}$ is a homeomorphism.
\end{example}

\tocsection{\textbf{11/9/26 Lecture 30:} Quotients, Products, and Compactness}

\textbf{Recap:} When given a sub-basis, we can make a topology using unions of finite intersections of elements of the sub-basis. When given a basis, we can do the same thing with just unions.\\

Moving onto homeomorphisms, consider the following example.

\begin{example}{\textbf{Example:} A Continuous Bijection That Isn't a Homeomorphism}
    Let $f : [0, 1) \to S^1$ by $f(x) = e^{2\pi ix}$. $f$ is continuous and bijective, but its inverse is not continuous. This is because $[0, 0.5)$ is open, but the arc $[1, -1) \subseteq S^1$ is not.
\end{example}

\subsection*{Quotient Topology}
\sectionspace

Let $(X, \cal{T})$ be a topological space and $\sim$ an equivalence relation on $X$. Now, consider
\[ X/{\sim} = \{[x] \col x \in X\} \]
We know there exists a canonical quotient map $q : X \to X/{\sim}$ by $x \mapsto [x]$. We define the topology on $X/{\sim}$ to be given by
\[ U \subseteq X/{\sim} \tt{ is open} \iff q^{-1}(U) \tt{ is open in } X \]

\begin{example}{\textbf{Example:} $\R/\Z$}
    Consider $\R$ with $x \sim y \iff x - y \in \Z$. Then, we see
    \[ \R/{\sim} = \{[x] \col x \in [0, 1)\} \]
    Another useful way to think of it is $[0, 1]$ but we identify $0$ and $1$. Note that this gives $[0, 1)$ a topology different from the subspace topology, since $[0, 0.5)$ is not open in this topology.
\end{example}

\begin{lemma}[\textbf{Proposition 30.1}]
    Let $X$ be a topological space with equivalence relation $\sim$. Let $f : X \to Y$ be continuous and constant on all equivalence classes. Then, there exists a continuous map $\tilde{f} : X/{\sim} \to Y$ such that $f = \tilde{f} \circ q$.
    \begin{center}
        \begin{tikzcd}
            X \arrow[d, "q"'] \arrow[dr, "f"] & \\
            X/{\sim} \arrow[r, "\tilde{f}"'] & Y
        \end{tikzcd}
    \end{center}
\end{lemma}
\begin{proof}{Proof:}
    Let $\tilde{f}([x]) = f(x)$, which is well defined as $f$ is constant on each equivalence class. Now, if $U$ is open in $Y$, then
    \[ q^{-1}(\tilde{f}^{-1}(U)) = f^{-1}(U) \]
    which is open as $f$ is continuous, so $\tilde{f}^{-1}(U)$ is open in $X/{\sim}$, since $q^{-1}(V)$ is open iff $V$ is open. 

    $\qed$
\end{proof}

\begin{example}{\textbf{Example:} $\R/\Z \cong S^1$}
    If $f(x) = e^{2\pi ix}$, it respects $x \sim y \iff x - y \in \Z$. Thus, we have a map $\tilde{f} : \R/\Z \to S^1$ that turns out to be a homeomorphism.
\end{example}

\subsection*{Product Topology}
\sectionspace

Let $X = \prod_{\alpha \in I} X_{\alpha}$ where each $X_{\alpha}$ is a topological space. We want to define a topology on $X$. To motivate, $\R^2 = \R \times \R$ is generated by the basis $U \times V$ ($U \subseteq \R$, $V \subseteq \R$ open).

\begin{definition}[\textbf{Definition 30.2:} Box Topology]
    The \emph{box topology} on $X$ is generated by the basis
    \[ \left\{\prod_{\alpha} U_{\alpha} \mid U_{\alpha} \subseteq X_{\alpha} \tt{ is open}\right\} \]
\end{definition}

This is not ideal, however, and we will show why.

\begin{example}{\textbf{Example:} The Issue with the Box Topology}
    Say we have $f_{\alpha} : Y \to X_{\alpha}$. We can form $f : Y \to X$ by $f(y) = \{f_{\alpha}(y)\}_{\alpha \in I}$. One desired property to want is that $f_{\alpha}$ continuous $\implies f$ continuous. It turns out the box topology doesn't give this.\\

    Let $X = \prod_{n = 1}^{\infty} \R$ with the box topology. Let $f_n(x) = x$, $f_n : \R \to \R$. Now, this gives us a map $f : \R \to X$ by
    \[ f(x) = (x, x, \ldots) \]
    It turns out this is not continuous! Consider
    \[ U = (-1, 1) \times \left(-\frac{1}{2}, \frac{1}{2}\right) \times \left(-\frac{1}{3}, \frac{1}{3}\right) \times \cdots \]
    This is open, but
    \[ f^{-1}(U) = \{0\} \]
    so $f$ is not continuous!
\end{example}

Now, given $X = \prod_{\alpha} X_{\alpha}$, we get a map $\pi_{\alpha} : X \to X_{\alpha}$ by $\pi_{\alpha}(x) = x(\alpha) := x_{\alpha}$.

\begin{definition}[\textbf{Definition 30.3:} Product Topology]
    The \emph{product topology} is the topology generated by the sub-basis
    \[ \bigcup_{\alpha \in I}\{\pi_{\alpha}^{-1}(U) \col U \subseteq X_{\alpha} \tt{ open}\} \]
    Note
    \[ \pi_{\alpha}^{-1}(U) = U \times \prod_{\alpha' \neq \alpha} X_{\alpha'} \]
    This topology can also be described as being generated by the basis
    \[ \left\{\prod_{\alpha} U_{\alpha} \col U_{\alpha} \subseteq X_{\alpha} \tt{ is open and equal to } X_{\alpha} \tt{ for all but finitely many } \alpha\right\} \]
\end{definition}

Note that the product and box topologies are the same on finite products. However, in general, the box topology is larger, or ``finer''.

\begin{definition}[\textbf{Definition 30.4:} Finer and Coarser Topologies]
    Let $\cal{T}$ and $\cal{T}'$ be topologies on the same space $X$. Suppose $\cal{T} \subseteq \cal{T}'$. We say $\cal{T}'$ is \emph{finer} than $\cal{T}$ and $\cal{T}$ is \emph{coarser} than $\cal{T}'$.
\end{definition}

\begin{lemma}[\textbf{Proposition 30.5}]
    Let $f : A \to \prod_{\alpha} X_{\alpha}$ such that $f(a) = (f_{\alpha}(a))_{\alpha}$. Then, $f$ is continuous with respect to the product topology iff $f_{\alpha}$ is for all $\alpha$.
\end{lemma}
\begin{proof}{Proof:}
    If $f$ is continuous, then as $f_{\alpha} = \pi_{\alpha} \circ f$, we see $f_{\alpha}$ is continuous.\\

    Conversely, if $f_{\alpha}$ is continuous for all $\alpha$, let
    \[ U = U_{\alpha_1} \times \cdots \times U_{\alpha_n} \times \prod_{\alpha \neq \alpha_i} X_{\alpha} \]
    be a basis element. Then,
    \[ f^{-1}(U) = \bigcap_{i = 1}^{n} f_{\alpha_i}^{-1}(U_{\alpha_i}) \]
    which is open. Thus, as the preimages of basis elements are open, $f$ is continuous.

    $\qed$
\end{proof}

\subsection*{Compactness}
\sectionspace

\begin{definition}[\textbf{Definition 30.6:} Open Covers]
    Let $(X, \cal{T})$ be a topological space. An \emph{open cover} of $X$ is a collection $\{U_{\alpha}\}$ of open sets such that $X = \bigcup_{\alpha} U_{\alpha}$. A \emph{finite subcover} is a finite subcollection $\{U_{\alpha_i}\}_{i = 1}^{n}$ such that $X = \bigcup_{i = 1}^{n} U_{\alpha_i}$.
\end{definition}

\begin{definition}[\textbf{Definition 30.7:} Compactness]
    A topological space $(X, \cal{T})$ is \emph{compact} if every open cover has a finite subcover.
\end{definition}

% NOTE (check): here the U_alpha are open sets of X (their traces U_alpha cap A are the open sets of A), and it's "iff": A is compact in the subspace topology exactly when every such cover of A has a finite subcover.
\textbf{Note:} If $A \subseteq X$, $A$ is compact in the subspace topology iff
\[ A \subseteq \bigcup_{\alpha} U_{\alpha} \implies A \subseteq \bigcup_{i = 1}^{n} U_{\alpha_i} \]

\begin{lemma}[\textbf{Proposition 30.8}]
    A closed subset of a compact set is compact.
\end{lemma}
\begin{proof}{Proof:}
    Let $A \subseteq X$ be closed. Then, if $\{U_{\alpha}\}$ is an open cover of $A$, then $\{U_{\alpha}\} \cup \{A^{c}\}$ is an open cover of $X$. Now, we get a finite subcover
    \[ X = \bigcup_{i = 1}^{n} V_{\alpha_i} \implies A \subseteq \bigcup_{i = 1}^{n} V_{\alpha_i} \]
    % FIXED: page has the V's drawn from {U_alpha} cup {A}; changed A to A^c (that's the set added to make the cover of X, and the one removed next).
    where $V_{\alpha_i}$ is drawn from $\{U_{\alpha}\} \cup \{A^{c}\}$. Removing $A^{c}$ from the collection, we still cover $A$ with a finite subcollection of $\{U_{\alpha}\}$.

    $\qed$
\end{proof}



\end{document}






